% BEGIN INCLUDED SGA4_English_current_buildable_source_20260906/source/06/drafts/SGA4_VI_LINES_44_390_English.texfrag
\setcounter{pageoriginal}{164}

\setcounter{section}{-1}
\section{引言}\etiquette[0]{VI.sec0}

本篇\pageoriginale{}研究两类密切相关的问题.

首先, 系统研究拓扑斯的有限性条件,
依次考察\emph{拟紧性（quasi-compactness）}、
\emph{拟分离性（quasi-separatedness）}与\emph{凝聚性（coherence）}
(= 拟紧性 $+$ 拟分离性) 的概念.
它们来自概形理论中熟悉的类似概念.
我们先对拓扑斯的对象 $X$ 及箭头 $u : X \to Y$
研究这些概念 (\hyperref[VI.sec1]{\S 1}),
再对拓扑斯 $E$ 本身 (\hyperref[VI.sec2]{\S 2}),
最后对拓扑斯态射 $f : E \to E'$ (\hyperref[VI.sec3]{\S 3}) 研究它们.
\hyperref[VI.sec4]{第 4 节} 研究由两个拓扑斯粘合 (IV 9) 得到的拓扑斯这一特殊情形.
为了阅读后续各节, 在这些稍显技术性的讨论中须记住的基本事实如下:
\emph{凝聚拓扑斯}是等价于某个形如 $C^{\sim}$ 的拓扑斯的拓扑斯,
其中 $C$ 是小景, 有限投射极限在其中可表示,
且每个覆盖族都含有一个有限覆盖子族;
凝聚拓扑斯之间的\emph{凝聚态射} $f : E \to E'$
是可以在等价意义下描述为由景态射 $f : C \to C'$ 伴随的态射,
其中 $C$ 与 $C'$ 如上, 且相应函子 $f^{\ast} : C'\to C$ 左正合.

指出下面一点无疑有益:
使用这里的有限性概念, 尤其是凝聚拓扑斯概念,
我们坚决地 (而且是第一次) 离开分析学家和“拓扑学家”熟悉的拓扑空间\footnote{引号表示这是一个不熟悉拓扑斯概念的“\ldots”
(\cf IV \hyperref[IV.subsec0.4]{0.4} 中“拓扑”一词的含义).}.
例如, 与一个分离拓扑空间 $X$ 相伴的拓扑斯
(IV \hyperref[IV.subsec2.1]{2.1}) 凝聚, 只有当 $X$ 为有限集时才可能!
换句话说, 本篇完全面向代数几何与代数中出现的各种拓扑斯.
事实上, 迄今在这些学科中使用的一切拓扑斯
(当然, 复数域上的代数几何超越方法除外!) 都是局部凝聚的.

其次, 对拓扑斯的上同调建立两个\emph{过渡到归纳极限的定理}.
其中一个 (\hyperref[VI.cor5.2]{5.2}) 表明,
凝聚拓扑斯上的函子 $H^{i}(E, -)$\pageoriginale{}
与 Abel 层的归纳极限交换.
另一个 (\hyperref[VI.subsec8.7]{8.7}) 基本上是说:
若拓扑斯 $X$ 表示为凝聚拓扑斯 $X_i$ 的滤过投射极限,
转移态射 $f_{ij} : X_j \to X_i$ 凝聚,
则 $X$ 的上同调 (取值于任意 Abel 层)
可作为各 $X_i$ 的上同调的归纳极限来计算.
要赋予这个陈述以精确含义,
主要须明确“拓扑斯的投射极限”的概念,
这将在 \S\hyperref[VI.sec7]{\S 7}, 8 中完成.
这是很有用的几何概念,
无疑在概形投射极限以外的背景中也同样有用
(概形投射极限曾是 M.\textsc{Artin} 最初定义这一概念时的模型).
例如, 我们在 (\hyperref[VI.subsec8.4]{8.4}) 中表明,
通过将局部化定义为一点的各“邻域”的投射极限,
可以统一各种“在一点局部化”的概念
(对于 Noether 空间、配备通常拓扑的概形, 甚至配备平展拓扑的概形).
这与一点处局部化的直观想法完全一致.

这两个过渡到极限的定理将在下一篇中,
针对概形的平展上同调这一特殊背景再次明确陈述
(\hyperrefext[VII][VII.prop3.3]{3.3} 及 \hyperref[VI.prop5.7]{5.7});
在此情形下, 本讨论班余下部分会不断使用它们,
而且迄今几乎所有运用平展上同调的地方也都如此.
因此, 如果读者愿意接受第 VII 篇陈述的这两个特殊情形的定理,
且只关心平展上同调的应用, 就可以略过本篇.
诚然, 拓扑斯的\emph{可构造对象（constructible object）}概念
(\ie 到终对象 $e$ 的结构态射 $X \to e$ 凝聚的对象)
也将在本讨论班余下部分发挥重要作用
(而它在本篇中只起很次要的作用,
仅仅因为在凝聚拓扑斯中, 它与这里居于核心的凝聚对象概念重合).
不过, 第 IX 篇在概形上的平展层这一特殊情形下,
不使用本篇的一般研究而独立发展了这一概念;
在此情形中, 它还呈现出很有用的特殊现象,
这些现象构成了第 IX 篇对它进行处理的基础
(从 \hyperref[IX.def2.3]{定义 IX 2.3} 开始).
由于第 IX 篇的文字比当前\pageoriginale{}第 VI 篇的文字早五年,
但读者使用起来仍无妨碍,
故我们只在第 IX 篇末证明第 IX 篇采用的可构造性概念与这里引入的概念等价.

\section{拓扑斯的对象与箭头的有限性条件}\etiquette[1]{VI.sec1}

\begin{enonce*}{定义 1.1}\etiquette[1.1]{VI.def1.1}
设 $C$ 是一个景.
$C$ 的对象 $X$ 称为拟紧的, 是指对每个覆盖族
$(X_i \to X)_{i\in I}$,
指标集 $I$ 都存在一个有限子集 $J$,
使子族 $(X_i \to X)_{i \in J}$ 仍覆盖.
\end{enonce*}

由于 $𝒰$-拓扑斯 $E$ 总被看作景
(IV \hyperref[IV.subsubsec1.1.1]{1.1.1}),
上述定义特别适用于 $E$ 的对象.
而且由下一个命题, 这一情形就是一般情形.

\begin{enonce*}{命题 1.2}\etiquette[1.2]{VI.prop1.2}
设 $C$ 是一个 $𝒰$-景,
$\epsilon : C \to C^{\sim}$ 为典范函子,
$X$ 是 $C$ 的一个对象.
则 $X$ 是景 $C$ 的拟紧对象,
当且仅当 $\epsilon(X)$ 是拓扑斯 $C^{\sim}$ 的拟紧对象.
\end{enonce*}

假定 $X$ 拟紧; 证明 $\epsilon(X)$ 也拟紧.
设 $(F_i \to \epsilon(X))_{i \in I}$ 为一个满族.
对每个 $i \in I$, 选取满族
$\epsilon(X_{ij}) \to F_i$, $j \in J_i$.
复合 $\epsilon(X_{ij}) \to \epsilon(X)$
等同于 $\epsilon(X)(X_{ij})$ 的一个元素 $\alpha_{ij}$.
一般不能断言它来自一个态射 $X_{ij} \to X$,
但由相伴层 $\epsilon(X)$ 的构造可知,
对固定的 $i$ 和 $j$, 可以找到覆盖族 $Y_k \to X_{ij}$,
使各 $\alpha_{ij}$ 在 $\epsilon(X)(Y_k)$ 中的像来自 $X(Y_k)$ 的元素,
\ie 来自态射 $Y_k \to X$.
因此, 加细满族 $\epsilon(X_{ij}) \to F_i$ 后,
可以假定各复合 $\epsilon(X_{ij}) \to \epsilon(X)$ 来自态射
$X_{ij} \to X$.\pageoriginale{}
当 $i$ 与 $j$ 变化时,
各 $\epsilon(X_{ij}) \to \epsilon(X)$ 构成满族,
故 (II \hyperref[VI.cor4.4]{4.4}) 各 $X_{ij} \to X$ 构成覆盖族.
由对 $X$ 的假设, 它具有有限覆盖子族,
相应的有限态射族 $\epsilon(X_{ij}) \to \epsilon(X)$ 仍为满族.
由于这些态射通过有限多个态射 $F_i \to \epsilon(X)$ 分解,
那一有限族就已是满族.
这就证明 $\epsilon(X)$ 拟紧.

反过来, $\epsilon(X)$ 的拟紧性蕴含 $X$ 的拟紧性,
由 \loccit{} 是显然的.

因此, 上述结果说明,
景中的拟紧性概念可归结为拓扑斯中的相应概念.

\begin{enonce*}{命题 1.3}\etiquette[1.3]{VI.prop1.3}
设 $C$ 是一个景, $(X_i \to X)_{i \in I}$ 为有限覆盖族,
且各 $X_i$ 拟紧. 则 $X$ 拟紧.
\end{enonce*}

为此, 由 \hyperref[VI.prop1.2]{1.2}, 可假定 $C$ 是一个拓扑斯.
设 $(X'_j \to X)$ 为 $C$ 中的一个覆盖族, \ie 满族.
对每个 $i \in I$, 经基变换 $X_{i} \to X$ 得到的族为满族,
因而 (由于 $X_i$ 拟紧) 具有一个有限满子族,
对应于指标集 $J$ 的一个有限子集 $J_i \subset J$.
置 $J'= \bigcup_{i\in I} J_i$.
由于 $I$ 有限, $J'$ 是 $J$ 的有限子集,
而族 $(X'_j \to X)_{j \in {J'}}$ 已为满族,
因为它局部是满族, \cqfd

\begin{enonce*}{推论 1.4}\etiquette[1.4]{VI.cor1.4}
设\pageoriginale{} $(X_i)_{i \in I}$ 是拓扑斯 $E$ 的一族对象,
$X$ 是它们的和.
则 $X$ 拟紧, 当且仅当所有 $X_i$ 都拟紧,
且除有限多个以外的 $X_i$ 都同构于 $E$ 的“空对象” $∅_{E}$.
\end{enonce*}

充分性立即来自 \hyperref[VI.prop1.3]{1.3}
(因为去掉为空的和项, 和不变).
证明必要性, 注意必须有有限子族 $(X_i)_{i \in J}$ 覆盖 $X$,
这蕴含对 $i \in I-J$, 有
$X_i = \fprod{X_{i}}{\amalg_{j=J}}{X}X_{j} =
\fprod{\amalg_{j=J}X_{i}}{X_j}{X} = ∅$
(II \hyperref[VI.lem4.5.1]{4.5.1}).
另一方面, 要证明每个 $X_i$ 拟紧,
注意 $X \simeq X_i \amalg {X'}_i$,
而 $X_i$ 的每个覆盖族 $Y_j \to X_i$
都定义了 $X$ 的一个覆盖族
$Y_j \amalg {X'}_i \to X_i \amalg {X'}_i = X$.
后者具有有限覆盖子族,
这显然蕴含相应的有限态射族 $Y_j \to X_i$ 也覆盖;
所以 $X_i$ 拟紧.

\begin{enonce*}[remark]{注 1.5.1}\etiquette[1.5.1]{VI.rem1.5.1}
拓扑斯 $E$ 的对象 $X$ 拟紧这一性质,
只依赖于诱导拓扑斯 $E_{/X}$,
且意味着后者的终对象拟紧.
同样, 景 $C$ 的对象 $X$ 拟紧,
当且仅当诱导景 $C_{/X}$ (III \hyperref[III.subsec5.1]{5.1}) 的终对象拟紧.
\end{enonce*}

\setcounter{subsection}{5}
\setcounter{subsubsection}{1}
\subsubsection{}\etiquette[1.5.2]{VI.subsubsec1.5.2}
设 $X$ 是拓扑斯 $E$ 的一个对象.
则 $X$ 拟紧, 当且仅当 $X$ 的每个覆盖 $X$ 的滤过递增子对象族
$(X_i)_{i \in I}$ 中都有某个 $X_i$ 等于 $X$.
因此, 这一条件\pageoriginale{}只依赖于 $X$ 的子对象所成的有序集
(\ie 诱导拓扑斯 $E_{/X}$ 的开子集).

\subsubsection{}\etiquette[1.5.3]{VI.subsubsec1.5.3}
考察 $E$ 的一个生成族 $(X_i)_{i \in I}$.
由定义和 \hyperref[VI.prop1.3]{1.3} 可知,
$E$ 的对象 $X$ 拟紧的必要条件是,
它具有由有限多个 $X_i$ 构成的覆盖族,
\ie $X$ 同构于 $E$ 的对象的一个有限和的商对象;
当各 $X_i$ 拟紧时, 这一条件也充分.
还要注意, 在对 $E$ 所作的假设下,
必要时将所考察的生成族换成一个子族后,
显然可以假定 $I \in 𝒰$.
利用 $E$ 的一个对象的商集为小集这一事实
(I \hyperref[I.cor7.5]{7.5}),
可知由 $E$ 的拟紧对象构成的 $E$ 的满子范畴 $E_{\qc}$
等价于一个 $\in 𝒰$ 的范畴,
\ie $E_{\qc}$ 的对象的同构类集的基数 $\in 𝒰$.

\begin{enonce*}[remark]{例 1.6.1}\etiquette[1.6.1]{VI.exem1.6.1}
设 $X$ 是一个拓扑空间.
$X$ 的开子集 $U$ 是 $X$ 的开子集景 $\Ouv(X)$ 的拟紧对象,
当且仅当 $U$ 配备由 $X$ 诱导的拓扑后是拟紧空间.
更一般地, $X$ 上的层 $F$ 是拓扑斯 $\Top(X)$ 的拟紧对象,
当且仅当与 $F$ 相伴的 $X$ 上的展空间 $X'$ 拟紧.
(这来自前一个断言、\hyperref[VI.rem1.5.1]{1.5}
及 $\Top (X)_{/F}$ 等价于 $\Top(X')$ 这一事实
(IV \hyperref[IV.subsec5.7]{5.7})).
\end{enonce*}

\setcounter{subsection}{6}
\setcounter{subsubsection}{1}
\subsubsection{}\etiquette[1.6.2]{VI.subsubsec1.6.2}
在\pageoriginale{}概形 (即 $𝒰$ 的元素) 范畴 ($\Sch$) 上,
考察 \cite{sga4.exp.VI.bib.6} \hyperref[VI.def6.3]{6.3} 中的一个拓扑 $T_i$
(例如 Zariski 拓扑、平展拓扑、fppf 拓扑或 fpqc 拓扑).
概形关于这个拓扑拟紧,
当且仅当它是通常意义下的拟紧概形.
(沿用 \loccit{} 的记号, 这里只用到每个族 $\in P'$ 都有限这一事实.)

\begin{enonce*}{定义 1.7}\etiquette[1.7]{VI.def1.7}
设 $E$ 是一个拓扑斯, $f : X \to Y$ 是 $E$ 的一个箭头.
若对每个 $Y'$ 拟紧的箭头 $Y'\to Y$,
对象 $X'= \fprod{X}{Y'}{Y}$ 也拟紧,
就称 $f$ 为拟紧态射.
若对角态射 $X \to \fprod{X}{X}{Y}$ 拟紧,
就称 $f$ 拟分离.
若 $f$ 拟紧且拟分离, 就称 $f$ 凝聚.
\end{enonce*}

\setcounter{subsection}{7}
\setcounter{subsubsection}{0}
\subsubsection{}\etiquette[1.7.1]{VI.subsubsec1.7.1}
展开拟分离态射 $f : X \to Y$ 的定义可见,
它也意味着: 对每对“在 $Y$ 上”的箭头
$g_1, g_2 : X'\rightrightarrows X$,
\ie 满足 $fg_1 = fg_2$ 的箭头,
$X'$ 的拟紧性蕴含 $\Ker(g_1, g_2)$ 的拟紧性.

\begin{enonce*}{命题 1.8}\etiquette[1.8]{VI.prop1.8}
设 $E$ 是一个拓扑斯.
\begin{enumerate}
\renewcommand{\theenumi}{\roman{enumi}}
\renewcommand{\labelenumi}{(\theenumi)}
\item $E$ 中的每个同构都是凝聚态射,
\ie 拟紧且拟分离.
两个拟紧 (\resp 拟分离, \resp 凝聚) 态射的复合
拟紧 (\resp 拟分离, \resp 凝聚).

\item 设\pageoriginale{} $f : X \to Y$ 和 $g : Y' \to Y$ 为 $E$ 中的态射,
$f': X'= \fprod{X}{Y'}{Y} \to Y'$ 由 $f$ 经基变换 $g$ 得到.
若 $f$ 拟紧 (\resp 拟分离, \resp 凝聚),
则 $f'$ 也如此.

\item 设 $f : X \to Y$ 和 $g : Y \to Z$ 是 $E$ 中的态射.
若 $gf$ 拟分离, 则 $f$ 拟分离.

\item 设 $f : X \to Y$ 和 $g : Y \to Z$ 是 $E$ 中的态射,
且 $g$ 拟分离.
若 $gf$ 拟紧 (\resp 凝聚), 则 $f$ 拟紧 (\resp 凝聚).
\end{enumerate}
\end{enonce*}

这些证明是熟知的 (\Cf \EGA IV 1 或 EGA I, 第 2 版).
在拟紧情形, 陈述 (i) 与 (ii) 由定义显然得到.
在拟分离情形, (ii) 中的基变换稳定性,
以及每个同构拟分离这一事实,
都立即来自定义及前一种情形.
复合 $X \to Y \to Z$ 的稳定性来自图
\[\vcenter{
\xymatrix@R=1.5cm@C=1.5cm{
X \ar[d] \ar@/_3.5pc/[dd]^(.2){\Delta_{X/Y}}_{\Delta_{X/Z}}&\\
\fprod{X}{X}{Y} \ar[d]^{u} \ar[rr]&  &Y \ar[d]^{\Delta_{Y/Z}}\\
\fprod{X}{X}{Z}  \ar[rr]&        & \fprod{Y}{Y}{Z}
}}\leqno{(*)}
\]
其中的方形\pageoriginale{}为笛卡儿方形.
由假设, $\Delta_{Y/Z}$ 拟紧,
故由它经基变换得到的 $u$ 也拟紧;
此外 $\Delta_{X/Y}$ 拟紧,
因此复合 $u\Delta_{X/Y} = \Delta_{X/Z}$ 也拟紧.
这就证明了 (i) 与 (ii);
凝聚情形由前两种情形合并得到.
要证明 (iv) 中首先考察的情形,
考察下列各方形均为笛卡儿方形的图:
\[\vcenter{
\xymatrix@C=2.75cm@R=1.5cm{
  &     X \ar[r]^{gf}          & Z\\
  X \ar[r]^{\Gamma_f=(\id_X,f)}  \ar[d]
  \ar@/_2pc/[rr]_(.55){f\qquad\qquad \fprod{f}{\id_{Y}}{Z}}  &\fprod{X}{Y}{Z} \ar[r]^{\pr_2}\ar[d]\ar[u]^{\pr_1}& Y\ar[u]_{g}\\
Y \ar[r]^{\Delta_{Y/Z}}&  \fprod{Y}{Y}{Z}& \quad .
}}\leqno{(**)}
\]
由假设, $g$ 拟分离,
故 $\Delta_{Y/Z}$ 拟紧,
从而经基变换得到的 $\Gamma_f$ 拟紧.
同样, 由假设 $gf$ 拟紧,
故经基变换得到的 $\pr_2$ 也拟紧.
因此, $f = \pr_2 \Gamma_f$ 作为两个拟紧态射的复合而拟紧.
(iv) 的相应情形立即来自前一种情形及尚待证明的 (iii).
为此回到图 $(\ast)$.
假设 $gf$ 拟分离意味着 $\Delta_{X/Z} = u \Delta_{X/Y}$ 拟紧,
而结论 $f$ 拟分离意味着 $\Delta_{X/Y}$ 拟紧.
一旦证明 $u$ 拟分离, 这就由 (iv) 得到.
不过 $u$ 显然\pageoriginale{}是单态射,
而我们确有如下显然结果:

\begin{enonce*}{推论 1.8.1}\etiquette[1.8.1]{VI.cor1.8.1}
每个单态射都拟分离.
\end{enonce*}

事实上, 若 $u : X \to Y$ 为单态射,
则 $\Delta_{X/Y}$ 是同构, 因而拟紧, \cqfd

\begin{enonce*}{推论 1.8.2}\etiquette[1.8.2]{VI.cor1.8.2}
设 $f:X \to Y$ 和 $f':X' \to Y'$ 是 $E$ 中的两个 $S$-态射.
若 $f$ 和 $f'$ 拟紧 (\resp 拟分离, \resp 凝聚),
则 $\fprod{f}{f'}{S} : \fprod{X}{X'}{S} \to \fprod{Y}{Y'}{S}$ 也如此.
\end{enonce*}

通常地合并 (i) 与 (ii) 即得.

\begin{enonce*}{注 1.9.1}\etiquette[1.9.1]{VI.rem1.9.1}

% END INCLUDED SGA4_English_current_buildable_source_20260906/source/06/drafts/SGA4_VI_LINES_44_390_English.texfrag


% BEGIN INCLUDED SGA4_English_current_buildable_source_20260906/source/06/drafts/SGA4_VI_LINES_391_677_English.texfrag
设 $E$ 是任意范畴,
$Q$ 是 $\ob E$ 的一个在同构下稳定的子集
(因此, 给定 $Q$ 就相当于给定 $E$ 的一个严格满子范畴),
扮演“拟紧”对象的角色.
假定 $E$ 中纤维积可表示.
于是可以仿照定义 \hyperref[VI.def1.7]{1.7},
定义 $Q$-拟紧、$Q$-拟分离及 $Q$-凝聚态射的概念.
此时命题 \hyperref[VI.prop1.8]{1.8}
及推论 \hyperref[VI.cor1.8.1]{1.8.1} 成立;
当两个对象的积在 $E$ 中可表示时,
\hyperref[VI.cor1.8.2]{1.8.2} 也成立.
\end{enonce*}

\setcounter{subsection}{9}
\setcounter{subsubsection}{1}
\subsubsection{}\etiquette[1.9.2]{VI.subsubsec1.9.2}
设 $S$ 是 $E$ 的一个对象,
$f:X\to Y$ 是 $E_{/S}$ 的一个态射.
由 \hyperref[VI.rem1.5.1]{注 1.5} 立即可知,
$f$ 拟紧 (\resp 拟分离, \resp 凝聚),
当且仅当 $E$ 中相应的态射\pageoriginale{}
(由“忘掉 $S$”函子作用于 $f$ 得到) 具有相同性质.
特别地, 若 $f:X\to Y$ 是 $E$ 中的态射,
则在上面取 $S=Y$ 可见,
$f$ 拟紧 (\resp 拟分离, \resp 凝聚),
当且仅当诱导拓扑斯 $E_{/Y}$ 的对象 $(X,f)$
到该拓扑斯的终对象的结构态射拟紧 (\resp 拟分离, \resp 凝聚).
也说 (借助结构态射 $f$)
$X$ \emph{在 $Y$ 上拟紧 (\emph{\resp} 拟分离, \emph{\resp} 凝聚)};
这种说法既可以理解为关于拓扑斯 $E$,
也可以理解为关于诱导拓扑斯 $E_{/Y}$
(其中 $Y$ 被看作终对象).

\subsubsection{}\etiquette[1.9.3]{VI.subsubsec1.9.3}
若拓扑斯 $E$ 的对象 $X$ 在终对象上凝聚,
就称其为\emph{可构造的}.
两个可构造对象之间的每个态射都凝聚
(\hyperref[VI.prop1.8]{1.8} (iv)).
由 $E$ 的可构造对象构成的 $E$ 的满子范畴 $E_{\cons}$
在有限逆极限下稳定,
这立即来自 \hyperref[VI.cor1.8.2]{1.8.2}
及 \hyperref[VI.prop1.8]{1.8} (ii).

\subsubsection{}\etiquette[1.9.4]{VI.subsubsec1.9.4}
\hyperref[VI.def1.7]{1.7} 的概念,
只在具有由拟紧对象构成的生成族的拓扑斯 $E$ 中才真正有意义.
在此情形, 这些概念“在底上局部”
(IV \hyperref[IV.subsubsec8.5.4]{8.5.4}), 现证明如下.

\begin{enonce*}{命题 1.10}\etiquette[1.10]{VI.prop1.10}
假定\pageoriginale{}拓扑斯 $E$ 具有一个由拟紧对象构成的生成子范畴 $C$,
并设 $f:X\to Y$ 是 $E$ 中的一个态射.
\begin{enumerate}
\renewcommand{\theenumi}{\roman{enumi}}
\renewcommand{\labelenumi}{(\theenumi)}
\item 态射 $f$ 拟紧, 当且仅当对满足 $S\in\Ob C$ 的每个态射 $S\to Y$,
对象 $\fprod{X}{S}{Y}$ 拟紧.

\item 设 $Y_i\to Y$, $i\in I$, 为一个覆盖族.
态射 $f$ 拟紧 (\emph{\resp} 拟分离, \emph{\resp} 凝聚),
当且仅当对每个 $i\in I$,
态射 $f_i:X_i=\fprod{X}{Y_i}{Y}\to Y_i$ 具有相同性质.
\end{enumerate}
\end{enonce*}

证明 (i). 必要性由定义立即得到.
证明充分性, 须证明对每个拟紧对象 $Y'$ 及每个态射 $Y'\to Y$,
对象 $X'=\fprod{X}{Y'}{Y}$ 拟紧.
存在各 $S_i$ 属于 $\Ob C$ 的覆盖族 $S_i\to Y'$;
由于 $Y'$ 拟紧, 可以取这一覆盖族为\emph{有限}族.
由假设, 各对象
${X'}_i=\fprod{X}{S_i}{Y}\simeq\fprod{X'}{S_i}{Y'}$
拟紧, 而且它们覆盖 $X'$ (因为各 $S_i$ 覆盖 $Y'$),
故由 \hyperref[VI.prop1.3]{1.3}, $X'$ 拟紧.

证明 (ii). 只须处理拟紧情形,
因为拟分离情形由定义 \hyperref[VI.def1.7]{1.7} 得到,
而凝聚情形由前两种情形合并得到.
必要性已在 \hyperref[VI.prop1.8]{1.8} (ii) 中证明.
只须证明若各 $f_i$ 拟紧, 则 $f$ 拟紧,
\ie 对每个拟紧的 $Y'$ 及每个态射 $Y'\to Y$,
对象 $X'=\fprod{X}{Y'}{Y}$ 拟紧.
由于 $(Y_i\to Y)$ 是覆盖族,\pageoriginale{}
存在 $Y$ 上的对象 $S_{\alpha}$ 覆盖 $Y'$,
且每个复合 $S_{\alpha}\to Y'\to Y$ 通过某个 $Y_i$ 分解.
由于 $C$ 生成, 可以取各 $S_{\alpha}$ 属于 $\Ob C$;
又因 $Y'$ 拟紧, 可以取该覆盖族有限.
于是各 $\fprod{X'}{S_{\alpha}}{Y'}$ 构成 $X'$ 的有限覆盖族,
由 \hyperref[VI.prop1.3]{1.3},
归结为证明每个
$\fprod{X'}{S_{\alpha}}{Y'}=\fprod{X}{S_{\alpha}}{Y}$ 拟紧.
选定 $S_{\alpha}\to Y$ 通过某个 $Y_i$ 的分解后,
所考察的对象也同构于 $\fprod{X_i}{S_{\alpha}}{Y_i}$;
它确实拟紧, 因为 $X_i$ 在 $Y_i$ 上拟紧, 而 $S_{\alpha}$ 拟紧.

\begin{enonce*}{推论 1.11}\etiquette[1.11]{VI.cor1.11}
设 $g:Y\to Z$ 是拓扑斯 $E$ 中的态射,
$(f_i:X_i\to Y)_{i\in I}$ 是以 $Y$ 为靶的一个覆盖态射族.
则:
\begin{enumerate}
\renewcommand{\theenumi}{\roman{enumi}}
\renewcommand{\labelenumi}{(\theenumi)}
\item 假定 $I$ 有限. 若各 $gf_i$ 拟紧, 则 $g$ 也拟紧.

\item 假定各 $f_i$ 拟紧.
若各 $gf_i$ 拟分离, 则 $g$ 也拟分离.
若 $I$ 有限且各 $gf_i$ 凝聚, 则 $g$ 凝聚.
\end{enumerate}
\end{enonce*}

情形 (i) 立即来自定义及 \hyperref[VI.prop1.3]{1.3},
无须对 $E$ 作任何假设.
证明 (ii), 由 (i), 只须处理非括号内情形.
考察图 \hyperref[VI.prop1.8]{1.8} $(\ast)$,
将 $X$ 换成某个 $X_i$.
由于各 $f_i$ 构成的族覆盖,
由它经基变换得到的各 $\fprod{f_i}{\id_Y}{Z}$ 构成的族也覆盖.
因此\pageoriginale{}, 由 \hyperref[VI.prop1.10]{1.10} (ii),
要证明 $\Delta_{Y/Z}$ 拟紧, 只须证明各 $\Gamma_{f_i}$ 拟紧.
由假设, $gf_i$ 拟分离,
故由它经基变换得到的 $\pr_2:\fprod{X_i}{Y}{Z}\to Y$ 也拟分离.
另一方面, $\pr_2\Gamma_{f_i}=f_i$ 由假设拟紧.
因此, 将 \hyperref[VI.prop1.8]{1.8} (iv)
用于 $\Gamma_{f_i}$ 与 $\pr_2$ 的复合, 即知 $\Gamma_{f_i}$ 拟紧, \cqfd

\begin{enonce*}{推论 1.12}\etiquette[1.12]{VI.cor1.12}
设 $(X_i)_{i\in I}$ 是 $E$ 的一个有限对象族,
$X$ 是各 $X_i$ 的余积,
$f_i:X_i\to Y$ 是 $E$ 中的态射,
$f:X\to Y$ 为相应态射.
态射 $f$ 拟紧 (\resp 拟分离, \resp 凝聚),
当且仅当各 $f_i$ 具有相同性质;
在拟分离情形, 可以去掉 $I$ 有限的假设.
\end{enonce*}

充分性是 \hyperref[VI.cor1.11]{1.11} 的特殊情形,
因为典范态射 $X_i\to X$ 构成有限覆盖族且为拟紧态射,
后一点由定义及 \hyperref[VI.cor1.4]{1.4} 立即得到.
必要性来自传递性 \hyperref[VI.prop1.8]{1.8} (i),
因为态射 $X_i\to X$ 甚至凝聚
(由 \hyperref[VI.cor1.8.1]{1.8.1}, 它们拟分离).

\begin{enonce*}{定义 1.13}\etiquette[1.13]{VI.def1.13}
设 $X$ 是拓扑斯 $E$ 的一个对象.
如果对 $E$ 的每个拟紧对象 $S$,
从 $S$ 到 $X$ 的每个态射都拟紧,
\ie (\hyperref[VI.def1.7]{1.7})
如果对 $E$ 的任意两个拟紧对象 $S,T$
及态射 $S\to X$, $T\to X$,
纤维积 $\fprod{S}{T}{X}$ 总拟紧,
就称 $X$ 为\emph{拓扑斯 $E$ 的拟分离对象}.
若 $X$ 拟紧且拟分离, 就称 $X$ 为拓扑斯 $E$ 的凝聚对象.
\end{enonce*}

\begin{enonce*}{命题 1.14}\etiquette[1.14]{VI.prop1.14}
设\pageoriginale{} $f:X\to Y$ 是拓扑斯 $E$ 中的态射.
\begin{enumerate}
\renewcommand{\theenumi}{\roman{enumi}}
\renewcommand{\labelenumi}{(\theenumi)}
\item 若 $Y$ 拟紧, 则 $f$ 拟紧蕴含 $X$ 拟紧.
若 $Y$ 拟分离, 则 $X$ 拟紧蕴含 $f$ 拟紧.
若 $Y$ 凝聚, 则 $f$ 拟紧当且仅当 $X$ 拟紧.

\item 若 $Y$ 拟紧 (\emph{\resp} 拟分离, \emph{\resp} 凝聚),
且 $f$ 拟紧 (\emph{\resp} 拟分离, \emph{\resp} 凝聚),
则 $X$ 拟紧 (\emph{\resp} 拟分离, \emph{\resp} 凝聚).
\end{enumerate}
\end{enonce*}

(i) 的前两个断言分别直接包含在拟紧态射的定义
(\hyperref[VI.def1.7]{1.7})
与拟分离对象的定义 (\hyperref[VI.def1.13]{1.13}) 中;
第三个由前两个合并得到.
(ii) 的第一种情形只是重述, 第三种仍由前两种合并得到.
只须处理拟分离情形.
设 $S$ 是 $E$ 的一个拟紧对象,
$g:S\to X$ 为态射; 证明 $g$ 拟紧.
由于 $Y$ 拟分离, $fg$ 拟紧;
又因 $f$ 拟分离, 由 (\hyperref[VI.prop1.8]{1.8} (iii)) 可知
$g$ 拟紧, \cqfd

\begin{enonce*}{推论 1.15}\etiquette[1.15]{VI.cor1.15}
$E$ 的拟分离对象的每个子对象都拟分离.
设 $(X_i)_{i\in I}$ 是 $E$ 的一族对象, $I\in 𝒰$;
则 $X=\amalg_iX_i$ 拟分离 (\emph{\resp} 凝聚),
当且仅当各 $X_i$ 也如此.
\end{enonce*}

第一个\pageoriginale{}断言来自 \hyperref[VI.prop1.14]{1.14} (ii)
及 \hyperref[VI.cor1.8.1]{1.8.1}, 但由定义也可立即得到.
对第二个断言的非括号内情形,
由于各 $X_i\to X$ 都是单态射, 只须证明充分性,
而这容易由定义及 \hyperref[VI.prop1.3]{1.3} 得到;
括号内情形来自刚处理的情形及 \hyperref[VI.cor1.4]{1.4}.

\begin{enonce*}{推论 1.16}\etiquette[1.16]{VI.cor1.16}
在 \hyperref[VI.prop1.10]{1.10} (ii) 的条件下,
若各 $Y_i$ 凝聚, 则 $f$ 拟紧,
当且仅当各 $X_i=\fprod{X}{Y_i}{Y}$ 是 $E$ 的拟紧对象.
\end{enonce*}

事实上, 由判据 \hyperref[VI.prop1.10]{1.10} (ii),
须表达各 $f_i:X_i\to Y_i$ 拟紧;
由 \hyperref[VI.prop1.14]{1.14} (i),
这等价于各 $X_i$ 拟紧.

\begin{enonce*}{推论 1.17}\etiquette[1.17]{VI.cor1.17}
设 $E$ 是一个具有由拟紧对象构成的生成族的拓扑斯,
$(g_i:Y_i\to Y)_{i\in I}$ 是 $E$ 中的覆盖族 (\resp 有限覆盖族),
且各 $Y_i$ 凝聚.
对象 $Y$ 拟分离 (\resp 凝聚),
当且仅当对每个 $i\in I$, $g_i$ 拟紧;
或等价地, 当且仅当对每对指标 $i,j\in I$,
$\fprod{Y_i}{Y_j}{Y}$ 为 $E$ 的拟紧对象.
\end{enonce*}

括号内情形立即来自非括号内情形及 \hyperref[VI.prop1.3]{1.3}.
第一个条件的必要性由定义立即得到; 证明其充分性.
设 $f:X\to Y$ 为态射, 且\pageoriginale{} $X$ 拟紧;
证明 $f$ 拟紧.
由于各 $g_i$ 拟紧, 各 $X_i=\fprod{X}{Y_i}{Y}$ 拟紧,
由 \hyperref[VI.cor1.16]{1.16} 可知 $f$ 拟紧.
最后, 用同一判据 \hyperref[VI.cor1.16]{1.16}
表达各 $g_i$ 的拟紧性,
就得到 \hyperref[VI.cor1.17]{1.17} 的第二个判据.
特别地:

\begin{enonce*}{推论 1.17.1}\etiquette[1.17.1]{VI.cor1.17.1}
设 $E$ 如 \hyperref[VI.cor1.17]{1.17},
$X$ 是 $E$ 的一个凝聚对象,
$\Rr\rightrightarrows X$ 是 $X$ 中的一个等价关系,
$Y=X/\Rr$. 则下列条件等价:
\begin{enumerate}
\renewcommand{\theenumi}{\roman{enumi}}
\renewcommand{\labelenumi}{(\theenumi)}
\item $Y$ 凝聚;
\item $f:X\to Y$ 拟紧;
\item $\Rr$ 拟紧.
\end{enumerate}
\end{enonce*}

\begin{enonce*}{推论 1.18}\etiquette[1.18]{VI.cor1.18}
设 $E$ 是具有由拟紧对象构成的生成子范畴 $C$ 的拓扑斯,
$Y$ 是 $E$ 的一个对象.
对象 $Y$ 拟分离,
当且仅当每个满足 $X\in\ob C$ 的态射 $X\to Y$ 拟紧.
\end{enonce*}

必要性由定义立即得到,
因为由假设各对象 $X\in\ob C$ 都拟紧.
证明充分性, 设 $Z$ 是 $E$ 的一个拟紧对象;
证明每个态射 $Z\to Y$ 拟紧.
由 \hyperref[VI.prop1.10]{1.10} (i),
只须验证对每个满足 $X\in\ob C$ 的 $X\to Y$,
对象 $\fprod{Z}{X}{Y}$ 拟紧.
这来自各 $X\to Y$ 由假设拟紧, 而 $Z$ 拟紧这一事实.

\begin{enonce*}{推论 1.19}\etiquette[1.19]{VI.cor1.19}
设\pageoriginale{} $E$ 是一个具有由拟紧对象构成的生成子范畴的拓扑斯,
$(g_i:Y_i\to Y)_{i\in I}$ 是 $E$ 中的覆盖族,
各 $Y_i$ 拟分离且各 $g_i$ 拟紧. 则 $Y$ 拟分离.
\end{enonce*}

事实上, 设 $X$ 是 $E$ 的拟紧对象;
证明每个态射 $f:X\to Y$ 拟紧.
由 \hyperref[VI.prop1.10]{1.10} (ii),
只须证明诱导态射 $f_i:X_i=\fprod{X}{Y_i}{Y}\to Y_i$ 拟紧.
由于 $g_i$ 拟紧, $X_i$ 拟紧 (因为 $X$ 拟紧),
又因 $Y_i$ 拟分离, 可知 $f_i$ 拟紧.

\begin{enonce*}[remark]{注 1.20.1}\etiquette[1.20.1]{VI.rem1.20.1}
习惯于代数几何中“拟紧”与“拟分离”术语的读者,
会预期这里考察的“绝对”概念 (关于拓扑斯 $E$ 的对象)
和“相对”概念 (关于 $E$ 的箭头) 之间还有某些其他关系;
例如, 若 $f:X\to Y$ 是 $E$ 中的态射,
且 $X$ 拟紧并拟分离, 则 $f$ 拟分离.
这类性质只在对拓扑斯 $E$ 施加限制性的有限性条件时才成立,
它们比 \hyperref[VI.prop1.10]{1.10} 中考虑的条件更强;
下一节将研究这些条件.
\end{enonce*}

\setcounter{subsection}{20}
\setcounter{subsubsection}{1}
\subsubsection{}\etiquette[1.20.2]{VI.subsubsec1.20.2}
由 \hyperref[VI.rem1.5.1]{1.5} 还可知,
拓扑斯 $E$ 的对象 $X$ 是否拟分离,
只依赖于诱导拓扑斯 $E_{/X}$,
并意味着该拓扑斯的终对象拟分离,
\ie $E_{/X}$ 中两个拟紧对象的积拟紧.


% END INCLUDED SGA4_English_current_buildable_source_20260906/source/06/drafts/SGA4_VI_LINES_391_677_English.texfrag


% BEGIN INCLUDED SGA4_English_current_buildable_source_20260906/source/06/drafts/SGA4_VI_LINES_678_1095_English.texfrag
\subsection{景的情形}\etiquette[1.21]{VI.subsec1.21}\pageoriginale

设 $C$ 是一个 $𝒰$-景,
$E$ 是 $C$ 上的 $𝒰$-层拓扑斯,
$f$ 是 $C$ 的一个箭头, $X$ 是 $C$ 的一个对象.
若 $\epsilon(f)$ 是拓扑斯 $E$ 的拟紧态射 (\resp 拟分离态射),
就称 $f$ 为\emph{景 $C$ 的拟紧态射} (\resp 拟分离态射);
同样, 若 $\epsilon(X)$ 是拓扑斯 $E$ 的拟分离对象,
就称 $X$ 为\emph{景 $C$ 的拟分离对象}
(比较 \hyperref[VI.prop1.2]{1.2}).
当 $C$ 是配备其典范拓扑的 $𝒰$-拓扑斯时,
典范函子 $\epsilon:C\to E$ 为范畴等价,
因而上述定义与先前的\hyperref[VI.def1.7]{定义 1.7}
及 \hyperref[VI.def1.13]{1.13} 相容.
还要注意, 将宇宙 $𝒰$ 换成另一个使 $C$ 也为 $𝒱$-景的宇宙 $𝒱$ 时,
所考察的定义不变, 至少在 $C$ 具有一个由拟紧对象构成的拓扑生成族时如此.
为此立即归结为情形 $𝒰\subset 𝒱$,
于是 $E$ 是一个 $𝒱$-景;
并须证明对 $𝒰$-拓扑斯 $E$ 的每个箭头 $f$,
$f$ 拟紧 (\resp 拟分离), 当且仅当它在 $E$ 上的 $𝒱$-层拓扑斯 $E'$
中定义一个拟紧箭头;
同样, 须证明对 $E$ 的每个对象 $X$,
$X$ 拟分离, 当且仅当 $E'$ 中的相应对象拟分离.
对拟紧对象的相应断言, 我们已经知道
(\hyperref[VI.prop1.2]{1.2}), 而且无须对 $E$ 作任何假设.
情形“$f$ 拟紧”由此得到,
只须将判据 \hyperref[VI.prop1.10]{1.10} (i)
先后在 $E$ 及\pageoriginale{} $E'$ 中
应用于 $E$ 的同一个由拟紧对象构成的生成族 $(X_i)_{i\in I}$.
这也给出了情形“$f$ 拟分离”,
即相应对角态射拟紧的情形.
最后, 在“$X$ 拟分离”的情形,
由判据 \hyperref[VI.cor1.18]{1.18},
仍归结为“$X$ 拟紧”的情形.

\subsection{}\etiquette[1.22]{VI.subsec1.22}
{\textit{\textbf{例.}}}
回到配备拓扑 $T_i$ 的范畴 $(\Sch)$
这一例子 \hyperref[VI.subsubsec1.6.2]{1.6.2}.
对适当的 $𝒱$, 它是一个 $𝒱$-景,
且具有以仿射概形子范畴为拓扑生成范畴,
其中的对象拟紧.
请读者验证, 此景的箭头 $f:X\to Y$ 拟紧 (\resp 拟分离),
当且仅当它是通常意义下的拟紧 (\resp 拟分离) 概形态射
(\cite{sga4.exp.VI.bib.3} 或 \cite{sga4.exp.VI.bib.4});
同样, 景的对象拟紧 (\resp 拟分离),
当且仅当它是 \loccit{} 通常意义下的拟紧 (\resp 拟分离) 概形.
实际上, 用我们这里的说法, 称\emph{凝聚概形态射、凝聚概形},
代替“拟紧且拟分离的概形态射”、“拟紧且拟分离的概形”, 是恰当的;
\EGA I, 第 $2^{\text{nd}}$ 版中也这样做.

\setcounter{subsection}{22}
\subsubsection{}\etiquette[1.22.1]{VI.subsubsec1.22.1}

对给定的概形 $X\in 𝒰$,
还可以考察与底拓扑空间相伴的拓扑斯 $\Top(X)$,
或与属于 $𝒰$ 的平展 $X$-概形
(\resp 在 $X$ 上局部\pageoriginale{}有限呈示的概形)
所成的景配备平展拓扑 (\resp fppf 拓扑) 后相伴的拓扑斯
$\Top(X_{\et})$ (\resp $\Top(X_{\text{fppf}})$).
这是一个具有由凝聚对象构成的生成族的 $𝒰$-拓扑斯,
这些对象对应于定义景中的仿射对象.

明确这一点后, 此拓扑斯的终对象拟紧 (\resp 拟分离, \resp 凝聚),
当且仅当概形 $X$ 拟紧 (\resp 拟分离, \resp 凝聚).
同样, $X$ 上展空间景
(\resp $X$ 上平展概形景 $X_{\text{et}}$,
\resp $X$ 上局部有限呈示概形景 $X_{\text{fppf}}$) 的态射
拟紧 (\resp 拟分离, \resp 凝聚),
当且仅当它是 \loccit{} 通常意义下的拟紧
(\resp 拟分离, \resp 凝聚) 概形态射.

\begin{enonce*}{定理 1.23}\etiquette[1.23]{VI.thm1.23}
设 $E$ 是一个 $𝒰$-拓扑斯, $X$ 是 $E$ 的一个对象.
\begin{enumerate}
\renewcommand{\theenumi}{\roman{enumi}}
\renewcommand{\labelenumi}{(\theenumi)}
\item $X$ 拟紧的充要条件是,
对 $E$ 的每个滤过归纳系 $(Y_i)_{i\in I}$, $I\in 𝒰$,
自然映射
\begin{equation*}
\varinjlim  \Hom(X, Y_i) \to \Hom(X, \varinjlim Y_i)\tag{1.23.1}
\end{equation*}\etiquette[1.23.1]{VI.eq1.23.1}
为单射.

\item 假定 $E$ 具有由拟紧对象构成的生成子范畴 $C$.
$X$ 凝聚的必要条件是,
对 \emph{(i)} 中的每个数据,
映射 \emph{\hyperref[VI.eq1.23.1]{(1.23.1)}} 为双射.
\end{enumerate}
\end{enonce*}

\begin{enumerate}
\renewcommand{\theenumi}{\roman{enumi}}
\renewcommand{\labelenumi}{(\theenumi)}
\item \underline{必要性}.
假定\pageoriginale{} $X$ 拟紧.
须证明, 若 $i\in I$ 及 $f_i,g_i:X\rightrightarrows Y_i$
使 $f_i,g_i$ 与 $Y_i\to Y=\varinjlim Y_j$ 的复合 $f,g$ 相等,
则存在 $j\geq i$,
使它们与 $Y_i\to Y_j$ 的复合 $f_j,g_j$ 就已相等.
由于 $E$ 中的滤过归纳极限与有限投射极限交换,
特别与核交换, 有
$$
\Ker (f, g) = \mathop{\lim}_{j\geq l} \Ker (f_j, g_j){\quoi}.
$$
由假设 $\Ker(f,g)=X$,
故 $X$ 是由其子对象 $\Ker(f_j,g_j)$ 构成的滤过递增族的极限.
由于 $X$ 拟紧, 它等于其中一个子对象;
故存在 $j\geq i$ 使得 $f_j=g_j$.

\underline{充分性}.
要证明 $X$ 拟紧,
等价地 (\hyperref[VI.subsubsec1.5.2]{1.5.2})
须证明对 $X$ 的每个 $\varinjlim$ 为 $X$ 的滤过递增子对象族
$(X_i)_{i\in I}$, 存在 $i\in I$ 使得 $X_i=X$.
由于 $E$ 中每个单态射都有效 (II \hyperref[II.sec4]{4.0}), 有
$$
X_i = \Ker (f_i, g_i : X \rightrightarrows X \amalg_{X_{i}} X){\quoi}.
$$
设 $Y_i=X\amalg_{X_i}X$,
并考察由各 $Y_i$ 构成的归纳系.
若 $Y$ 为其归纳极限, 则
$Y\simeq X\amalg_{\varinjlim X_i}X=X\amalg_X X=X$.
因此, 对固定的 $i\in I$,
$f_i,g_i$ 给出 $\Hom(X,Y)$ 的同一元素,
故由假设, 存在 $j\geq i$ 使得 $f_j=g_j$,
\ie $X_j=X$, \cqfd

\item 只须\pageoriginale{}证明,
在所给条件下, 每个态射 $X\to Y$
都来自适当 $i\in I$ 的某个态射 $X\to Y_i$.
由于 $C$ 生成 $E$, 各 $Y_i$ 覆盖 $Y$,
可以用 $C$ 的对象 $X_{\alpha}$ 覆盖 $X$,
使每个复合 $X_{\alpha}\to X\to Y$ 通过某个 $Y_i$ 分解.
由于 $X$ 拟紧, 可以假定各 $X_{\alpha}$ 构成的族有限,
从而各态射 $X_{\alpha}\to Y$ 通过一个固定的 $Y_i$
分解为态射 $g_{\alpha i}:X_{\alpha}\to Y_i$.
由于 $X$ 拟分离, 各 $\fprod{X_{\alpha}}{X_{\beta}}{X}$ 拟紧.
此外, 对每个 $(\alpha,\beta)$,
$\fprod{X_{\alpha}}{X_{\beta}}{X}$ 上的 $\pr_1$ 和 $\pr_2$
与 $X_{\alpha},X_{\beta}\to Y_i$ 的复合
\[
\xymatrix{
\fprod{X_{\alpha}}{X_{\beta}}{X} \ar[rd] \ar[dd]& &\\
                            &X_{\alpha} \ar[ld] \ar[rrrr]&   & & & Y_i
\ar[lld]\\
X    \ar[rrr]                       &          & & Y  &
}
\]
与 $Y_i\to Y$ 复合后相等.
由 (i), 因而存在 $j\geq i$,
使它们与 $Y_i\to Y_j$ 的复合就已相等.
由于对 $(\alpha,\beta)$ 所成的集合有限,
可以不依赖于这一对选取 $j$.
因此, 态射 $g_{\alpha j}$
由一个态射 $g_j:X\to Y_j$ 复合而来;
后者与 $Y_j\to Y$ 的复合又为给定态射 $X\to Y$,
因为与覆盖映射 $X_{\alpha}\to X$ 复合后确实如此.
这就完成了证明.
\end{enumerate}

\begin{enonce*}[remark]{注 1.23.2}\etiquette[1.23.2]{VI.rem1.23.2}
(ii) 中对 $E$ 的假设可换为\emph{对 $X$ 的}一个附加假设,
即 $X$ 具有一个由拟紧对象组成的覆盖族基本系.
\end{enonce*}

\begin{enonce*}{推论 1.24}\etiquette[1.24]{VI.cor1.24}
设\pageoriginale{} $E$ 是一个拓扑斯,
$C$ 是 $E$ 的一个由凝聚对象构成的严格满子范畴.
利用 $E$ 中的 $𝒰$-归纳极限可表示这一事实,
得到 \emph{(I \hyperref[I.sec8.6.1]{8.6.1})} 典范函子
\begin{equation*}
\Ind (C) \to E{\quoi}.\tag{1.24.1}
\end{equation*}\etiquette[1.24.1]{VI.eq1.24.1}
这个函子全忠实.
若 $E$ 具有一个由凝聚对象构成的生成子范畴,
且 $C$ 对直因子稳定, 则下列条件等价:
\begin{enumerate}
\renewcommand{\theenumi}{\roman{enumi}}
\renewcommand{\labelenumi}{(\theenumi)}
\item 函子 \emph{\hyperref[VI.eq1.24.1]{(1.24.1)}} 本质满,
\ie 它为范畴等价.

\item $C$ 是 $E$ 的生成子范畴, 且对有限归纳极限稳定.

\item 范畴 $C$ 等于 $E$ 中由 $E$ 的所有凝聚对象构成的满子范畴 $E_{\coh}$,
而且 $E$ 的对象凝聚的必要条件
\emph{\hyperref[VI.thm1.23]{1.23} (ii)} 也充分.
\end{enumerate}
\end{enonce*}

设 $E_{PF}$ 为 I \hyperref[I.cor8.7.6]{8.7.6} 中定义的 $E$ 的满子范畴.
由 I \hyperref[I.prop8.7.5]{8.7.5} a),
陈述 \hyperref[VI.thm1.23]{1.23} (ii) 等价于
$$
C \subset E_{\coh} \subset E_{PF}{\quoi},
$$
中的第二个包含;
而复合包含 $C\supset E_{PF}$ 意味着
\hyperref[VI.eq1.24.1]{(1.24.1)} 全忠实.
由于\pageoriginale{}由假设, $E$ 的子范畴 $E_{\coh}$ 生成,
$E_{PF}$ 更是如此;
此外 $E$ 中有限投射极限可表示,
故可以应用 I \hyperref[VI.cor8.7.7]{8.7.7}.
\hyperref[VI.cor1.24]{1.24} 的条件 (iii)
也可写为 $C=E_{\coh}=E_{PF}$,
并等价于条件 $C=E_{PF}$,
这正是 \loccit{} 的条件 (ii bis)
(这里用到由假设, $E_{\coh}$
因而 $E_{PF}$ 是 $E$ 的生成子范畴).
同样, \hyperref[VI.cor1.24]{1.24} 的条件 (ii)
正是 \loccit{} 的条件 (iii bis).
因此, \loccit{} 中条件 (i)、(ii bis) 与 (iii bis) 的等价性
给出 \hyperref[VI.cor1.24]{1.24} 中条件 (i)、(ii)、(iii) 的等价性, \cqfd

该证明 (更确切地说,
引用 I \hyperref[I.cor8.7.7]{8.7.7} (ii bis)) 还给出了下述推论的 a) 部分.

\begin{enonce*}{推论 1.24.2}\etiquette[1.24.2]{VI.cor1.24.2}
设 $E$ 是一个具有由凝聚对象构成的生成子族的拓扑斯,
$E_{PF}$ 是由 $E$ 中满足如下条件的对象 $X$ 构成的严格满子范畴:
由 $X$ 表示的协变函子 $\Hom(X,-)$ 与滤过归纳极限交换.
则:
\begin{enumerate}
\renewcommand{\theenumi}{\alph{enumi}}
\renewcommand{\labelenumi}{\theenumi)}
\item 自然函子
$$
\Ind (E_{PF}) \to E
$$
为范畴等价.

\item $E$ 的对象 $X$ 属于 $E_{PF}$ 的充要条件是,
它同构于某个双箭头 $X_1\rightrightarrows X_0$ 的余核\pageoriginale{}的一个直因子
($=$ 投影子的像 (\hyperref[I.sec10.6]{I, 10.6})),
其中 $X_1$ 和 $X_0$ 凝聚.
\item $E$ 的子范畴 $E_{PF}$ 对有限归纳极限稳定,
并且等价于一个小范畴.
\end{enumerate}
\end{enonce*}

还须证明 b) 和 c).
c) 的第一个断言立即来自 I \hyperref[I.prop2.8]{2.8},
并由 \hyperref[VI.thm1.23]{1.23} (ii)
蕴含 b) 的“若”方向.
证明“仅当”方向, 使用 \hyperref[VI.thm1.23]{1.23} (i),
它表明 $X$ 拟紧,
从而由于 $E_{\coh}$ 生成 $E$ 且对有限和稳定,
存在满态射 $X_0\to X$, 其中 $X_0$ 凝聚.
设 $R=\fprod{X_0}{X_0}{X}$,
则 $X$ 等同于 $R\rightrightarrows X_0$ 的余核.
由 a), 有 $R=\varinjlim R_i$,
为各 coker $(R_i\rightrightarrows X_0)=X_i$ 的滤过归纳极限.
由假设 $X\in\Ob E_{PF}$,
对充分大的 $i$, 态射 $u_i:X'_i\to X$ 有截面 $v:X\to X'_i$.
设 $w=vu_i:X'_i\to X'_i$,
则 $w^2=w$, 且
$X\simeq\Coker(w,\id_{X'_i}:X'_i\rightrightarrows X'_i)$.
由于 $R_i$ 仍拟紧, 用 $X_1\to R_i$ 覆盖它,
其中 $X_i$ 凝聚,
从而 $X_i\simeq\Coker(X_1\rightrightarrows X_0)$,
这就证明了 b).

最后, 由于 $E_{PF}\subset E_{\text{qu.cpct}}$,
c) 的最后一个断言来自 \hyperref[VI.subsubsec1.5.3]{1.5.3}.

\begin{enonce*}{推论 1.25}\etiquette[1.25]{VI.cor1.25}
设 $E$ 是一个拓扑斯,
由凝聚对象构成的 $E$ 的满子范畴 $C$ 生成.
考察典范函子
$$
\varphi : \Ind(c) \to E{\quoi}.
$$
下列\pageoriginale{}条件等价:
\begin{itemize}
\item[(i)] 函子 $\varphi$ 本质满,
\ie $E$ 的每个对象都是凝聚对象的滤过归纳极限.

\item[(i bis)] 函子 $\varphi$ 为范畴等价.

\item[(ii)] $E$ 中凝聚对象的每个有限归纳极限都是凝聚对象.

\item[(ii bis)] $E$ 中凝聚对象的双箭头的余核是凝聚对象.

\item[(iii)] \emph{\hyperref[VI.thm1.23]{1.23} (ii)} 的逆命题成立.
\end{itemize}
\end{enonce*}

这是 \hyperref[VI.cor1.24]{1.24} 的特殊情形,
因为 $C$ 对有限和稳定 (\hyperref[VI.cor1.4]{1.4}),
从而 (ii) 与 (ii bis) 等价, 并且有如下引理.

\begin{enonce*}{引理 1.25.1}\etiquette[1.25.1]{VI.lem1.25.1}
设 $E$ 是一个拓扑斯.
$E$ 的凝聚对象 $X$ 的每个直因子 $X'$ 都凝聚.
\end{enonce*}

事实上, $X'$ 作为 $X$ 的商而拟紧
(\hyperref[VI.prop1.3]{1.3}),
作为 $X$ 的子对象而拟分离 (\hyperref[VI.cor1.15]{1.15}).

\begin{enonce*}{推论 1.26}\etiquette[1.26]{VI.cor1.26}
设 $E$ 是满足等价条件 \emph{\hyperref[VI.cor1.25]{1.25}} 的拓扑斯.
则每个诱导拓扑斯也如此.
\end{enonce*}

这由判据 (ii) 立即得到,
只须考虑 \hyperref[VI.subsubsec1.20.2]{1.20.2},
以及由凝聚对象构成的 $E$ 的生成子范畴的存在性,
显然蕴含诱导拓扑斯的相同性质这一事实.

\begin{enonce*}[remark]{注 1.27}\etiquette[1.27]{VI.rem1.27}
在\pageoriginale{}满足 \hyperref[VI.cor1.25]{1.25} 的条件的重要拓扑斯例子中,
可以列举 Noether 拓扑斯 (\hyperref[VI.subsec2.14]{2.14}),
以及概形 $X$ 的 Zariski 或平展拓扑斯
(\hyperrefext[VII][VII.sec1]{1}) (\cf IX, 注, 第 42 页).
不过, 即使拓扑斯 $E$ \emph{凝聚}
(\hyperref[VI.def2.3]{2.3})
(\ie 凝聚对象的满子范畴 $C$ 生成且对有限投射极限稳定),
$C$ 也未必\emph{完美（perfect）}
(\hyperref[VI.def2.9.1]{2.9.1}),
\ie 未必满足 \hyperref[VI.cor1.25]{1.25} 的等价条件;
\cf 下文 \hyperref[VI.exer1.28]{1.28}.
\end{enonce*}

\begin{enonce*}[remark]{习题 1.28}\etiquette[1.28]{VI.exer1.28}
设 $E$ 是一个拓扑斯,
$p_1,p_2:X_1\rightrightarrows X_0$ 是 $E$ 中的双箭头.
定义 $X_0\times X_0$ 的子对象的递增序列 $R_n$ $(n\geq0)$,
使得 $R_0=\Sup(\text{Im}(p_1,p_2),\text{Im}(p_2,p_1))$,
并且对每个 $n\geq1$,
$R_n=\pr_{13}(\pr_{12}^{-1}(R_{n-1}),\pr_{23}^{-1}(R_{n-1}))$,
其中 $\pr_{ij}$ $(1\leq i<j\leq3)$
表示从 $X_0$ 的三重积到二重积的显然投影.
\begin{enumerate}
\renewcommand{\theenumi}{\alph{enumi}}
\renewcommand{\labelenumi}{\theenumi)}
\item 证明这样得到 $X_0\times X_0$ 的一个递增子对象序列,
$R=\varinjlim R_n$ 是 $X_0$ 中的等价图式
(I \hyperref[I.sec10.4]{10.4}), 并且典范态射
$$
X = \Coker (p_1, p_2) \to X_0 /R
$$
为同构.

\item 推出, 若 $X_0$ 拟紧且 $X=\Coker(p_1,p_2)$ 拟分离
(从而凝聚), 则序列 $(R_n)$ 稳定.
反过来\pageoriginale{}, 若这一序列稳定,
且假定 $X_0$ 凝聚、$X_1$ 拟紧, 则 $X$ 凝聚.
(证明最后一点, 写成
$R_n\simeq\fprod{(R_{n-1},\pr_2)}{(R_{n-1},\pr_1)}{X_0}$,
推出对每个 $n$, $R_n$ 在 $X_0$ 上拟紧,
然后应用 \hyperref[VI.cor1.19]{1.19}.)

\item 构造两个集合映射 $p_1,p_2:X_1\rightrightarrows X_0$ 的例子,
使序列 $(R_n)$ 不稳定.

\item 设 $C$ 是一个小范畴,
其中有限归纳极限与有限投射极限均可表示,
且每个态射都分解为一个有效满态射后接一个单态射.
对 $C$ 的双箭头 $(p_1,p_2)$ 重复各 $R_n$ 的构造.
给 $C$ 配备典范拓扑,
证明典范函子 $\epsilon:C\to C^{\sim}$ 正合,
从而与各 $R_n$ 的形成交换.

\item 取 $C$ 为 $(\Ens)$ 的一个小满子范畴,
它在同构意义下对有限投射极限与有限归纳极限稳定,
并含有 c) 中的集合 $X_0,X_1$.
推出拓扑斯 $C^{\sim}$
(它是\emph{凝聚}拓扑斯 (\hyperref[VI.def2.3]{2.3}),
\ie 由其凝聚对象子范畴生成, 而该子范畴对有限投射极限稳定)
不满足 \hyperref[VI.cor1.25]{1.25} 的等价条件.

\item 设 $k$ 是一个域,
$C$ 是 $k$ 上局部有限呈示概形的 fppf 景 (\SGA 3 IV 6.3).
证明存在一个 $k$ 上的约化仿射概形 $X$,
具有一个阶无限的自同构 $f$
(例如取仿射平面 $E^2$, 配备自同构 $f(x)=x$, $f(y)=x+y$).\pageoriginale{}
证明若 $p_1,p_2:X_1\rightrightarrows X_0$ 是 $C$ 中的双箭头,
使其在集合范畴中的相应双箭头
$(\Ens) p_1(\bar{k}):p_2(\bar{k}):
  X_1(\bar{k})\rightrightarrows X_0(\bar{k})$
给出不稳定序列 $(R_n)$,
则 $p_1,p_2$ 定义的 $C^{\sim}$ 中的双箭头也如此.
因此, 若 $X_0$ 在 $k$ 上有限型,
则 $C^{\sim}$ 中 $(p_1,p_2)$ 的余核不凝聚.
推出 $C^{\sim}$ 中对 $(f,\id_X)$ 的余核不凝聚,
从而拓扑斯 $C^{\sim}$ 不满足 \hyperref[VI.cor1.25]{1.25} 的等价条件.
更一般地推出, 若 $S$ 为非空概形,
$C$ 是 $S$ 上局部有限呈示概形的 fppf 景,
则拓扑斯 $C^{\sim}$ 不满足 \hyperref[VI.cor1.25]{1.25} 的条件.
(须注意, 与表面印象相反,
$C^{\sim}$ 为 Noether (\hyperref[VI.def2.11]{2.11})
只有在 $S$ 为空时才可能,
这来自前面的讨论及 \hyperref[VI.subsec2.14]{2.14}.)
\end{enumerate}
\end{enonce*}

\begin{enonce*}{定义 1.30}\etiquette[1.30]{VI.def1.30}
设 $E$ 是一个拓扑斯.
$E$ 的对象 $X$ 称为拓扑斯 $E$ 的\emph{预 Noether 对象（pre-Noetherian object）}\footnote{更强的 \emph{Noether} 对象概念在下文
\hyperref[VI.def2.11]{2.11} 给出.},
是指它满足下列两个等价条件:
\begin{enumerate}
\renewcommand{\theenumi}{\roman{enumi}}
\renewcommand{\labelenumi}{(\theenumi)}

\item $X$ 的每个子对象都拟紧.

\item $X$ 的每个递增子对象序列都稳定.
\end{enumerate}
\end{enonce*}

\setcounter{subsection}{30}
\setcounter{subsubsection}{0}
\subsubsection{}\etiquette[1.30.1]{VI.subsubsec1.30.1}
条件 (i) 与 (ii) 的等价性是显然的.
注意, 这些条件仍然只依赖于诱导拓扑斯 $E_{/X}$,
甚至只依赖于 $E_{/X}$ 的开子集有序集,
而它同构于 $X$ 的子对象有序集.
这些条件在过渡到 $X$ 的子对象时稳定.



% END INCLUDED SGA4_English_current_buildable_source_20260906/source/06/drafts/SGA4_VI_LINES_678_1095_English.texfrag


% BEGIN INCLUDED SGA4_English_current_buildable_source_20260906/source/06/drafts/SGA4_VI_LINES_1096_1382_English.texfrag
\subsection{}\etiquette[1.31]{VI.subsec1.31}
仿照\pageoriginale{} \hyperref[VI.prop1.3]{1.3}
及 \hyperref[VI.cor1.4]{1.4}, 可证明下列事实:
若 $(X_i\to X)$ 是有限覆盖族, 且各 $X_i$ 为预 Noether,
则 $X$ 为预 Noether;
特别地, $E$ 的预 Noether 对象的商仍为预 Noether.
若 $(X_i)_{i\in I}$ 是 $E$ 的一族对象,
则其余积 $X$ 是 $E$ 的预 Noether 对象,
当且仅当除有限多个外的 $X_i$ 都同构于 $∅_E$,
且所有 $X_i$ 都为预 Noether.

\subsection{}\etiquette[1.32]{VI.subsec1.32}
显然, 拓扑斯 $E$ 的预 Noether 对象拟紧.
当 $E$ 具有一个由预 Noether 对象构成的生成族时,
逆命题成立 (这立即来自 \hyperref[VI.subsec1.31]{1.31}):
此时 $E$ 的预 Noether 对象恰为其拟紧对象.

假定 $E$ 具有一个由拟紧对象构成的生成族 $(X_i)_{i\in I}$.
则下列三个条件等价:
\begin{enumerate}
\renewcommand{\theenumi}{\roman{enumi}}
\renewcommand{\labelenumi}{(\theenumi)}
\item 各 $X_i$ 为预 Noether.

\item $E$ 的每个拟紧对象都为预 Noether.

\item $E$ 中的每个单态射都拟紧.
\end{enumerate}
\noindent
(我们刚已证明 (i) 与 (ii) 等价;
(ii) $\Rightarrow$ (iii) 由定义立即得到,
而 (iii) $\Rightarrow$ (ii) 来自定义 \hyperref[VI.def1.30]{1.30} (i).)

注意\pageoriginale{}, (iii) 蕴含 \emph{$E$ 的每个态射都拟分离};
从而 (\hyperref[VI.prop1.14]{1.14} (ii))
$E$ 中位于 $E$ 的某个拟分离对象之上的每个对象本身也拟分离.

\begin{enonce*}[remark]{例 1.33}\etiquette[1.33]{VI.exem1.33}
(\emph{群的分类拓扑斯}).
设 $G$ 是一个 $\in 𝒰$ 的群,
$B_G$ 为其分类拓扑斯 (IV \hyperref[IV.subsec2.4]{2.4}),
\ie 由配备 $G$ 的左作用的 $\in 𝒰$ 的集合构成的范畴.
$B_G$ 的非空连通对象
(IV \hyperref[IV.subsubsec4.3.5]{4.3.5})
是 $G$ 传递作用的非空 $G$-集合;
$B_G$ 的每个对象 $E$ 都是其连通分支的余积,
这些连通分支为 $E$ 的极小非空子对象
(即 $G$ 在 $E$ 中的轨道).
因此, $E$ 是 $B_G$ 的拟紧对象,
当且仅当 $G$ 在 $E$ 中的轨道集 $\pi_0(E)$ 有限;
此时 $E$ 甚至是 $B_G$ 的预 Noether 对象.
这些对象 (甚至单个非空连通对象 $G_s$)
构成 $B_G$ 的生成子范畴.
$B_G$ 的终对象 $e$ 拟分离,
\ie $B_G$ 中两个拟紧对象的积拟紧,
当且仅当积 $G_s\times G_s$ 拟紧,
这显然等价于 $G$ 有限.
更一般地, 可知 $B_G$ 的对象 $E$ 拟分离,
当且仅当其各点的稳定子是 $G$ 的\emph{有限}子群,
\ie 它同构于齐性空间 $G/H$ 的一个余积,
其中 $H$ \emph{有限}.
事实上, 由 \hyperref[VI.cor1.15]{1.15},
在此判据中可以假定 $E$ 连通;
但若 $x\in E$ 的稳定子群为 $G_x$,
已知 (IV \hyperref[IV.subsec5.8]{5.8})\pageoriginale{}
$B_G/E$ 等价于拓扑斯 $B_{G_x}$,
故其终对象拟分离, 当且仅当 $G_x$ 有限.
这一判据表明,
$B_G$ 中位于一个拟分离对象之上的每个对象都拟分离;
更有甚者, $B_G$ 中两个拟分离对象的积拟分离.
它还表明, $B_G$ 的凝聚对象,
是 $B_G$ 中那些同构于齐性空间 $G/H$ 的有限余积的对象,
其中 $H$ \emph{有限}.
由于 $G_s$ 凝聚, $B_G$ 的凝聚对象已经构成一个生成子范畴.
此外, 此子范畴对纤维积稳定
(这来自位于一个拟分离对象之上的每个对象都拟分离这一事实).
(用下文 \hyperref[VI.def2.11]{2.11} 的术语,
$B_G$ 是一个\emph{局部 Noether} 拓扑斯,
但它拟分离仅当 $G$ 有限,
此时 $B_G$ 甚至为 Noether.)

立即可知, 态射 $f:E'\to E$ 拟紧,
当且仅当轨道集上诱导的映射
$\pi_0(f):\pi_0(E')\to\pi_0(E)$ 为固有的,
\ie 其纤维有限.
我们重新得出单态射拟紧这一事实
(\hyperref[VI.subsec1.32]{1.32}),
从而每个态射拟分离.
因此, $B_G$ 中的凝聚态射就是拟紧态射,
我们刚刚刻画了它们.

设 $E$ 是拓扑斯 $B_G$ 的一个对象.
容易验证, $E$ 是 $B_G$ 的凝聚对象的滤过归纳极限,
当且仅当 $E$ 的各点的稳定子为 ind-有限群
(\ie 其有限子群的滤过\pageoriginale{}归纳极限).
特别地, 若 $G$ 不是 ind-有限的 (例如 $G=\ZZ$),
则拓扑斯 $B_G$ 的终对象不是凝聚对象的归纳极限.
更确切地说, 拓扑斯 $B_G$ 满足 \hyperref[VI.cor1.25]{1.25} 的等价条件,
当且仅当其终对象是凝聚对象的滤过归纳极限,
或者等价地, 当且仅当群 $G$ 为 ind-有限.
\end{enonce*}


\section{拓扑斯的有限性条件}\etiquette[2]{VI.sec2}

\begin{enonce*}{命题 2.1}\etiquette[2.1]{VI.prop2.1}
设 $E$ 是一个 $𝒰$-拓扑斯.
下列条件等价:
\begin{enumerate}
\renewcommand{\theenumi}{\roman{enumi}}
\renewcommand{\labelenumi}{(\theenumi)}
\item $E$ 中存在一个由拟紧对象构成且对纤维积稳定的满生成子范畴 $C$.

\item 存在一个 $𝒰$-景 $C$,
使 $C$ 的每个对象拟紧、$C$ 中纤维积可表示,
且 $E$ 等价于层范畴 $C^{\sim}$.
\end{enumerate}

假定这些条件成立.
则在 \emph{(i) (\resp (ii))} 中甚至可以选取范畴 $C$ 为 $𝒰$-小的.
此外, 对每个 $X\in\ob C$,
拓扑斯 $E$ 的对象 $X$ \emph{(\resp $\epsilon(X)$)} \emph{凝聚}.
\end{enonce*}

蕴含 (i) $\Rightarrow$ (ii) 来自 IV \hyperref[IV.cor1.2.1]{1.2.1};
而在 (i) (\resp (ii)) 中可以取 $C$ 为 $𝒰$-小的这一事实,
立即来自将 IV \hyperref[IV.subsubsec1.2.3]{1.2.3} 的论证
应用于 $C$ 的一个拓扑意义下
(II \hyperref[II.def3.0.1]{3.0.1}) 的小满生成子范畴.
只须证明 (ii) $\Rightarrow$ (i)
及 \hyperref[VI.prop2.1]{2.1} 的最后一个断言;
它们立即来自下面的结果.

\begin{enonce*}{推论 2.1.1}\etiquette[2.1.1]{VI.cor2.1.1}
设\pageoriginale{} $C$ 是一个其中纤维积可表示且每个对象拟紧的 $𝒰$-景,
$E=C^{\sim}$, $\varepsilon:C\to E$ 为典范函子.
则对每个 $X\in\ob C$, $\varepsilon(X)$ 是 $E$ 的凝聚对象.
$E$ 中由位于某个 $\varepsilon(X)$ 之上的 $E$ 的凝聚对象构成的满子范畴
对纤维积稳定 (并且显然在 $E$ 中生成).
\end{enonce*}

已经知道 $\epsilon(X)$ 拟紧 (\hyperref[VI.prop1.2]{1.2});
下面证明它拟分离.
设 $F,G$ 是 $E$ 的两个拟紧对象,
$F\to\varepsilon(X)$, $G\to\epsilon(X)$ 是两个态射.
须证明纤维积 $\fprod{F}{G}{\epsilon(X)}$ 拟紧.
用有限多个形如 $\epsilon(Y_i)$ 及分别形如 $\epsilon(Z_j)$ 的对象
覆盖 $F$ 及 $G$,
并仿照 \hyperref[VI.prop1.2]{1.2},
归结为复合 $\epsilon(Y_i)\to\epsilon(X)$
及 $\epsilon(Z_j)\to\epsilon(X)$
分别来自态射 $Y_i\to X$ 及 $Z_j\to X$ 的情形,
由 \hyperref[VI.prop1.3]{1.3} 可归结为
态射 $F\to\epsilon(X)$ 和 $G\to\epsilon(X)$
由 $\epsilon$ 作用于态射 $Y\to X$ 和 $Z\to X$ 得到的情形.
但因 $C$ 中纤维积可表示, 且 $\epsilon$ 与纤维积交换,
归结为证明 $\epsilon(\fprod{Y}{Z}{X})$ 拟紧;
对每个形如 $\epsilon(U)$ 的对象, 这已在上面指出.

再设 $G\to F$, $H\to F$ 是 $E$ 的凝聚对象之间的态射,
并假定对某个 $X\in\ob C$ 存在态射 $F\to\epsilon(X)$.
证明 $\fprod{G}{H}{F}$ 凝聚.
由于 $G$ 与 $H$ 拟紧而 $F$ 拟分离,
由定义已知此纤维积拟紧.
还须\pageoriginale{}证明它拟分离,
或等价地 (\hyperref[VI.cor1.15]{1.15}),
证明 $\fprod{G}{H}{\epsilon(X)}$ 拟分离.
这来自更一般的事实:

\begin{enonce*}{推论 2.1.2}\etiquette[2.1.2]{VI.cor2.1.2}
在 \hyperref[VI.cor2.1.1]{2.1.1} 的条件下,
对 $E$ 的每个拟分离对象 $G$,
每个态射 $G\to\epsilon(X)$ 都拟分离.
若 $H\to\epsilon(X)$ 为第二个态射, $H$ 拟分离,
则 $\fprod{G}{H}{\epsilon(X)}$ 拟分离.
\end{enonce*}

第一个断言蕴含第二个:
经基变换, 它蕴含 $\fprod{G}{H}{\epsilon(X)}\to H$ 拟分离;
又因 $H$ 拟分离,
$\fprod{G}{H}{\epsilon(X)}$ 也拟分离
(\hyperref[VI.prop1.14]{1.14} (ii)).
要证明 $G\to\epsilon(X)$ 拟分离, 将用下面的引理.

\begin{enonce*}{引理 2.1.3}\etiquette[2.1.3]{VI.lem2.1.3}
设 $E$ 是一个拓扑斯, $f:X\to Y$ 是 $E$ 中的态射.
假定 $E$ 有由拟紧对象构成的生成族,
且 $X$ 有一个覆盖族 $(g_i:X_i\to X)_{i\in I}$,
使各 $X_i$ 拟紧, 且对每对指标 $i,j\in I$,
$\fprod{X_i}{X_j}{Y}$ 拟分离.
则若 $X$ 是 $E$ 的拟分离对象, $f$ 为拟分离态射.
\end{enonce*}

事实上, 族
$(\fprod{g_i}{g_j}{Y}:\fprod{X_i}{X_j}{Y}\to
\fprod{X}{X}{Y})_{(i,j)\in I\times I}$ 覆盖.
因此, 要证明态射 $\Delta_{X/Y}:X\to\fprod{X}{X}{Y}$ 拟紧,
只须证明对每个 $(i,j)\in I\times I$,
由它沿 $\fprod{g_i}{g_j}{Y}$ 经基变换\pageoriginale{}所得的态射拟紧
(\hyperref[VI.prop1.10]{1.10} (ii)).
后者正是典范态射
$\fprod{X_i}{X_j}{X}\to\fprod{X_i}{X_j}{Y}$.
由假设 $X$ 拟分离且各 $X_i$ 拟紧,
故 $\fprod{X_i}{X_j}{X}$ 拟紧;
而由假设 $\fprod{X_i}{X_j}{Y}$ 拟分离,
故态射 $\fprod{X_i}{X_j}{X}\to\fprod{X_i}{X_j}{Y}$ 确实拟紧, \cqfd

现在可完成 \hyperref[VI.cor2.1.2]{2.1.2} 的证明,
即证明每个 $G$ 拟分离的态射 $f:G\to\epsilon(X)$ 拟分离.
为此, 用对象 $\epsilon(X_i)$ 覆盖 $G$,
使复合态射 $\epsilon(X_i)\to\epsilon(X)$ 来自态射 $X_i\to X$.
由 \hyperref[VI.lem2.1.3]{2.1.3},
因各 $\epsilon(X_i)$ 拟紧,
只须验证纤维积
$\fprod{\epsilon(X_i)}{\epsilon(X_j)}{\epsilon(X)}=
\epsilon(\fprod{X_i}{X_j}{X})$ 拟分离.
我们已知 $E$ 中每个形如 $\epsilon(U)$ 的对象都凝聚.
这就证明 \hyperref[VI.cor2.1.2]{2.1.2},
并完成 \hyperref[VI.prop2.1]{2.1} 的证明.

\begin{enonce*}{命题 2.2}\etiquette[2.2]{VI.prop2.2}
设 $E$ 是一个含有由凝聚对象构成的满生成子范畴 $C$ 的拓扑斯.
下列关于 $E$ 的条件等价:
\begin{itemize}
\item[(i)] $E$ 的每个拟分离对象都在终对象上拟分离.

\item[(i bis)] 对 $E$ 中的每个态射 $f:X\to Y$,
若 $X$ 拟分离, 则 $f$ 拟分离.

\item[(i ter)] $C$ 的每个对象 $X$ 都在 $E$ 的终对象上拟分离,
\ie 对角态射 $X\to X\times X$ 拟紧.

\item[(ii)] $E$ 中两个拟分离对象的积\pageoriginale{}拟分离.

\item[(ii bis)] $C$ 的两个对象在 $E$ 中的积拟分离.
\end{itemize}

这些条件蕴含 $E$ 中由 $E$ 的凝聚对象构成的满子范畴 $E_{\coh}$
对纤维积稳定
(从而 $E$ 更满足 \hyperref[VI.prop2.1]{2.1} 的等价条件).
\end{enonce*}

显然 (i bis) $\Rightarrow$ (i),
逆向蕴含来自 \hyperref[VI.prop1.8]{1.8} (iii).
用已使用过的论证可得 (i) $\Rightarrow$ (ii):
若 $X$ 与 $Y$ 拟分离, 则由假设 (i), $X$ 在 $e$ 上拟分离,
故经基变换 $X\times Y$ 在 $Y$ 上拟分离,
从而由 \hyperref[VI.prop1.14]{1.14} (ii), $X\times Y$ 拟分离.
同一论证表明 (i ter) $\Rightarrow$ (ii bis),
而 (i) $\Rightarrow$ (i ter) 及 (ii) $\Rightarrow$ (ii bis) 显然.
因此, 只须证明 (ii bis) $\Rightarrow$ (i).
若 $X$ 是 $E$ 的拟分离对象,
用 $C$ 的对象 $X_i$ 覆盖 $X$
(由于 $C$ 生成, 这可以做到).
由于对 $C$ 的假设, 各 $X_i$ 拟紧,
可以将 \hyperref[VI.lem2.1.3]{2.1.3} 用于态射 $X\to e$;
因此, 要证明此态射拟分离,
恰好归结为验证各 $X_i\times X_j$ 拟分离,
这正是假设 (ii bis).

\begin{enonce*}[remark]{注 2.2.1}\etiquette[2.2.1]{VI.rem2.2.1}
条件 (i ter) 也等价于下列条件 (i quater):
对 $C$ 中的每个双箭头 $g_1,g_2:Y\rightrightarrows X$,
其在 $E$ 中的核拟紧
(或等价地 (\hyperref[VI.cor1.15]{1.15}), 凝聚).
\end{enonce*}

为\pageoriginale{}证明这一点,
只须将判据 \hyperref[VI.prop1.10]{1.10} (i)
用于 (i ter) 中考察的对角态射 $X\to X\times X$.

\begin{enonce*}{定义 2.3}\etiquette[2.3]{VI.def2.3}
满足 \emph{\hyperref[VI.prop2.2]{2.2}
(\resp \hyperref[VI.prop2.1]{2.1})} 的等价条件的拓扑斯 $E$,
称为\emph{代数拓扑斯（algebraic topos）
(\resp 局部代数拓扑斯, 或局部凝聚拓扑斯)}.
拓扑斯 $E$ 的对象 $X$ 称为\emph{拓扑斯 $E$ 的代数对象},
是指诱导拓扑斯 $E_{/X}$ 为代数的.
拓扑斯 $E$ 称为\emph{拟分离的 (\resp 凝聚的)},
是指它为代数拓扑斯,
且其终对象拟分离 (\resp 凝聚).
\end{enonce*}

\setcounter{subsection}{3}
\subsubsection{}\etiquette[2.3.1]{VI.subsubsec2.3.1}
在代数拓扑斯中,
第 ${\rm n^{\circ}}$~1 节发展的有限性概念具有完全令人满意的性质,
与概形范畴中同名概念的性质类似.
实际遇到的所有局部代数 (= 局部凝聚) 拓扑斯事实上都是代数拓扑斯,
因此这一概念当前的实际意义似乎颇为有限.
不过, 可以构造局部代数但不是代数的拓扑斯
(\hyperref[VI.exer2.17]{2.17} f).


% END INCLUDED SGA4_English_current_buildable_source_20260906/source/06/drafts/SGA4_VI_LINES_1096_1382_English.texfrag


% BEGIN INCLUDED SGA4_English_current_buildable_source_20260906/source/06/drafts/SGA4_VI_LINES_1383_1678_English.texfrag
\subsection{}\etiquette[2.4]{VI.subsec2.4}
下面对 \hyperref[VI.def2.3]{2.3} 引入的概念作若干一般说明,
希望读者由此相信, 那里引入的术语是一致的.

\subsubsection{}\etiquette[2.4.1]{VI.subsubsec2.4.1}
在 \hyperref[VI.prop2.1]{2.1} (i) \resp (ii) 的条件下,
对每个 $X\in C$, $X$ (\resp $\epsilon(X)$)
是 $E$ 的一个代数 (凝聚) 对象
(\hyperref[VI.cor2.1.2]{2.1.2} 及 1.19.2);
换句话\pageoriginale{}说,
诱导拓扑斯 $E_{/X}$
(\resp $E_{\epsilon(X)}\approxeq(C_{/X})^{\sim}$)
为代数的, 从而\emph{凝聚}.
等价地, \emph{当 $C$ 具有终对象时, $E$ 本身为代数的, 从而凝聚.}

\subsubsection{}\etiquette[2.4.2]{VI.subsubsec2.4.2}
若 $E$ 是代数拓扑斯,
则每个诱导拓扑斯 $E_{/X}$ 也为代数的
(\hyperref[VI.prop2.2]{2.2} (i bis)).
拓扑斯 $E$ 局部代数 $=$ 局部凝聚的充要条件是:
存在 $E$ 的一个覆盖终对象的对象族 $(X_i)_{i\in I}$,
使各 $X_i$ 为代数的 (\resp 凝聚的),
\ie 各诱导拓扑斯 $E_{/X_i}$ 为代数的 (\resp 凝聚的).
充分性由定义立即得到,
必要性来自 \hyperref[VI.subsubsec2.4.1]{2.4.1} 的观察.
当然, 代数拓扑斯局部代数;
若 $E$ 为局部代数拓扑斯, 每个诱导拓扑斯 $E_{/X}$ 都局部代数.

\subsubsection{}\etiquette[2.4.3]{VI.subsubsec2.4.3}
若 $X$ 是拓扑斯 $E$ 的代数对象,
则 $E$ 中位于 $X$ 之上的每个对象 $X'$ 仍为代数的
(\cf \hyperref[VI.subsubsec2.4.2]{2.4.2}).
特别地, 代数拓扑斯的每个对象都为代数的; 逆命题当然也成立.

\subsubsection{}\etiquette[2.4.4]{VI.subsubsec2.4.4}
$E$ 中由\emph{代数的拟分离}对象构成的满子范畴,
对纤维积及子对象稳定 (\hyperref[VI.cor1.15]{1.15}),
因而也对核稳定.
若 $E$ 为代数的,
此范畴 (这时它就是\pageoriginale{} $E$ 的拟分离对象范畴)
还对两个对象的积稳定.
若 $E$ 拟分离 (且仅在此时),
这一范畴对任意有限投射极限稳定,
或者等价地, 含有 $E$ 的终对象.

$E$ 中由\emph{代数的凝聚}对象构成的满子范畴 $C$ 对纤维积稳定.
若 $E$ 局部代数, \ie 局部凝聚,
称 $E$ 为代数的, 相当于称 $C$ 还对双箭头的核稳定
(\hyperref[VI.rem2.2.1]{2.2.1});
称 $E$ 拟分离, 相当于再要求 $C$ 对两个对象的积稳定
(\cf \hyperref[VI.cor1.17]{1.17}).
最后, 称 $E$ 凝聚,
相当于称 $C$ 对任意有限投射极限稳定,
或者等价地, 含有 $E$ 的终对象.

\subsubsection{}\etiquette[2.4.5]{VI.subsubsec2.4.5}
设 $E$ 是一个拓扑斯.
$E$ 局部凝聚 (\resp 代数, \resp 拟分离, \resp 凝聚)
的充要条件是, 它具有一个由拟紧对象构成的满生成子范畴 $C$
(这些对象自动凝聚 (\hyperref[VI.prop2.1]{2.1})),
且\emph{对}纤维积
(\resp 对纤维积与核,
\resp 对纤维积及两个对象的积,
\resp 对任意有限投射极限) \emph{稳定};
或等价地, $E$ 等价于一个拓扑斯 $C^{\sim}$,
其中 $C$ 为每个对象都拟紧的 $𝒰$-景,
且其中纤维积及两个对象的积
(\resp 所有有限投射极限) 可表示.
(证明必要性, 取\pageoriginale{} $C$ 为 $E$ 的凝聚且\emph{代数的}对象范畴,
并应用 \hyperref[VI.subsec2.4]{2.4};
证明充分性, 应用 \hyperref[VI.prop2.2]{2.2} (i ter).)
在这一陈述中, 当然还可假定 $C$ 为 $𝒰$-小的.

\subsubsection{}\etiquette[2.4.6]{VI.subsubsec2.4.6}
设 $E$ 是一个拓扑斯.
若 $E$ 为代数的, $E$ 的对象 $X$ 拟分离 (\resp 凝聚),
当且仅当诱导拓扑斯 $E_{/X}$ 也如此.
当不再假定 $E$ 为代数的时,
只能说 $E_{/X}$ 拟分离 (\resp 凝聚),
当且仅当 $X$ 拟分离 (\resp 凝聚) 且\emph{为代数的}.
这无关大局, 因为我们大概无须在 $E$ 为代数的情形之外使用这些概念.

\subsubsection{}\etiquette[2.4.7]{VI.subsubsec2.4.7}
为了使拓扑斯 $E$ 的凝聚 (\resp 拟分离) 对象概念具有令人满意的性质,
须限于同时为代数的对象;
若 $E$ 为代数的, 这一条件自动满足.
由于我们大概始终无须处理非代数的局部代数拓扑斯 $E$,
修改 \hyperref[VI.def1.13]{1.13} 引入的术语,
可能将它仅保留给代数对象的问题, 因而并不十分迫切!

此外, 我们没有定义拟紧拓扑斯的概念.
按本段的思路, 自然会把这一名称给予终对象拟紧的\emph{代数}拓扑斯.

我们\pageoriginale{}对引入这样的用法有所犹豫,
因为若 $X$ 是拓扑空间,
则 $X$ 拟紧当且仅当相伴拓扑斯 $\Top(X)$ 拟紧这一说法就不成立.
关于这一点, 注意 $\Top(X)$ 局部凝聚,
当且仅当 $X$ 具有一个由拟紧开集 $U$ 构成的开集基,
使对 $U_j,U_k\subset U_i$, $U_j\cap U_k$ 仍拟紧
(这时 $\Top(X)$ 甚至为代数的).
概形的底拓扑空间满足这一条件,
但分析学家遇到的空间很少满足它, 即使空间本身紧致也如此.
为了便于查阅, 明确陈述如下:

\begin{enonce*}{命题 2.5}\etiquette[2.5]{VI.prop2.5}
设 $E$ 是一个代数拓扑斯, $f:X\to Y$ 是 $E$ 中的态射.
若 $Y$ 是 $E$ 的拟分离 (\resp 凝聚) 对象,
则 $X$ 是 $E$ 的拟分离 (\resp 凝聚) 对象的充要条件是,
态射 $f$ 拟分离 (\resp 凝聚).
\end{enonce*}

充分性已在 \hyperref[VI.prop1.14]{1.14} (ii) 中见到,
且不需对 $E$ 作任何假设.
非括号内情形的必要性, 甚至无须假定 $Y$ 拟分离就成立,
它通过 \hyperref[VI.prop2.2]{2.2} i bis)
就是定义 \hyperref[VI.def2.3]{2.3}.
只须证明若 $Y$ 凝聚且 $X$ 凝聚, 则 $f$ 凝聚.
既已知道 $f$ 拟分离, 只须证明 $f$ 拟紧;
这来自 $X$ 拟紧而 $Y$ 拟分离
(\hyperref[VI.def1.13]{1.13}).

\begin{enonce*}{推论 2.6}\etiquette[2.6]{VI.cor2.6}
设\pageoriginale{} $E$ 是一个代数拓扑斯,
$f:X\to Y$ 是 $E$ 中的态射,
$(Y_i\to Y)_{i\in I}$ 为覆盖族,
且各 $Y_i$ 拟分离 (\resp 凝聚).
$f$ 拟分离 (\resp 拟紧, \resp 凝聚) 的充要条件是,
对每个 $i\in I$,
$X_i=\fprod{X}{Y_i}{Y}$ 是 $E$ 的拟分离
(\resp 拟紧, \resp 凝聚) 对象.
\end{enonce*}

在“拟分离”或“凝聚”情形,
只须合并 \hyperref[VI.prop1.10]{1.10} (ii)
及 \hyperref[VI.prop2.5]{2.5};
“拟紧”情形已在 \hyperref[VI.cor1.16]{1.16} 中见到.
注意, \hyperref[VI.cor2.6]{2.6} 中所述条件的必要性,
也无须对拓扑斯 $E$ 作任何假设就成立.

作为 \hyperref[VI.cor2.6]{2.6} 的特殊情形, 得到:

\begin{enonce*}{推论 2.7}\etiquette[2.7]{VI.cor2.7}
设 $E$ 是一个凝聚拓扑斯, $X$ 是 $E$ 的一个对象.
$X$ 是 $E$ 的拟紧 (\resp 拟分离, \resp 凝聚) 对象的充要条件是,
它在终对象上拟紧 (\resp 拟分离, \resp 凝聚)
(\hyperref[VI.subsubsec1.9.2]{1.9.2}).
\end{enonce*}

特别地, $X$ 可构造 (\hyperref[VI.subsubsec1.9.2]{1.9.2}),
当且仅当它凝聚.

\begin{enonce*}{推论 2.8}\etiquette[2.8]{VI.cor2.8}
设 $E$ 是一个局部凝聚拓扑斯,
$f:X\to Y$ 是 $E$ 中的态射.
$f$ 拟分离的充要条件是,
对 $E_{/Y}$ 的每个拟分离对象 $Y'$,
$f^{\ast}(Y')=\fprod{X}{Y'}{Y}$
是 $E_{/X}$ 的拟分离对象
(或者是 $E$ 的拟分离对象, 两者相同
(\hyperref[VI.subsubsec1.20.2]{1.20.2})).
\end{enonce*}

“仅当”\pageoriginale{}方向无须对 $E$ 作任何条件就成立,
因为若 $Y'$ 拟分离,
则经基变换 $f':X'=\fprod{X}{Y'}{Y}\to Y'$ 拟分离,
从而 $X'$ 拟分离 (\hyperref[VI.prop1.14]{1.14} (ii)).
证明逆向, 注意存在一个由拟分离代数对象构成的 $Y$ 的覆盖族 $Y_i\to Y$.
由 \hyperref[VI.prop1.10]{1.10} (ii),
要验证 $f$ 拟分离, 只须验证各
$f_i:X_i=\fprod{X}{X_i}{Y}\to Y_i$ 拟分离.
但对 $f^{\ast}$ 的假设蕴含各 $X_i$ 拟分离,
因而各 $f_i$ 也拟分离,
因为 $E_{/Y_i}$ 为代数拓扑斯.

\begin{enonce*}[remark]{注 2.8.1}\etiquette[2.8.1]{VI.rem2.8.1}
将假设“$E$ 为代数的”换为更一般的假设“$Y$ 为代数的”时,
陈述 \hyperref[VI.prop2.5]{2.5} 及 \hyperref[VI.cor2.6]{2.6} 仍成立;
只须将原陈述应用于为代数的诱导拓扑斯 $E_{/Y}$, 立即可见.
同样, 在 \hyperref[VI.cor2.8]{2.8} 中,
只须假定 $E_{/Y}$ (而非 $E$) 局部凝聚.
\end{enonce*}

\begin{enonce*}{命题 2.9}\etiquette[2.9]{VI.prop2.9}
设 $E$ 是一个拓扑斯, $C$ 是 $E$ 的一个子范畴.
下列条件等价:
\begin{enumerate}
\renewcommand{\theenumi}{\roman{enumi}}
\renewcommand{\labelenumi}{(\theenumi)}
\item $E$ 局部凝聚 (\resp 凝聚),
满足 \hyperref[VI.cor1.25]{1.25} 的等价条件,
且 $C$ 是 $E$ 中由凝聚对象构成的满子范畴 $E_{\coh}$.

\item $C$ 是 $E$ 的一个严格满生成子范畴,
对有限归纳极限和纤维积
(\resp 以及有限投射极限) 稳定,
且由拟紧对象构成.
\end{enumerate}
\end{enonce*}

蕴含\pageoriginale{} (i) $\Rightarrow$ (ii)
立即来自定义以及 \hyperref[VI.cor1.25]{1.25} 中所考察条件的表述
\hyperref[VI.cor1.25]{1.25} (ii).
证明逆向蕴含.
$C$ 是一个由拟紧对象构成且对纤维积
(\resp 有限投射极限) 稳定的满生成子范畴,
由定义就蕴含 $E$ 局部凝聚 (\resp 凝聚).
再由 $C$ 严格满且对有限归纳极限稳定,
可知 $C=E_{\coh}$
(由 \hyperref[VI.cor1.24]{1.24} (ii) $\Rightarrow$ (iii)),
且条件 \hyperref[VI.cor1.25]{1.25} (ii) 满足, \cqfd

\begin{enonce*}{定义 2.9.1}\etiquette[2.9.1]{VI.def2.9.1}
$𝒰$-拓扑斯 $E$ 称为完美的,
是指它凝聚 (\hyperref[VI.def2.3]{2.3})
且满足 \hyperref[VI.cor1.25]{1.25} 的等价条件,
\ie $E$ 中由 $E$ 的凝聚对象构成的子范畴 $E_{\coh}$ 对有限归纳极限稳定.
若 $E$ 存在一个覆盖终对象的对象族 $(X_i)_{i\in I}$,
使对每个 $i\in I$, 诱导拓扑斯 $E_{/X_i}$ 完美,
就称 $E$ 局部完美.
\end{enonce*}

\setcounter{subsection}{9}
\setcounter{subsubsection}{1}
\subsubsection{}\etiquette[2.9.2]{VI.subsubsec2.9.2}

由此定义立即可知,
完美 (\resp 局部完美) 拓扑斯凝聚 (\resp 局部凝聚).
完美 (\resp 局部完美) 拓扑斯的例子,
可以列举下文 \hyperref[VI.def2.11]{2.11} 引入的
Noether (\resp 局部 Noether) 拓扑斯
(\cf \hyperref[VI.subsec2.14]{2.14}),
以及凝聚概形 (\resp 任意概形) 的 Zariski 与平展拓扑斯
(IX, 注, 第 42 页).
反例\pageoriginale{}则可列举概形 $S$ 的 fppf 拓扑斯,
它局部凝聚
(而且若 $S$ 为凝聚概形, \ie 拟紧且拟分离, 则甚至凝聚)
(\hyperrefext[VII][VII.lem5.6]{5.6}),
但若 $S\neq∅$, 它不局部完美 (\hyperref[VI.exer1.28]{1.28} f)).

在局部完美拓扑斯中, 可构造对象的概念特别稳定
(下文 \hyperref[VI.prop2.9.3]{2.9.3}),
这是此概念似乎有用的一个理由.
另一个理由是, 与凝聚或局部凝聚拓扑斯概念一样,
完美或局部完美拓扑斯概念在形成某个开集的补闭子拓扑斯时稳定
(\hyperref[VI.prop4.6]{4.6}),
且在拓扑斯粘合运算下的行为特别简单
(\hyperref[VI.cor4.11]{4.11}, \hyperref[VI.cor4.14]{4.14}).

\begin{enonce*}{命题 2.9.3}\etiquette[2.9.3]{VI.prop2.9.3}
设 $E$ 是一个局部完美拓扑斯 (\hyperref[VI.def2.9.1]{2.9.1}).
则 $E$ 中由 $E$ 的可构造对象
(\hyperref[VI.subsubsec1.9.3]{1.9.3}) 构成的满子范畴 $E_{\cons}$
对有限归纳极限稳定
(当然, 由 (\hyperref[VI.subsubsec1.9.3]{1.9.3}), 也对有限投射极限稳定).
\end{enonce*}

由于 $E$ 的对象可构造的性质在 $E$ 上局部
(\hyperref[VI.prop1.10]{1.10} (ii)),
归结为 $E$ 是完美拓扑斯的情形.
但此时 $E_{\cons}=E_{\coh}$ (\hyperref[VI.cor2.7]{2.7}),
结论由定义 \hyperref[VI.def2.9.1]{2.9.1} 得到.

\begin{enonce*}{推论 2.9.4}\etiquette[2.9.4]{VI.cor2.9.4}
设\pageoriginale{} $E$ 是一个局部完美拓扑斯.
则 $E$ 完美, 当且仅当 $E$ 凝聚.
对 $E$ 的每个对象 $X$, 诱导拓扑斯 $E_{/X}$ 局部完美.
\end{enonce*}

\begin{enonce*}[remark]{问题 2.9.5}\etiquette[2.9.5]{VI.prob2.9.5}
可以设想, 完美拓扑斯足够特殊,
容许同样明确的\emph{结构理论}
(沿着 M.\textsc{Artin} 在口头讨论班期间提出的模型).
若一个拓扑斯等价于形如 $\widehat{C}$ 的拓扑斯,
其中 $C$ 是有限范畴 (\cf 习题 \hyperref[VI.exer3.11]{3.11}),
就称其为\emph{有限的};
若一个拓扑斯是有限拓扑斯的滤过投射极限,
就称其为\emph{pro-有限的}
(取下文第 6 节的意义, 由 \hyperref[VI.exer3.12]{3.12} c) 可应用).
由于有限拓扑斯为 Noether (\hyperref[VI.exer2.17]{2.17} g)),
因而完美, 而完美拓扑斯的滤过投射极限完美 (6),
故 pro-有限拓扑斯完美.
问题是, 反过来每个完美拓扑斯是否 pro-有限.
这等价于问, $E_{\coh}$ 的每个有限子范畴
是否都包含在 $E_{\coh}$ 的一个满子范畴 $F_0$ 中,
使其等价于某个范畴 $F_{\coh}$,
其中 $F$ 为有限拓扑斯 (\cf \hyperref[VI.exer3.11]{3.11}).
(若答案是否定的,
就须找到对完美拓扑斯的便于处理的内在附加条件,
保证它为 pro-有限.)
最后, 还须用适当的“Karoubi pro-有限范畴”概念
研究 pro-有限拓扑斯的结构,
这一概念受 IV \hyperref[IV.exer7.6]{7.6} h) 启发.
即使在 $E$ 是某个良好的凝聚概形 $X$ 的 Zariski 或平展拓扑斯的特殊情形,
这一研究\pageoriginale{}也尚未开展;
它应当给出比与 $X$ 相伴的紧空间 $X_{\cons}$
(\EGA IV 1.) 更精细的不变量.

同样地, 可以提出下列问题:
设 $E$ 是完美拓扑斯,
$E'$ 是由 $E$ 的一个开集诱导的拓扑斯,
$E''$ 是相应的闭子拓扑斯
(IV \hyperref[IV.subsubsec9.3.5]{9.9}).
若 $E'$ 与 $E''$ 为有限拓扑斯, $E$ 是否也如此?
\end{enonce*}

% END INCLUDED SGA4_English_current_buildable_source_20260906/source/06/drafts/SGA4_VI_LINES_1383_1678_English.texfrag


% BEGIN INCLUDED SGA4_English_current_buildable_source_20260906/source/06/drafts/SGA4_VI_LINES_1679_2073_English.texfrag
\begin{enonce*}{命题 2.10}\etiquette[2.10]{VI.prop2.10}
设 $E$ 是一个拓扑斯. 下列条件等价:
\begin{itemize}
\item[(i)] $E$ 具有一个对纤维积稳定、
由 $E$ 的预 Noether 对象构成的满生成子范畴.

\item[(ii)] $E$ 具有一个由拟分离预 Noether 对象
(\ie 凝聚预 Noether 对象) 构成的满生成子范畴.

\item[(iii)] $E$ 局部凝聚 (\hyperref[VI.def2.3]{2.3}),
且 $E$ 的每个拟紧对象都为预 Noether.

\item[(iii bis)] $E$ 局部凝聚,
且具有由 $E$ 的预 Noether 对象构成的生成族.
\end{itemize}
\end{enonce*}

(iii) 与 (iii bis) 的等价性来自 \hyperref[VI.subsec1.32]{1.32},
而 (i) $\Rightarrow$ (iii bis) 及 (iii) $\Rightarrow$ (i) 显然;
故 (i) 等价于 (iii) 及 (iii bis).
最后, (i) $\Rightarrow$ (ii)
来自 \hyperref[VI.prop2.1]{2.1} 的最后一个断言,\pageoriginale{}
只须证明 (ii) $\Rightarrow$ (i).
这一蕴含来自以下事实:
$E$ 中由拟分离 Noether 对象构成的满子范畴若生成, 则对纤维积稳定:
这样的纤维积由定义拟紧,
故由 \hyperref[VI.subsec1.32]{1.32} 的第一个断言为 Noether;
而由 \hyperref[VI.subsec1.32]{1.32} 的最后一个断言, 它拟分离.

\begin{enonce*}{定义 2.11}\etiquette[2.11]{VI.def2.11}
拓扑斯 $E$ 称为局部 Noether 的,
是指它满足 \hyperref[VI.prop2.10]{2.10} 的等价条件;
若其终对象还凝聚
(或等价地 (\hyperref[VI.subsec1.32]{1.32}), 为预 Noether 且拟分离),
就称它为 Noether 的.
拓扑斯 $E$ 的对象 $X$ 称为 $E$ 的 Noether 对象,
是指诱导拓扑斯 $E_{/X}$ 为 Noether 拓扑斯.
\end{enonce*}

\setcounter{subsection}{11}
\subsection{}\etiquette[2.12]{VI.subsec2.12}
若 $E$ 是局部 Noether 拓扑斯,
则每个诱导拓扑斯 $E_{/X}$ 也如此.
反过来, 若能找到覆盖终对象的一个族 $(X_i)$,
使各 $E_{/X_i}$ 局部 Noether, 则 $E$ 也如此.
若 $C$ 是 $E$ 的一个由预 Noether 对象构成、对纤维积稳定的满生成子范畴
(\hyperref[VI.def2.11]{2.11} (i)),
则对每个 $X \in C$, $E_{/X}$ 为 Noether 拓扑斯,
\ie $C$ 的对象甚至是 \emph{Noether} 的
(因为若 $E$ 局部 Noether 且 $X \in \Ob E$,
则 $E_{/X}$ 为 Noether 的, 当且仅当 $X$ 凝聚).
由此可知, 拓扑斯 $E$ 局部 Noether,
当且仅当 $E$ 的终对象可以由 Noether 对象 $X_i$ 覆盖.

\subsection{}\etiquette[2.13]{VI.subsec2.13}
在局部 Noether 拓扑斯 $E$ 中,\pageoriginale{}
每个单态射拟紧, 每个态射拟分离
(\hyperref[VI.subsec1.32]{1.32}).
更有甚者, $E$ 为\emph{代数}拓扑斯
(\hyperref[VI.def2.3]{2.3}),
而不只是局部凝聚, \ie 局部代数.
因此, 拓扑斯 $E$ 为 Noether 的,
当且仅当它局部 Noether 且\emph{凝聚}
(\hyperref[VI.def2.3]{2.3}).
换句话说, 若 $E$ 是局部 Noether 拓扑斯,
则 $E$ 的 Noether 对象恰为 $E$ 的\emph{凝聚}对象.

\subsection{}\etiquette[2.14]{VI.subsec2.14}
设 $E$ 是一个局部 Noether 拓扑斯.
若 $E$ 拟分离, \ie 其终对象拟分离,
则 $E$ 的\emph{每个}对象都拟分离
(\hyperref[VI.subsec1.32]{1.32});
故 $E$ 的预 Noether 对象与 $E$ 的 Noether
(\ie 凝聚) 对象相同.
由于预 Noether 对象的商对象仍为预 Noether,
可知 $E$ 满足 \hyperref[VI.cor1.25]{1.25} 的等价条件
(因为它满足条件 \hyperref[VI.cor1.25]{1.25} (ii bis)).
特别地, 若 $C$ 是 $E$ 中由 $E$ 的 Noether (凝聚) 对象构成的满子范畴,
则自然函子
$$
\Ind(C) \to E
$$
为范畴等价.
得到结论: \emph{Noether 拓扑斯完美}
(\hyperref[VI.def2.9.1]{2.9.1}),
因而局部 Noether 拓扑斯局部完美.

若只假定 $E$ 局部 Noether,
$E$ 的每个对象仍是其预 Noether 子对象的归纳极限
(因为在具有由\pageoriginale{}拟紧对象构成的生成子范畴的拓扑斯中,
每个对象都是其拟紧子对象的滤过归纳极限).
但由于 $E$ 的预 Noether 对象不再一定凝聚,
这个陈述此时的用途很有限;
由 (\hyperref[VI.exem1.33]{1.33}) 已知,
\hyperref[VI.cor1.25]{1.25} 的条件不再一定满足.

\begin{enonce*}[remark]{例 2.15}\etiquette[2.15]{VI.exem2.15}
(\emph{拓扑空间}.) 设 $X$ 是一个拓扑空间.
则下列条件等价:
(i) 拓扑空间 $X$ 为 Noether 的
(\ie $X$ 的每个开子集拟紧,
\ie $X$ 的每个递增开子集序列稳定);
(ii) 拓扑斯 $\Top(X)$ 为 Noether 的;
(iii) $\Top(X)$ 的终对象 $X$ 为 Noether 的;
(iv) $\Top(X)$ 的终对象 $X$ 为预 Noether 的.

事实上, 显然有 (ii) $\Leftrightarrow$ (iii) $\Rightarrow$ (iv)
$\Leftrightarrow$ (i).
另一方面, (i) 蕴含 $\Top(X)$ 的对有限投射极限稳定的生成子范畴 Ouv(X)
由预 Noether 对象构成, 从而得到 (ii).

同样可见, 拓扑空间 $X$ 局部 Noether
(\ie 为 Noether 开子集的并),
当且仅当拓扑斯 $\Top(X)$ 局部 Noether.
\end{enonce*}

\setcounter{subsection}{15}
\setcounter{subsubsection}{0}
\subsubsection{}\etiquette[2.15.1]{VI.subsubsec2.15.1}
因此, 定义 \hyperref[VI.def1.30]{1.30}
及 \hyperref[VI.def2.11]{2.11}
与拓扑空间的通用术语一致.
另一方面, 在一般情形,
我们选择赋予“Noether 对象”一词比
\hyperref[VI.def1.30]{1.30} 中更朴素的概念更强的\pageoriginale{}含义,
使与这一概念相联的所有性质
(而不只是定义的语法结构)
都与 \emph{Noether} 拓扑空间及概形所附有的直觉完全一致.

读者可在下面的习题中找到更多例子.

\begin{enonce*}[remark]{习题 2.16}\etiquette[2.16]{VI.exer2.16}
(\emph{具有算子群的空间与分类拓扑斯}.)

设 $X$ 是一个拓扑空间, 离散群 $G$ 在其上作用,
得到 (IV \hyperref[IV.subsec2.3]{2.3}) 拓扑斯 $E = \Top(X, G)$.
此拓扑斯的对象 $X'$ 解释为 $X$ 上配备算子群 $G$ 的展空间;
并注意诱导拓扑斯 $E_{/X'}$ 典范等价于 $\Top(X', G)$.
\begin{enumerate}
\renewcommand{\theenumi}{\alph{enumi}}
\renewcommand{\labelenumi}{\theenumi)}
\item 拓扑斯 $E$ 的终对象拟紧,
当且仅当商拓扑空间 $X/G$ 拟紧
(利用 IV \hyperref[IV.subsubsec8.4.1]{8.4.1}).

\item 设 $P$ 为 $E$ 的对象,
由配备 $G$ 的对角作用的平凡展空间 $X \times G$ 定义.
证明诱导拓扑斯 $E_{/P}$ 等价于拓扑斯 $\Top(X)$
(比较 IV \hyperref[IV.subsubsec5.8.3]{5.8.3}).
推出 $E$ 局部凝聚 (\resp 局部 Noether),
当且仅当 $\Top(X)$ 局部凝聚
(\cf \hyperref[VI.subsubsec2.4.7]{2.4.7})
(\resp 局部 Noether (\cf \hyperref[VI.exem2.15]{2.15})).
证明若 $E$ 局部凝聚, \ie 局部代数, 则它实际上为代数的;
利用由对象 $T_U$ 构成\pageoriginale{}的生成族,
其中 $U$ 为 $X$ 的凝聚开子集,
$T_U = G \times U$ 配备 $G$ 的作用
$g \cdot (u, g') = (u, gg')$,
结构态射 $p_U : T_U \to X$
由 $p_U(u, g) = g\cdot u$ 定义.

\item 假定 $E$ 为代数的,
\ie $\Top(X)$ 为代数的 (\hyperref[VI.subsubsec2.4.7]{2.4.7}).
证明结构态射 $P \to e$ 拟分离.
证明 $E$ 拟分离,
当且仅当 $\Top(X)$ 拟分离
(或等价地, 拓扑空间 $X$ 拟分离,
\ie $X$ 的两个拟紧开子集的交是拟紧开子集),
\emph{并且} $G$ 有限或 $X$ 为空.
(比较 \hyperref[VI.exem1.33]{1.33}.)

\item 证明结构态射 $P \to e$ 拟紧,
当且仅当 $G$ 有限.
证明 $E$ 凝聚 (\resp Noether),
当且仅当 $\Top(X)$ 也如此,
\emph{并且} $G$ 有限或 $X$ 为空.

\item 设 $E$ 是一个拓扑斯, $G$ 是 $E$ 的一个群,
得到分类拓扑斯 $B_G$ (IV \hyperref[IV.subsec2.4]{2.4}).
以前述讨论为指引, 研究 $B_G$ 中的有限性条件.
从 $E$ 的严格 pro-群 $𝒢 = (G_i)_{i \in I}$ 出发提出同样的问题.
特别地, 证明若 $E$ 凝聚 (\resp Noether)
且各 $G_i$ 凝聚, 则 $B_{𝒢}$ 凝聚 (\resp Noether).
\end{enumerate}
\end{enonce*}

\begin{enonce*}[remark]{习题 2.17}\etiquette[2.17]{VI.exer2.17}
(\emph{形如 $\widehat{C}$ 的拓扑斯}.)
设 $C$ 是一个等价于某个 $\in 𝒰$ 的范畴的范畴,
$E = \widehat{C}$ 为 $C$ 上的预层范畴.
借助典范函子 $\varepsilon : C \to \widehat{C}$,
将 $C$ 等同于 $E$ 的一个满子范畴.
\begin{enumerate}
\renewcommand{\theenumi}{\alph {enumi}}
\renewcommand{\labelenumi}{(\theenumi)}
\item $C$ 是\pageoriginale{} $E$ 的一个由拟紧对象构成的生成子范畴.
这些对象为预 Noether
(\cf \hyperref[VI.subsec1.32]{1.32}),
当且仅当对 $C$ 的每个对象 $X$,
$X$ 上的每个筛都由 $C$ 的一个有限态射族 $X_i \to X$ 生成.
$E$ 的终对象拟紧 (\resp 预 Noether),
当且仅当存在一个对象集有限的 $C$ 的终子范畴
(\resp $C$ 上的每个筛都由 $C$ 的一个对象集有限的子范畴生成).

\item 拓扑斯 $E$ 具有由凝聚
(或等价地, 拟分离) 对象构成的生成子范畴,
当且仅当 $C$ 的对象在 $E$ 中凝聚;
或等价地, 对 $C$ 中的两个箭头
$Y \xrightarrow{f} X$, $Z \xrightarrow{g} X$,
$E$ 的对象 $\fprod{Y}{Z}{X}$ 拟紧,
\ie 能找到 $C$ 中有限个交换方形
\[
\xymatrix@=1cm{
 & T_i \ar[ld] \ar[rd]& \\
Y\ar[rd]_{f}&    & Z\ar[ld]^{g}\\
 & X  &,
}
\]
使 $C$ 中的每个其他交换方形
\[
\xymatrix@=1cm{
 & T \ar[ld] \ar[rd]& \\
Y\ar[rd]_{f}&    & Z\ar[ld]^{g}\\
 & X  &,
}
\]
都来自\pageoriginale{}某个态射 $T \to T_i$.
(注意, 若 $(F_i)_{i \in I}$ 为 $E$ 的生成族,
则每个 $X \in \ob C$ 都同构于某个 $F_i$ 的一个直因子.)

\item 拓扑斯 $E$ 局部凝聚,
\ie 由一族凝聚且\emph{代数}的对象生成,
当且仅当 $C$ 的对象 $X$ 为 $E$ 的凝聚代数对象,
\ie 条件 b) 及下列条件成立:
对 $C$ 中的每对\emph{在 $C$ 的某个对象 $X$ 上的}箭头
$f$, $g : Y \rightrightarrows Z$,
这对箭头在 $E$ 中的核为 $E$ 的拟紧对象,
\ie $C$ 中存在有限个交换图
\[
\xymatrix@C=1.5cm{
T_i \ar[r]& Y \ar@<2pt>[r]^{f,g}\ar@<-2pt>[r]&Z
}
\]
使 $C$ 中的每个其他交换图
\[
\xymatrix@C=1.5cm{
T \ar[r]& Y \ar@<2pt>[r]^{f,g}\ar@<-2pt>[r]&Z
}
\]
都来自某个态射 $T \to T_i$.
拓扑斯 $E$ 为代数的,
当且仅当它满足条件 b) 及前述条件,
但其中 $(f, g)$ 为 $C$ 的任意箭头对
(不一定在 $C$ 的某个对象 $X$ 上).
拓扑斯 $E$ 拟分离,
当且仅当它满足 b) 的条件,
且对两个对象 $X$, $Y\in\ob C$,
$E$ 中的积 $X\times Y$ 拟紧.
拓扑斯 $E$ 凝聚,
当且仅当它满足前两个条件,
且 $C$ 有一个底集有限的终子范畴.

\item 给出\pageoriginale{}一个例子,
使 $E$ 的生成子范畴 $C$ 的对象为预 Noether,
但 $E$ 不具有由凝聚对象构成的生成族
(\ie b) 不成立, 或 $C$ 的对象不全凝聚).
为此取 $C$ 为下列范畴:
对象为 $X$, $T_i$ $(i = 0, 1, \ldots)$
及 $e$ (终对象);
除到 $e$ 的结构态射及恒等态射外,
这些对象之间的态射只有满足 $u_i v_i = \id_X$ 的态射
$u_i : T_i \to X$ 及 $v_i : X \to T_i$;
最后置 $p_i = v_i u_i$ (它必然满足 $p_i^2 = p_i$).
验证 $X$ 或任意 $T_i$ 上的筛只有两个平凡筛,
而 $e$ 恰有一个非平凡筛.
因此由 a), $C$ 的对象都是 $E$ 的预 Noether 对象.
不过, $E$ 中的积 $X \times X$ 不拟紧,
因为它不满足 b) 的判据.

\item 给出一个例子,
使 $E$ 的生成子范畴 $C$ 由 $E$ 的凝聚对象构成,
但 $E$ 不局部凝聚.
为此取 $C$ 为下列范畴:
对象集由相异对象 $e$ (终对象)、$Y$、$Z$
及 $T_i$ $(i=0, 1, \ldots)$ 构成;
除到终对象 $e$ 的态射及恒等态射外,
这些对象间的态射只有满足 $fu_i = gu_i$ 的态射
$f$, $g : Y \rightrightarrows Y$.
验证 $C$ 满足条件 b)
(稍加耐心可得, $C$ 的对象的所有纤维积都同构于 $∅_E$
或属于 $C$,\pageoriginale{}
但 $\fprod{Z}{Z}{e}$ 和 $\fprod{Y}{Z}{e}$ 除外,
它们由 $C$ 的两个元素覆盖),
而显然 $\Ker(f, g)$ 不满足条件 c).

\item 给出一个例子, 使 $C$ 对纤维积稳定
(从而 $E$ 更是局部凝聚, \ie 局部代数拓扑斯),
但 $E$ 不是代数拓扑斯.
(取 d) 给出的例子,
以及 $E$ 中由对象 $Y$、$Z$、$T_{i}$ $(i > 0)$
及 $∅_{E}$ 构成的满子范畴 $C'$.)

\item 假定范畴 $C$ 有限.
证明拓扑斯 $E$ 为 Noether 的;
其 Noether (\ie 拟紧) 对象为那些对 $C$ 的每个对象 $X$,
$F(X)$ 都是\emph{有限}集的反变函子
$F : C^{\circ} \to (\Ens)$;
而 $E$ 上的纤维函子将 Noether 对象映为有限集
(\cf IV \hyperref[IV.exer7.6]{7.6}. h)).

\item 假定 $C$ 的每个态射 $f: X \to X$
都可分解为复合 $ip$, 其中 $i$ 为单态射, $p$ 有右逆.
设 $X$ 为 $C$ 的对象,
$I(X)$ 为 $C$ 意义下其子对象所成的集合, 看作有序集;
并设 $S(I(X))$ 为 $I(X)$ 的那些子集 $U$ 所成的集合,
使对两个元素 $X'$, $X'' \in I(X)$,
条件 $X'\in U$ 和 $X''\leq X'$ 蕴含 $X''\in U$.
证明 $X$ 在 $E$ 中的子对象有序集
同构于集合 $S(I(X))$ (按包含关系排序).
推出若 $I(X)$ 有限, 则 $X$ 为 $E$ 的预 Noether 对象.
特别地, 若 (除上述分解条件外)
$C$ 对纤维积 (\resp 有限投射极限) 稳定,
并且\pageoriginale{}对每个 $X \in \ob C$,
$X$ 在 $C$ 中的子对象集 $I(X)$ 有限,
则拓扑斯 $E$ 局部 Noether (\resp Noether).

\item 设 $C$ 为 $\in 𝒰$ 的有限集
(或者非空有限集) 范畴,
于是 $E = \widehat{C}$ 为增广单纯集
(\resp 通常的单纯集) 范畴.
证明 $E$ 是 Noether 拓扑斯. (利用 h).)
取 $C$ 的非本质常值 pro-对象 $(X_i)_{i \in I}$,
证明 $E$ 具有并不将 Noether 对象映为 Noether 对象
(\ie 有限集) 的纤维函子.

\item 考察拓扑斯 $E$ 的各性质之间的下列蕴含图:
\[
\xymatrix@R=.75cm{
**[r]\txt{E 为 Noether} \ar@{=>}[d] \ar@{=>}[r]& **[r]\txt{E
局部 Noether} \ar@{=>}[d] & \ar@{=>}[rdd]\\
**[r]\txt{E 凝聚} \ar@{=>}[r] & **[r]\txt{E
为代数的} \ar@{=>}[d]\\
 & **[r]\txt<9.5pc>{E 局部凝聚, 即\\E \raisebox{-1ex}{cohalg}
  生成 E} \ar@{=>}[d]& & \txt{E \raisebox{-1ex}{pre-Noeth}
  生成 E} \ar@{=>}[ldd] \\
 & **[r]\txt{E \raisebox{-1ex}{coh} 生成 E} \ar@{=>}[d]&\\
 & **[r]\txt{E \raisebox{-1ex}{qucpct} 生成 E}  &&.\\
}
\]
证明此图中的每个蕴含都严格,
而且所考察概念之间除上图中各蕴含的复合外没有其他蕴含.
这里 $E_{\coh}, E_{\text{qucpct}}$, $E_{\text{prénoeth}}$,
$E_{\text{cohalg}}$\pageoriginale{}
分别表示 $E$ 中由 $E$ 的凝聚、\resp 拟紧、
\resp 预 Noether、\resp 凝聚且代数的对象构成的满子范畴.
(可以只考虑形如 $\widehat{C}$ 的拓扑斯,
利用 d)、e)、f) 中所述结果.)

\item 回答下列问题 (本段作者不知道答案):
$E = \widehat{C}$ 能否局部 Noether,
而各 $\Hom(X, Y)$ (其中 $X$, $Y \in \ob C$) 不为有限集?
\end{enumerate}
\end{enonce*}

\begin{enonce*}[remark]{习题 2.18}\etiquette[2.18]{VI.exer2.18}
设 $E$ 是一个拓扑斯.
\begin{enumerate}
\renewcommand{\theenumi}{\alph{enumi}}
\renewcommand{\labelenumi}{\theenumi)}
\item 设 $X$ 为 $E$ 的预 Noether 对象.
证明 $X$ 同构于 $E$ 的有限个连通对象
(IV \hyperref[IV.subsubsec4.3.5]{4.3.5}) 的和.
(用反证法证明, $X$ 的递增分拆序列
(IV \hyperref[IV.exer8.7]{8.7}) 稳定.)

\item 设 $X$ 是 $E$ 的一个对象,
$(f_i : X_i \to X)_{i \in I}$ 为 $X$ 的一个覆盖族.
证明若每个 $X_i$ 都同构于 $E$ 的连通对象的一个和,
则 $X$ 也如此.
(归结为各 $X_i$ 连通的情形,
再归结为它们是 $X$ 的子对象的情形,
并在指标集 $I$ 上考察由关系 $X_i \cap X_j \neq \phi_E$
生成的等价关系 $R$.
证明若 $J = I/R$, 则 $X$ 是各
$X(i) = \displaystyle{\mathop{\Sup}_{i \in I}} X_i .)$ 的和.

\item 由 b) 推出,
若 $(X_i)_{i \in I}$ 表示 $E$ 的生成族,
则 $E$ 局部连通 (IV \hyperref[IV.exer8.7]{8.7}. $\ell$)),
当且\pageoriginale{}仅当每个 $X_i$ $(i \in I)$
都同构于 $E$ 的连通对象的和.

\item 由 a) 及 c) 推出,
若 $E$ 具有由预 Noether 对象构成的生成子范畴,
特别地, 若 $E$ 局部 Noether,
则 $E$ 局部连通,
从而更同构于一族连通拓扑斯
(IV \hyperref[IV.exer8.7]{8.7} e))
的和拓扑斯 (IV \hyperref[IV.exer8.7]{8.7} b)),
\ie 各终对象连通的拓扑斯.
(特别地, 对每个态射 $f : P \to E$,
其中 $P$ 是一个“非空连通且单连通”的拓扑斯
(IV \hyperref[IV.exer2.7.5]{2.7.5}),
可以给 $E$ 对应一个基本 pro-群 $\pi_1(E, f)$.
注意, 由 \textsc{Deligne} 的一个结果 ([9]),
若 $E$ 局部 Noether 且“非空”,
总能找到这样的 $f$, 其中 $P$ 为点拓扑斯.
\end{enumerate}
\end{enonce*}

% END INCLUDED SGA4_English_current_buildable_source_20260906/source/06/drafts/SGA4_VI_LINES_1679_2073_English.texfrag


% BEGIN INCLUDED SGA4_English_current_buildable_source_20260906/source/06/drafts/SGA4_VI_LINES_2074_2418_English.texfrag
\section{拓扑斯态射的有限性条件}\etiquette[3]{VI.sec3}

\begin{enonce*}{定义 3.1}\etiquette[3.1]{VI.def3.1}
设 $f:E'\to E$ 是一个拓扑斯态射.
若对 $E$ 的每个拟紧 (\resp 拟分离) 对象 $X$,
$f^{\ast}(X)$ 拟紧 (\resp 拟分离),
就称 $f$ \emph{拟紧 (\resp 拟分离)}.
若 $f$ 拟紧且拟分离, 就称 $f$ 凝聚.
\end{enonce*}

\setcounter{subsection}{1}
\setcounter{subsubsection}{0}
\subsubsection{}\etiquette[3.1.1]{VI.subsubsec3.1.1}
注意, 若 $f$ 具有上述任一性质,
则与 $f$ 同构的每个拓扑斯态射都具有相同性质.
还立即可知, 两个拟紧 (\resp 拟分离, \resp 凝聚) 拓扑斯态射的复合,
仍拟紧 (\resp 拟分离, \resp 凝聚).

\begin{enonce*}{命题 3.2}\etiquette[3.2]{VI.prop3.2}
沿用\pageoriginale{} \hyperref[VI.def3.1]{3.1} 的记号,
设 $(X_i)_{i\in I}$ 为 $E$ 的一个由拟紧
(\resp 凝聚) 对象构成的生成族.
在括号内情形, 假定 $E'$ 具有由拟紧对象构成的生成族.
则 $f$ 拟紧 (\resp 凝聚),
当且仅当对每个 $i\in I$, $f^{\ast}(X_i)$ 拟紧 (\resp 凝聚).
\end{enonce*}

必要性是显然的.
证明第一种情形的充分性,
注意对 $E$ 的每个拟紧对象 $X$,
存在有限满态射族 $X_i\to X$.
因此存在有限满族 $f^{\ast}(X_i)\to f^{\ast}(X)$,
其中各 $f^{\ast}(X_i)$ 由假设拟紧,
故 $f^{\ast}(X)$ 拟紧 (\hyperref[VI.prop1.3]{1.3}).
在第二种情形, 还须证明若 $X$ 是 $E$ 的拟分离对象,
则 $f^{\ast}(X)$ 拟分离.
对 $X$ 的假设可以表述为 (\hyperref[VI.cor1.17]{1.17}):
存在满态射族 $X_i\to X$,
使各 $\fprod{X_i}{X_j}{X}$ 拟紧.
于是有覆盖族 $f^{\ast}(X_i)\to f^{\ast}(X)$,
使纤维积
$\fprod{f^{\ast}(X_i)}{f^{\ast}(X_j)}{f^{\ast}(X)}$ 拟紧
(因为由第一种情形, $f^{\ast}$ 将拟紧对象映为拟紧对象).
由于各 $f^{\ast}(X_i)$ 凝聚,
结论仍由 \hyperref[VI.cor1.17]{1.17} 得到.

\begin{enonce*}{推论 3.3}\etiquette[3.3]{VI.cor3.3}
设 $C$ 与 $C'$ 是两个 $𝒰$-景,
其中纤维积可表示, 且每个覆盖族都有有限覆盖子族,
并设 $g:C\to C'$ 为景态射
\emph{(IV \hyperref[IV.subsubsec4.9.1]{4.9.1})}.
则 $g$ 定义的拓扑斯态射
$f:C^{\sim}\to{C'}^{\sim}$ 凝聚.
\end{enonce*}

事实上\pageoriginale{},
$\epsilon_C(C)$ (\resp $\epsilon_{C'}(C')$)
是 $C^{\sim}$ (\resp ${C'}^{\sim}$) 的生成族,
满足 \hyperref[VI.prop3.2]{3.2} 的相应条件
(\hyperref[VI.cor2.1.1]{2.1.1}).

\begin{enonce*}{命题 3.4}\etiquette[3.4]{VI.prop3.4}
设 $E$ 是一个局部凝聚拓扑斯
(\hyperref[VI.def2.3]{2.3}),
$f:X\to Y$ 是 $E$ 的箭头,
从而有诱导拓扑斯的态射
\emph{(IV \hyperref[IV.eq5.5.2]{(5.5.2)})} $E_{/X}\to E_{/Y}$.
这个拓扑斯态射拟紧 (\resp 拟分离, \resp 凝聚),
当且仅当 $f$ 为拟紧 (\resp 拟分离, \resp 凝聚) 态射.
\end{enonce*}

拟紧情形由定义立即得到 (无须对 $E$ 作任何条件);
拟分离情形正是 \hyperref[VI.cor2.8]{2.8};
凝聚情形由前两种情形合并得到.

\begin{enonce*}{命题 3.5}\etiquette[3.5]{VI.prop3.5}
设 $f:E'\to E$ 是局部凝聚拓扑斯之间的态射.
设 $(X_i)_{i\in I}$ 为 $E$ 中由凝聚代数对象构成的生成族.
下列条件等价:
\begin{itemize}
\item[(i)] 对 $E$ 的每个拟紧箭头 $u$, $f^{\ast}(u)$ 拟紧.

\item[(i bis)] 对每个箭头 $u:X_i\to X_j$, $f^{\ast}(u)$ 拟紧.

\item[(ii)] 对 $E'$ 的每个凝聚对象 $Y'$、
$E$ 的每个凝聚代数对象 $Y$,
以及每个箭头 $v:Y'\to f^{\ast}(Y)$,
相应拓扑斯态射
$$
f_v:{E'}_{/Y'}\to E_{/Y}
$$
凝聚; 这里该态射由从 $v$ 得到的局部化态射\pageoriginale{}
${E'}_{/Y'}\to{E'}_{/f^{\ast}(Y)}$
\emph{(IV \hyperref[IV.eq5.5.2]{(5.5.2)})}
与 $f$ 诱导的态射
${E'}_{/f^{\ast}(Y)}\to E_{/Y}$
\emph{(IV \hyperref[IV.eq5.10.1]{(5.10.1)})} 复合而成.

\item[(ii bis)] 与 \emph{(ii)} 相同的条件,
但只限于一组数据 $v_{\alpha}:{Y'}_{\alpha}\to f^{\ast}(Y_{\alpha})$,
使各 ${Y'}_{\alpha}$ 为代数的, 且覆盖 $E'$ 的终对象.

\item[(ii ter)] (给定 $E'$ 中由凝聚对象构成的生成族
$(X'_{i'})_{i'\in I'}$,
并对每个 $i'\in I'$ 给定一个 $i\in I$
及箭头 $v_{i'}:X'_{i'}\to f^{\ast}(X_i)$ 时.)
对每个 $i'\in I'$,
$v_{i'}$ 定义的拓扑斯态射 $E'_{/{X'}_{i'}}\to E_{/X_i}$ 凝聚.
\end{itemize}
\end{enonce*}

由于每个箭头 $X_i\to X_j$ 都拟紧,
显然 (i) $\Rightarrow$ (i bis).
反过来, 假定 (i bis) 成立, 证明 (i).
设 $u:X\to Y$ 为 $E$ 的拟紧箭头.
通过一个满族 $X_i\to Y$,
对 $u$ 的假设就是说各 $\fprod{X}{X_i}{Y}$ 拟紧
(\hyperref[VI.cor1.16]{1.16}),
\ie 对每个 $i$, 存在有限满态射族
$X_j\to\fprod{X}{X_i}{Y}$ (\hyperref[VI.prop1.3]{1.3}).
于是有相应的满族 $f^{\ast}(X_i)\to f^{\ast}(Y)$,
使对每个 $i$ 存在有限满族
$f^{\ast}(X_j)\to f^{\ast}(\fprod{X}{X_i}{Y})\simeq
\fprod{f^{\ast}(X)}{f^{\ast}(X_i)}{f^{\ast}(Y)}$;
由假设 (i bis), 其与到 $f^{\ast}(X_i)$ 的投影 $\pr_2$ 的复合,
是拟紧态射 $f^{\ast}(X_j)\to f^{\ast}(X_i)$.
由此 (\hyperref[VI.cor1.11]{1.11} (i)) 可知,
对每个 $i$,
$\pr_2:\fprod{f^{\ast}(X)}{f^{\ast}(X_i)}{f^{\ast}(Y)}$ 拟紧,
故 (\hyperref[VI.prop1.10]{1.10} (ii))
$f^{\ast}(u):f^{\ast}(X)\to f^{\ast}(Y)$ 拟紧.

还须\pageoriginale{}证明 (i) $\Rightarrow$ (ii)
及 (ii bis) $\Rightarrow$ (i);
蕴含 (ii) $\Rightarrow$ (ii bis) 是显然的,
而 (ii ter) 是 (ii bis) 的特殊情形.
蕴含 (i) $\Rightarrow$ (ii) 是立即的:
事实上, $E_{/Y}$ 的对象 $X$ 拟紧 (\resp 拟分离),
当且仅当其结构态射 $X\to Y$
(\resp 其对角态射 $X\to\fprod{X}{X}{Y}$) 拟紧
(\hyperref[VI.cor2.7]{2.7});
同样的拟紧性与拟分离性判据也适用于 $f^{\ast}_Y(X')$.
最后, 证明 (ii bis) $\Rightarrow$ (i).
设 $u:X\to Y$ 是 $E$ 的拟紧箭头;
须证明 $f^{\ast}(u)$ 拟紧.
由于各 $Y'_{\alpha}$ 覆盖 $E'$ 的终对象,
只须证明各 $f^{\ast}(u)|Y'_{\alpha}$ 拟紧
(\hyperref[VI.prop1.10]{1.10} (ii)).
这些态射正是各 $f^{\ast}_{v_\alpha}(u|Y_{\alpha})$,
故归结为证明下面的结果.

\begin{enonce*}{推论 3.6}\etiquette[3.6]{VI.cor3.6}
设 $f:E'\to E$ 是局部凝聚拓扑斯之间的凝聚态射.
则对 $E$ 的每个拟紧箭头 $u$, $f^{\ast}(u)$ 拟紧.
\end{enonce*}

由已证的 \hyperref[VI.prop3.5]{3.5} 的蕴含
(i bis) $\Rightarrow$ (i),
只须证明, 若 $u:X\to Y$ 是 $E$ 的箭头, $X,Y$ 凝聚,
则 $f^{\ast}(u)$ 是拟紧态射.
这立即来自对 $f$ 的假设,
它蕴含 $f^{\ast}(X)$ 与 $f^{\ast}(Y)$ 也凝聚.

\begin{enonce*}{定义 3.7}\etiquette[3.7]{VI.def3.7}
设 $f:E'\to E$\pageoriginale{} 是局部凝聚拓扑斯之间的态射.
若 $f$ 满足 \hyperref[VI.prop3.5]{3.5} 的等价条件,
就称其为\emph{局部凝聚的}.
\end{enonce*}

注意, “局部”指在源上局部.

\setcounter{subsection}{7}
\setcounter{subsubsection}{0}
\subsubsection{}\etiquette[3.7.1]{VI.subsubsec3.7.1}
这一概念仍然只依赖于拓扑斯态射 $f$ 的同构类,
且显然在拓扑斯态射的复合下稳定.
此外, 由 \hyperref[VI.prop3.5]{3.5},
若 $f$ 凝聚, 则它局部凝聚;
若 $E'$ 和 $E$ 凝聚, 则逆命题成立
(由 \hyperref[VI.prop3.5]{3.5} (ii bis)).
判据 \hyperref[VI.prop3.5]{3.5} (ii bis) 立即给出下列判据:

\begin{enonce*}{推论 3.7}\etiquette[3.7]{VI.cor3.7}
设 $f:E'\to E$ 是局部凝聚拓扑斯之间的态射,
$(Y_i)_{i\in I}$ 与 $({Y'}_i)_{i\in I}$
为 $E$ 及 $E'$ 的两个对象族;
并对每个 $i\in I$, 设 $v_i:Y'_i\to f^{\ast}(Y_i)$ 为态射,
从而有拓扑斯态射
$$
f_{v_i}:E'_{/Y'_i}\to E_{/Y_i}{\quoi}.
$$
假定各 ${Y'}_i$ 覆盖 $E'$ 的终对象.
则 $f$ 局部凝聚,
当且仅当对每个 $i\in I$, $f_{v_i}$ 局部凝聚.
\end{enonce*}

特别地, 将此用于诱导拓扑斯的恒等态射, 得到:

\begin{enonce*}{推论 3.8}\etiquette[3.8]{VI.cor3.8}
设 $E$ 是一个局部凝聚拓扑斯, $v:X\to Y$ 是 $E$ 的箭头.
则 $v$ 定义的诱导拓扑斯态射 $E_{/X}\to E_{/Y}$ 局部凝聚.
\end{enonce*}

\begin{enonce*}{注 3.9}\etiquette[3.9]{VI.rem3.9}
看来\pageoriginale{}实际遇到的所有代数拓扑斯态射都局部凝聚
(因此“局部凝聚”一词大概不会经常用到).
等价地, 实际遇到的凝聚拓扑斯之间的态射都是凝聚的.
不过, 可以构造凝聚拓扑斯之间的非凝聚态射;
更具体地说, 可以构造从点拓扑斯
\emph{(IV \hyperref[IV.subsec2.2]{2.2})}
到某个 Noether 拓扑斯 $E$ 的非凝聚态射,
\emph{\cf \hyperref[VI.exer2.17]{2.17} i)}.
\end{enonce*}

\begin{enonce*}[remark]{例 3.10}\etiquette[3.10]{VI.exem3.10}
(\emph{与概形相伴的拓扑斯}.)
回到范畴 $(\Sch)$ 配备 \SGA IV 6.3 的各拓扑 $T_i$
这一例子 \hyperref[VI.subsec1.22]{1.22}.
对每个概形 $X$,
设 $X_{T_i}$ 为诱导景 $(\Sch)_{/X}$ 定义的 $𝒱$-拓扑斯.
概形态射 $f:X\to Y$ 因而定义拓扑斯态射
$f_{T_i}:X_{T_i}\to Y_{T_i}$;
由 \hyperref[VI.prop3.4]{3.4} 及 \hyperref[VI.subsec1.22]{1.22} 可知,
这个拓扑斯态射拟紧 (\resp 拟分离, \resp 凝聚),
当且仅当概形态射 $f$ 在 \EGA IV 1 的通常意义下
拟紧 (\resp 拟分离, \resp 凝聚).
无论如何, 由 \hyperref[VI.cor3.8]{3.8}, $f_{T_i}$ 局部凝聚.

还可以如 \hyperref[VI.subsubsec1.22.1]{1.22.1}
给概形 $X$ 对应 $𝒰$-拓扑斯 $\Top(X)$ 与 $\Top(X_{\text{ét}})$,
并给每个概形态射 $f:X\to Y$
对应相应拓扑斯之间的态射 $\Top(f)$ 与 $\Top(f_{\text{ét}})$.
以 $T(f)$ 表示这两个拓扑斯态射之一.\pageoriginale{}
由 \hyperref[VI.prop3.2]{3.2} 立即可验证,
此拓扑斯态射总局部凝聚,
并且它拟紧 (\resp 拟分离, \resp 凝聚),
当且仅当概形态射 $f$ 在通常意义下拟紧
(\resp 拟分离, \resp 凝聚).

这些观察再次表明,
本节引入的术语与概形理论的通用术语相容,
因而即使最不情愿的读者也应当接受它们.
\end{enonce*}

\begin{enonce*}[remark]{习题 3.11}\etiquette[3.11]{VI.exer3.11}
(\emph{凝聚拓扑斯与预拓扑斯}.)
\begin{enumerate}
\renewcommand{\theenumi}{\alph {enumi}}
\renewcommand{\labelenumi}{\theenumi)}
\item $𝒰$-预拓扑斯（pretopos）
(或简称预拓扑斯) 是满足下列条件的范畴 $C$:
\begin{enumerate}
\renewcommand{\theenumii}{\arabic{enumii}}
\renewcommand{\labelenumii}{\theenumii)}
\item $C$ 中有限投射极限可表示.

\item $C$ 中有限余积可表示, 且它们不交且泛.

\item $C$ 中的等价关系有效, 且 $C$ 中的每个满态射泛有效.

\item $C$ 等价于某个 $\in 𝒰$ 的范畴.
\end{enumerate}
证明若 $E$ 是凝聚 $𝒰$-拓扑斯,
则 $E$ 中由其凝聚对象构成的满子范畴 $E_{\coh}$ 为 $𝒰$-预拓扑斯,
而包含函子 $E_{\coh}\to E$ 左正合,
与\pageoriginale{}有限余积交换,
并与过渡到等价关系的商交换.
(利用 \hyperref[VI.subsubsec1.5.3]{1.5.3}、
\hyperref[VI.cor1.15]{1.15} 及 \hyperref[VI.cor1.17.1]{1.17.1}.)
证明 $E$ 在 $C$ 上诱导的拓扑为“预典范”拓扑,
\ie 其覆盖族 $X_i\to X$
恰为那些具有关于 $C$ 的典范拓扑覆盖的\emph{有限}子族的族.
因此, 可以在等价意义下由预拓扑斯 $C=E_{\coh}$
恢复 $E$ 为拓扑斯 $C^{\sim}$
(其中 $C$ 配备其预典范拓扑).

\item 设 $C$ 是 $𝒰$-预拓扑斯.
给 $C$ 配备预典范拓扑.
证明 $C$ 的每个态射 $f$ 都分解为 $f''f'$,
其中 $f'$ 为满态射, $f''$ 为单态射.
证明拓扑斯 $E=C^{\sim}$ 凝聚,
且典范函子 $\epsilon:C\to E$
诱导 $C$ 到 $E$ 的子范畴 $E_{\coh}$ 的等价.
(先证明 $\varepsilon$ 为全忠实左正合函子,
与有限余积以及过渡到等价关系的商交换;
再证明 $C$ 的一个对象在 $E$ 中的凝聚子对象属于 $C$ 的本质像.)
因此, 可以在范畴等价意义下,
由其相伴 $𝒰$-拓扑斯 $E=C^{\sim}$ 恢复 $𝒰$-预拓扑斯 $C$.

\item 设 $C$ 和 $C'$ 是两个 $𝒰$-预拓扑斯.
证明函子 $\varphi:C'\to C$
为从 $C$ 到 $C'$ 的景态射 (关于预典范拓扑),
当且仅当 $\varphi$ 左正合,
且与有限余积以及过渡到等价关系的商交换.

\item 沿用\pageoriginale{} c) 的记号,
假定 $C$ 和 $C'$ 如 a) 中那样与两个凝聚 $𝒰$-拓扑斯 $E,E'$ 相伴.
证明典范函子 (IV \hyperref[IV.subsubsec4.9.3]{4.9.3})
$$
\Mmorsite(C,C')\to\Hhomtop(E,E')
$$
从其源 (在 c) 中明确描述) 到其靶中
由\emph{凝聚}拓扑斯态射 $E\to E'$ 构成的满子范畴诱导等价.
给每个凝聚拓扑斯态射 $f:E\to E'$
对应由 $f^{\ast}$ 诱导的函子 $E'_{\coh}=C'\to E_{\coh}=C$,
即可得到一个拟逆.

\item 设 $C$ 是一个 $𝒰$-预拓扑斯, $E=C^{\sim}$.
直接用 $C$ 的正合性表述 \hyperref[VI.cor1.25]{1.25} 的等价条件
(参看 \hyperref[VI.exer1.28]{1.28} b).
证明 $C$ 的对象 $X$ 在 $E$ 中为 Noether 的,
当且仅当 $X$ 在 $C$ 中的每个递增子对象序列稳定;
从而 $E$ 为 Noether 的,
当且仅当 $C$ 的每个对象满足上述条件.
\end{enumerate}
\end{enonce*}

\begin{enonce*}[remark]{习题 3.12}\etiquette[3.12]{VI.exer3.12}
若存在有限范畴 $C$, 使拓扑斯 $E$ 等价于 $\widehat{C}$,
就称 $E$ 为\emph{有限}拓扑斯.
\begin{enumerate}
\renewcommand{\theenumi}{\alph{enumi}}
\renewcommand{\labelenumi}{\theenumi)}
\item 拓扑斯 $E$ 有限,
当且仅当 $E$ 有足够多的本质点
(IV \hyperref[IV.exer7.6]{7.6} b)),
且 $E$ 的本质点范畴 $\Ppointess(E)$ 等价于有限范畴.
(利用 IV \hyperref[IV.exer7.6]{7.6} d) 及 h).)

\item 假定\pageoriginale{} $E$ 有限.
证明 $E$ 的每个点都本质.
证明 2-函子 $E\mapsto\Ppoint(E)$ 及 $C\mapsto\widehat{C}$
在有限拓扑斯 $E$ 的 2-范畴
与等价于有限范畴的 Karoubi 范畴
(IV \hyperref[IV.exer7.5]{7.5} a)) $C$ 的 2-范畴之间建立 2-等价,
且有限拓扑斯之间的每个态射都本质
(IV \hyperref[IV.exer7.6]{7.6} a)).
(利用 a) 及 IV \hyperref[IV.exer7.6]{7.6} h).)

\item 假定 $E=\widehat{C}$, $C$ 等价于有限范畴,
$F$ 是凝聚拓扑斯.
回忆 (IV \hyperref[IV.eq4.6.3.1]{(4.6.3.1)}),
拓扑斯态射 $f:E\to F$ 对应于函子 $u:C\to\Ppoint(F)$.
证明 $f$ 凝聚,
当且仅当 $f$ 将 $E$ 的每个凝聚点
(\hyperref[VI.def3.1]{3.1}) 映为 $F$ 的凝聚点;
或者等价地, 当且仅当对每个 $X\in\ob C$,
$u(X)$ 为 $F$ 的\emph{凝聚}点.
(利用 \hyperref[VI.prop3.2]{3.2}
及 \hyperref[VI.exer2.17]{2.17} g).)
推出从有限拓扑斯到有限拓扑斯的每个态射 $f:E\to F$ 都凝聚.

\item 设 $X$ 是一个清醒拓扑空间
(IV \hyperref[IV.subsubsec4.2.1]{4.2.1}).
拓扑斯 $\Top(X)$ (IV \hyperref[IV.subsec2.1]{2.1}) 有限,
当且仅当 $X$ 为有限集.
(利用 IV \hyperref[IV.subsubsec7.1.6]{7.1.6}.)
\end{enumerate}
\end{enonce*}


% END INCLUDED SGA4_English_current_buildable_source_20260906/source/06/drafts/SGA4_VI_LINES_2074_2418_English.texfrag


% BEGIN INCLUDED SGA4_English_current_buildable_source_20260906/source/06/drafts/SGA4_VI_LINES_2419_2936_English.texfrag
\section{粘合得到的拓扑斯中的有限性条件}\etiquette[4]{VI.sec4}

本节在本讨论班后文中不会用到.

\subsection{}\etiquette[4.1]{VI.subsec4.1}
设 $E$ 是一个拓扑斯, $U$ 是 $E$ 的一个开集
(\ie $E$ 的终对象的子对象),
并考察 $E$ 的相应开子拓扑斯与闭子拓扑斯\pageoriginale{} (IV 9)
$$
E'= E_{/U}, E'' = E_{/U}{\quoi}.
$$
以
$$
j : E'\to E, i : E''\to E
$$
表示典范拓扑斯态射.
回忆 (\hyperref[VI.prop3.4]{3.4}),
$j$ 为拟紧 (\resp 凝聚) 拓扑斯态射的充要条件是,
$U$ 到 $E$ 的终对象的包含
$$
j : U \to e
$$
(也记为 $j$) 为 $E$ 中的拟紧 (\resp 凝聚) 态射;
此外, $j:E'\to E$ 总拟分离
(因为 $j:U\to e$ 拟分离 (\hyperref[VI.cor1.8.1]{1.8.1})).
我们拟给出拓扑斯 $E''$ 凝聚
以及拓扑斯态射 $i:E''\to E$ 凝聚的判据.
先注意:

\begin{enonce*}{命题 4.2}\etiquette[4.2]{VI.prop4.2}
沿用 \hyperref[VI.subsec4.1]{4.1} 的记号,
包含态射 $i:E''\to E$ 拟紧,
\ie 对 $E$ 的每个拟紧对象 $X$,
$E''$ 的对象 $i^{\ast}(X)$ 拟紧;
此外, 若 $X$ 为预 Noether
(\hyperref[VI.def1.30]{1.30}),
则 $i^{\ast}(X)$ 为预 Noether.
\end{enonce*}

这来自\hyperref[VI.def1.1]{定义 1.1} 以及下述引理.

\begin{enonce*}{引理 4.2.1}\etiquette[4.2.1]{VI.lem4.2.1}
映射\pageoriginale{} $Y\mapsto i^{\ast}(Y)$
在 $X$ 中包含 $X$ 的子对象 $X_U=X\times U$ 的子对象集
与 $i^{\ast}(X)$ 的子对象集之间诱导有序集同构.
\end{enonce*}

注意, 由于 $i_{\ast}$ 保守
(因为它全忠实) 且左正合
(特别地, 它与纤维积交换),
立即可知 $E''$ 中的态射 $u:Y\to Z$ 为单态射,
当且仅当 $i_{\ast}(u):i_{\ast}(Y)\to i_{\ast}(Z)$ 为单态射.
因此, 用 $i_{\ast}$ 将 $E''$ 等同于 $E$ 中
由满足 IV \hyperref[IV.subsec9.7]{9.7} 1) 的等价条件的对象 $Z$
构成的满子范畴后,
$E''$ 的对象 $Z$ 在 $E''$ 中的子对象,
等同于 $Z$ 在 $E$ 中属于 $E''$ 的子对象 $Y$,
或等价地, 包含 $Z_U\xrightarrow{\sim}U$ 的子对象.
对形如由融合和 IV \hyperref[IV.subsubsec9.5.1]{(9.5.1)} 给出的
$i_{\ast}i^{\ast}(X)=X_{\complement U}$ 的 $Z$,
$$
i_{\ast} i^{\ast} (X) = X_{\complement U} = X \amalg_{X_U} U {\quoi},
$$
给定后者在 $E$ 中的子对象 $Y$,
等价于给定由 $X$ 的子对象 $Y_1$ 及 $U$ 的子对象 $Y_2$ 构成的一个对,
条件是这些子对象在 $X_U$ 中的逆像相同.
此时 $Y$ 包含 $U$ 的条件意味着 $Y_2=U$,
而对 $Y_1$ 剩下的条件就是它包含 $X_U$.
由此在 $i^{\ast}(X)$ 于 $E''$ 中的子对象集
与 $X$ 的包含 $X_U$ 的子对象集之间建立了双射.
显然\pageoriginale{}, 这是有序集同构,
其逆映射确为 \hyperref[VI.lem4.2.1]{4.2.1} 中所述映射.

\begin{enonce*}{推论 4.3}\etiquette[4.3]{VI.cor4.3}
\begin{enumerate}
\renewcommand{\theenumi}{\alph{enumi}}
\renewcommand{\labelenumi}{\theenumi)}
\item 设 $X$ 是 $E$ 的对象.
$X$ 拟紧的充分条件是 $i^{\ast}(X)$ 与 $j^{\ast}(X)$ 拟紧;
若 $j:U\to e$ 拟紧, 则此条件也必要.

\item 设 $Y$ 是 $E''$ 的对象.
$Y$ 拟紧的充分条件是 $i_{\ast}(Y)$ 拟紧;
若 $U$ 拟紧, 则此条件也必要.
\end{enumerate}
\end{enonce*}

\noindent
\underline{证明}
\begin{enumerate}
\renewcommand{\theenumi}{\alph{enumi}}
\renewcommand{\labelenumi}{\theenumi)}
\item 充分性来自\hyperref[VI.def1.1]{定义 1.1}
及函子对 $(i^{\ast},j^{\ast})$ 保守这一事实
(IV \hyperlink{sga4.exp.IV.proposition.9.4.1.item.c}{9.11} 3)).
必要性来自 $i$ 与 $j$ 拟紧这一事实
(由 \hyperref[VI.prop4.2]{3.14} 及 $j$ 拟紧的假设).

\item 由于 $Y\simeq i^{\ast}i_{\ast}(Y)$,
充分性来自 \hyperref[VI.prop4.2]{4.2}.
必要性来自 a) 的充分性,
只须考虑 $i^{\ast}i_{\ast}(Y)\simeq Y$
及 $j^{\ast}i_{\ast}(Y)\simeq U$.
\end{enumerate}

\begin{enonce*}{推论 4.4}\etiquette[4.4]{VI.cor4.4}
设 $Y$ 是 $E''$ 的对象.
若 $U$ 拟紧, 或 $E$ 有一个由拟紧对象构成的生成族,
则 $Y$ 拟分离 (\resp 凝聚) 的\emph{充分}条件是
$i_{\ast}(Y)$ 拟分离 (\resp 凝聚).
若 $U$ 拟分离 (\resp 凝聚) 且 $j:U\to e$ 为拟紧态射,
则 $Y$ 拟分离 (\resp 凝聚) 的必要条件是
$i_{\ast}(Y)$ 拟分离 (\resp 凝聚).
\end{enonce*}

括号内\pageoriginale{}情形来自非括号内情形,
只须考虑 \hyperref[VI.cor4.3]{4.3} b).
假定 $i_{\ast}(Y)$ 拟分离,
证明在所述两个假设之一成立时 $Y$ 拟分离.
因此须证明, 对 $Y$ 上的两个拟紧对象 $Y'$ 与 $Y''$,
纤维积 $\fprod{Y'}{Y''}{Y}$ 拟紧.
当假定 $U$ 拟紧时,
注意只须证明 $i_{\ast}(\fprod{Y'}{Y''}{Y})$ 拟紧
(\hyperref[VI.cor4.3]{4.3} b));
这来自 $i_{\ast}(Y')$ 与 $i_{\ast}(Y'')$ 拟紧
(\hyperref[VI.cor4.3]{4.3} b)) 这一事实,
这里用到 $U$ 的拟紧性、$i_{\ast}$ 与纤维积交换,
以及 $i_{\ast}(Y)$ 拟分离的假设.
当假定 $E$ 具有由拟紧对象构成的生成族时,
$i_{\ast}(Y')$ 是其拟紧子对象 ${X'}_{\alpha}$ 的滤过归纳极限;
故 $Y'\simeq i^{\ast}i_{\ast}(Y')$
是子对象 $i^{\ast}({X'}_{\alpha})$ 的滤过归纳极限;
由于 $Y'$ 拟紧, 它等于其中一个 $i^{\ast}({X'}_{\alpha})$.
同样, $Y''$ 形如 $i^{\ast}({X''}_{\beta})$,
其中 ${X''}_{\beta}$ 是 $i_{\ast}(Y'')$ 的拟紧子对象.
但此时
$\fprod{Y'}{Y''}{Y}=
i^{\ast}(\fprod{X'_{\alpha}}{X''_{\beta}}{i_{\ast}Y})$;
而由 $i_{\ast}(Y)$ 拟分离的假设,
$\fprod{X'_{\alpha}}{X''_{\beta}}{i_{\ast}(Y)}$ 拟紧,
故由 \hyperref[VI.prop4.2]{4.2},
$\fprod{Y'}{Y''}{Y}$ 也拟紧.

反过来, 假定 $U$ 拟分离且 $j:U\to e$ 拟紧,
并假定 $Y$ 拟分离.
证明 $i_{\ast}(Y)$ 也拟分离,
\ie 对 $i_{\ast}(Y)$ 上的两个拟紧对象 $X',X''$,
纤维积 $\fprod{X'}{X''}{i_{\ast}(Y)}$ 拟紧.
由 \hyperref[VI.cor4.3]{4.3} a),
只须证明其在 $i^{\ast}$ 下的像拟紧\pageoriginale{}
(这来自 \hyperref[VI.prop4.2]{4.2} 及 $Y$ 拟分离的假设),
\emph{并且}其在 $j^{\ast}$ 下的像拟紧.
后者为 $j^{\ast}(X')\times j^{\ast}(X'')$,
它确实拟紧,
因为 $j^{\ast}(X')$ 与 $j^{\ast}(X'')$ 拟紧
($j$ 拟紧), 而 $U$ 拟分离.

\begin{enonce*}{推论 4.5}\etiquette[4.5]{VI.cor4.5}
假定 $E$ 具有一个由拟紧
(\emph{\resp} 预 Noether) 对象构成的生成子范畴.
则 $E''$ 也如此.
\end{enonce*}

这来自 \hyperref[VI.prop4.2]{4.2} 及下述引理.

\begin{enonce*}{引理 4.5.1}\etiquette[4.5.1]{VI.lem4.5.1}
若 $(X_{\alpha})_{\alpha}$ 是 $E$ 的生成族,
则 $(i^{\ast}(X_{\alpha}))_{\alpha}$ 是 $E''$ 的生成族.
\end{enonce*}

事实上, 这意味着 $E''$ 上的函子族
$$
Y \to \Hom (i^{\ast}(X_{\alpha}), Y)
$$
保守.
但 $\Hom(i^{\ast}(X_{\alpha}),Y)\simeq\Hom(X_{\alpha},i_{\ast}(Y))$,
所以只须利用函子 $i_{\ast}$ 保守这一事实.

\begin{enonce*}{命题 4.6}\etiquette[4.6]{VI.prop4.6}
沿用 \hyperref[VI.subsec4.1]{4.1} 的记号.
\begin{enumerate}
\renewcommand{\theenumi}{\alph{enumi}}
\renewcommand{\labelenumi}{\theenumi)}
\item 若 $E$ 具有由凝聚对象构成的生成族,
则闭子拓扑斯 $E''$ 也如此,
且包含态射 $i:E''\to E$ 凝聚.

\item 假定 $E$ 中的态射 $j:U\to e$ 拟紧.
若\pageoriginale{} $E$ 局部凝聚
(\emph{\resp} 凝聚,
\emph{\resp} 拟分离,
\emph{\resp} 局部 Noether,
\emph{\resp} Noether,
\emph{\resp} 局部完美,
\emph{\resp} 完美),
则 $E''$ 也如此, 且包含态射 $i:E''\to E$ 凝聚.
\end{enumerate}
\end{enonce*}

\begin{enumerate}
\renewcommand{\theenumi}{\alph{enumi}}
\renewcommand{\labelenumi}{\theenumi)}
\item 第一个断言将来自第二个断言及 \hyperref[VI.lem4.5.1]{4.5.1}.
无论如何, \hyperref[VI.cor4.5]{4.5}
蕴含 $E''$ 具有一个由拟紧对象构成的生成族,
故 \hyperref[VI.prop3.2]{3.2} 可应用,
并表明只须验证对 $E$ 的每个凝聚对象 $X$,
$E''$ 的对象 $i^{\ast}(X)$ 凝聚.
利用开集的补拓扑斯的形成与局部化相容这一事实
(IV \hyperref[IV.cor9.4.3]{9.13} b)),
归结为 $X$ 为 $E$ 的终对象的情形,
从而归结为证明 $E''$ 的终对象凝聚.
现在 \hyperref[VI.cor4.4]{4.4} 可应用,
故只须证明 $i_{\ast}(e_{E''})=e_E$ 凝聚, 而这确实成立.

\item 在 b) 的每一个假设下,
$E$ 都具有由凝聚对象构成的生成族,
故由 a) 已知 $i:E''\to E$ 凝聚.
此外, \hyperref[VI.lem4.5.1]{4.5.1} 因而蕴含,
其一个生成子范畴为凝聚对象构成的满子范畴 ${E''}_{\coh}$.
假定 $E$ 凝聚, 证明 $E''$ 也凝聚,
\ie (\hyperref[VI.subsubsec2.4.5]{2.4.5})
$E''$ 的子范畴 ${E''}_{\coh}$ 对有限投射极限稳定,
已知 $E$ 的子范畴 $E_{\coh}$ 具有此性质.
这立即来自括号内判据 \hyperref[VI.cor4.4]{4.4}
(它可以应用, 因为现在 $U$ 凝聚,
这是因为 $e$\pageoriginale{} 凝聚而 $j:U\to e$ 拟紧),
只须考虑 $i_{\ast}$ 左正合.
经局部化, 再次使用 IV \hyperref[IV.cor9.4.3]{9.13} b),
可知若 $E$ 局部凝聚, 则 $E''$ 也如此.
若 $E$ 拟分离, 则 $E''$ 也拟分离:
事实上, 由 \hyperref[VI.cor4.4]{4.4},
因 $E$ 的终对象拟分离, $E''$ 的终对象拟分离;
还须证明 $E''$ 的两个拟分离对象的积拟分离
(已知 $E$ 中如此);
这再次来自判据 \hyperref[VI.cor4.4]{4.4}
以及 $i_{\ast}$ 与积交换这一事实
(注意 $e$ 的子对象 $U$ 拟分离, 因为 $e$ 拟分离,
故 \hyperref[VI.cor4.4]{4.4} 可应用).
\end{enumerate}

由于拓扑斯局部 Noether (\resp Noether),
当且仅当它局部凝聚 (\resp 凝聚)
且具有由预 Noether 对象构成的生成族
(\hyperref[VI.prop2.10]{2.10} (iii bis)),
由前面的讨论及 \hyperref[VI.cor4.5]{4.5} 可知,
若 $E$ 局部 Noether (\resp Noether), 则 $E''$ 也如此.
再假定 $E$ 完美, 证明 $E''$ 完美.
既已知道它凝聚,
只须验证它满足判据 \hyperref[VI.cor1.25]{1.25} (i),
\ie $E''$ 的每个对象 $Y$ 都是凝聚对象的滤过归纳极限,
已知 $E$ 中的类似陈述成立.
由于 $i_{\ast}(Y)$ 是凝聚对象 $X_{\alpha}$ 的滤过归纳极限,
$Y\simeq i^{\ast}i_{\ast}(Y)$
是凝聚对象 $i^{\ast}(X_{\alpha})$ 的滤过归纳极限.
经局部化 (IV \hyperref[IV.cor9.4.3]{9.13} b)),
可知若 $E$ 局部完美, 则 $E''$ 也如此.
这就完成 \hyperref[VI.prop4.6]{4.6} 的证明.

\begin{enonce*}[remark]{注 4.6.1}\etiquette[4.6.1]{VI.rem4.6.1}
事实上\pageoriginale{},
\hyperref[VI.prop4.6]{4.6} b) 中 $j:U\to e$ 拟紧的假设并不必要,
或许“$E$ 拟分离”情形除外.
确实, 由前面的证明,
只须证明若 $E$ 凝聚, 则 $E''$ 凝聚.
$U$ 是其拟紧子对象 $U_{\alpha}$ 的滤过归纳极限;
而在 \hyperref[VI.sec7]{7} 中将验证,
$E''$ 此时等同于闭子拓扑斯 $E''_i$ 的投射极限拓扑斯
(取 \hyperref[VI.sec7]{7} 的意义),
其中这些子拓扑斯是各 $E_{/U_i}$ 的补,
由 \hyperref[VI.prop4.6]{4.6} b),
它们为凝聚拓扑斯, 且转移态射凝聚.
于是可知 $E''$ 为凝聚拓扑斯 (\hyperref[VI.sec7]{7}).
\end{enonce*}

\begin{enonce*}{推论 4.7}\etiquette[4.7]{VI.cor4.7}
假定 $E$ 中由凝聚对象构成的子范畴 $E_{\coh}$ 生成,
并设 $X$ 是 $E$ 的一个对象.
\begin{enumerate}
\renewcommand{\theenumi}{\alph{enumi}}
\renewcommand{\labelenumi}{\theenumi)}
\item 若 $X$ 拟分离, 则 $i^{\ast}(X)$ 与 $j^{\ast}(X)$ 拟分离;
若 $j:U\to e$ 拟紧, 此条件也充分.

\item 假定 $j:U\to e$ 拟紧.
则 $X$ 凝聚, 当且仅当 $i^{\ast}(X)$ 与 $j^{\ast}(X)$ 凝聚.

\item 假定 $E$ 凝聚且 $j:U\to e$ 拟紧.
则 $E$ 完美 (\hyperref[VI.def2.9.1]{2.9.1}),
当且仅当 $E'$ 和 $E''$ 完美.
\end{enumerate}
\end{enonce*}

\begin{enumerate}
\renewcommand{\theenumi}{\alph{enumi}}
\renewcommand{\labelenumi}{\theenumi)}
\item 在 \hyperref[VI.subsec4.1]{4.1} 中已指出 $j:E'\to E$ 拟分离,
由 \hyperref[VI.prop4.6]{4.6} a) 知 $i$ 也如此,
这就给出必要性.
证明充分性, 设 $X'$ 和 $X''$ 是 $X$ 上的拟紧对象.
须证明 $\fprod{X'}{X''}{X}$ 拟紧;
为此\pageoriginale{}, 只须证明其在 $i^{\ast}$ 及 $j^{\ast}$ 下的像拟紧
(\hyperref[VI.cor4.3]{4.3} a)).
由于 $i$ 及 $j$ 拟紧
(由 \hyperref[VI.prop4.2]{4.2} 及 $j:U\to e$ 拟紧的假设),
结论来自 $i^{\ast}$ 及 $j^{\ast}$ 与纤维积交换,
以及 $i^{\ast}(X)$ 与 $j^{\ast}(X)$ 拟分离的假设.

\item 合并 a) 与 \hyperref[VI.cor4.3]{4.3} a) 即得.

\item 必要性已见到 (\hyperref[VI.prop4.6]{4.6} b)).
充分性来自以下事实:
由于 $E$ 凝聚, $E'$ 与 $E''$ 凝聚
(\hyperref[VI.prop4.6]{4.6} b)),
故对所考察的任一拓扑斯,
完美意味着其凝聚对象范畴对有限 $\varinjlim$ 稳定.
于是由 b) 得到结论.
\end{enumerate}

\begin{enonce*}[remark]{习题 4.8}\etiquette[4.8]{VI.exer4.8}
\begin{enumerate}
\renewcommand{\theenumi}{\alph{enumi}}
\renewcommand{\labelenumi}{\theenumi)}
\item 解决下列问题
(写下这段话的人惭愧地承认不知道答案):
沿用 \hyperref[VI.subsec4.1]{4.1} 的记号,
若 $E$ 拟分离 (\resp 代数), $E''$ 是否也如此?

\item 假定 $E$ 等价于形如 $\widehat{C}$ 的拓扑斯,
其中 $C$ 是等价于某个 $\in 𝒰$ 的范畴的范畴.
利用 \hyperref[VI.exer2.17]{2.17} c) 及 IV 9.24,
直接证明若 $E$ 局部凝聚
(\resp 凝聚, \resp 代数, \resp 拟分离,
\resp 局部 Noether, \resp Noether),
则 $E''$ 也如此.
\end{enumerate}
\end{enonce*}

\setcounter{subsection}{8}
\subsection{}\etiquette[4.9]{VI.subsec4.9}
沿用 \hyperref[VI.subsec4.1]{4.1} 的记号,
并回忆 (IV \hyperref[IV.theo9.5.4]{9.16}),
给定由一个拓扑斯 $E$ 和 $E$ 的一个开集 $U$ 构成的情形 $(E,U)$,
本质上等价于给定一个三元组 $(E',E'',f)$,
其中 $E'$ 和 $E''$ 为\pageoriginale{}拓扑斯,
\begin{equation*}
f : E'\to E''\tag{4.9.1}
\end{equation*}\etiquette[4.9.1]{VI.eq4.9.1}
为“粘合函子”, \ie \emph{左正合且可及}的函子;
从 \hyperref[VI.subsec4.1]{4.1} 中的 $(E,U)$ 出发,
相伴粘合函子给出为
\begin{equation*}
  f = i^{\ast} j_{\ast}{\quoi}.\tag{4.9.2}
\end{equation*}\etiquette[4.9.2]{VI.eq4.9.2}

我们拟用数据 $E'$、$E''$ 及 $f$
表达 $E$ 为凝聚 (\resp 完美) 拓扑斯且 $U$ 为拟紧对象这一事实,
或者等价地, 表达 $E$ 与 $E'$ 凝聚这一事实.
已经知道, 这蕴含 $E''$ 也凝聚
(\hyperref[VI.prop4.6]{4.6} b)),
因此问题变为:
给定两个\emph{凝聚}拓扑斯 $E'$ 与 $E''$
及粘合函子 $f$ \eqhyperref[VI.eq4.9.1]{4.9.1},
$f$ 须满足什么条件, 才能使粘合拓扑斯 $E$ 凝聚 (\resp 完美)?
还知道 (\hyperref[VI.cor4.7]{4.7} c)),
$E$ 完美, 当且仅当 $E$ 凝聚且 $E'$ 与 $E''$ 完美;
故问题的括号内情形归结为非括号内情形.
不过应当指出, 若已假定 $E'$ 完美,
后者的解答 (\hyperref[VI.prop4.10]{4.10}) 更为优美.

\setcounter{subsection}{9}
\setcounter{subsubsection}{2}
\subsubsection{}\etiquette[4.9.3]{VI.subsubsec4.9.3}

先注意, 若 $E$ 凝聚,
则 $E$ 中由 $E$ 的凝聚对象
$X=(X',X'',u:X''\to f(X'))$ 构成的满子范畴 $E_{\coh}$,
可恢复为 $E$ 的子范畴 $E_0$,
其对象为\pageoriginale{}使
$X'=j^{\ast}(X)$ 与 $X''=i^{\ast}(X)$ 凝聚的 $X$
(\hyperref[VI.cor4.7]{4.7} b)).
因此, $E$ 凝聚的必要条件是前述子范畴 $E_0$ 生成.
此条件也充分,
因为由 \hyperref[VI.cor4.3]{4.3} a),
$E_0$ 由拟紧对象构成;
而且 $E_0$ 显然对有限投射极限稳定,
故由 \hyperref[VI.subsubsec2.4.5]{3.4.5} 得到结论.

再回忆 (IV \hyperref[IV.subsec9.6]{9.18}),
给定粘合函子 \eqhyperref[VI.eq4.9.1]{4.9.1},
等价于给定 $E''$ 上取值于 $\Pro(E')^0$ 的一个\emph{层}:
\begin{equation*}
G \in \Ffaisc (E'', \Pro(E')^{0}), G : {E''}^{0} \to \Pro({E'}^{0})
\tag{4.9.4}{\quoi},
\end{equation*}\etiquette[4.9.4]{VI.eq4.9.4}
或等价于给定一个\emph{与归纳极限交换}的函子 $g=G^0$,
\begin{equation*}
g:E''\to \Pro (E')\tag{4.9.5}
\end{equation*}\etiquette[4.9.5]{VI.eq4.9.5}
若 $C$ 是 $E''$ 的生成子范畴, 配备诱导拓扑,
则已知 (II \hyperref[VI.subsec6.10]{6.10}),
给定 $G$ 也等价于
(在唯一同构意义下) 给定 $C$ 上取值于 $\Pro(E')^0$ 的一个层
\begin{equation*}
G_C \in \Ffaisc (C, \Pro(E')^0), G_C : C^{\circ} \to \Pro (E')^{0}{\quoi},\tag{4.9.6}
\end{equation*}\etiquette[4.9.6]{VI.eq4.9.6}
或等价于给定一个函子 $g_C=G_C^0$,
\begin{equation*}
g_C : C \to \Pro (E'){\quoi},\tag{4.9.7}
\end{equation*}\etiquette[4.9.7]{VI.eq4.9.7}
满足\pageoriginale{}熟悉的左正合条件
(II \hyperlink{sga4.exp.II.remark.6.2.item.1}{6.2.1})).
当然, $g_C$ 就是\emph{$g$ $\hyperref[VI.eq4.9.5]{(4.9.5)}$ 在 $C$ 上的限制}.
我们最感兴趣的情形是取 $C=E''_{\coh}$,
得到对象
\hypertarget{sga4.exp.VI.equation.4.9.8}{}\begin{align*}
& G_0 \in \Ffaisc (E''_{\coh}, \Pro (E')^{0}), G_0 : {E''}_{\coh}^0 \to
 \Pro (E')^0 {\quoi},\tag{4.9.8}\\
& g_0 = G^0_0 : E''_{\coh} \to \Pro(E'){\quoi},\tag{4.9.9}
\end{align*}\hypertarget{sga4.exp.VI.equation.4.9.9}{}
它们中的每一个仍等价于给定 $f$.

\setcounter{subsubsection}{9}
\subsubsection{}\etiquette[4.9.10]{VI.subsubsec4.9.10}

假定 $E''$ 的满生成子范畴 $C$ 由拟紧对象构成,
且在 $E''$ 中对有限和、过渡到等价关系的商以及纤维积稳定;
例如 $C=E''_{\coh}$ 就是如此
(\hyperref[VI.cor1.15]{1.15} 及 \hyperref[VI.cor1.17.1]{1.17.1}).
若 $P$ 是有限投射极限在其中可表示的范畴
(例如 $P=\Pro(E')^0$),
则立即可知,
函子 $G_C:C^{\circ}\to P$ 是取值于 $P$ 的层,
当且仅当函子 $g_C=G_C^0:C\to P^0$
与有限和以及过渡到等价关系的商交换.
这特别说明了哪些对象是层 \eqhyperref[VI.eq4.9.6]{4.9.6}
(从而表达了粘合函子 $f:E'\to E''$).

明确上述事实后, 可给出 \hyperref[VI.subsec4.9]{4.9} 中所提问题的解答:

\begin{enonce*}{命题 4.10}\etiquette[4.10]{VI.prop4.10}
设 $E'$ 与 $E''$\pageoriginale{} 是两个凝聚拓扑斯,
$f:E'\to E''$ 为粘合函子
\emph{(IV \hyperref[IV.eq9.3.6.3]{9.10})},
\ie 左正合且可及的函子.
设 $E$ 为由它经粘合得到的拓扑斯
\emph{(IV \hyperref[IV.theo9.5.4]{9.16})},
以 $E'_{\coh}$ (\resp $E'_{PF}$)
表示 $E'$ 中由凝聚对象
(\resp 由对象 $X$ 所表示的协变函子 $\Hom(X,-)$
与小滤过归纳极限交换的对象 $X$)
构成的严格满子范畴.
考察下列条件:
\begin{itemize}
\item[(i)] $E$ 凝聚.

\item[(ii)] 对 $E'$ 的每个对象 $X'$,
存在以 $X'$ 为靶、以 $E'$ 的凝聚对象为源的态射族
$$
v_{\alpha} : Y'_{\alpha} \to X'
$$
使 $E''$ 中各
$$
f(v_{\alpha}) : f(Y'_{\alpha}) \to f(X')
$$
构成的族覆盖.

\item[(ii bis)] $f$ 与 (小) 滤过归纳极限交换,
且对每个 $X'\in\Ob E'_{PF}$, \emph{(ii)} 成立.

\item[(ii ter)] 由粘合函子 $f$ 定义的典范函子
\emph{(IV \hyperref[IV.eq9.7.2.3]{9.20.1})}
\begin{equation*}
\pi : \Ppoint(E'') \to \Pro(E')\tag{4.10.1}
\end{equation*}\etiquette[4.10.1]{VI.eq4.10.1}
在同构意义下通过 $\Pro(E'_{\coh})$ 分解
(借助由包含 $E'_{\coh}\to E'$ 得到的全忠实函子
$\Pro(E'_{\coh})\to\Pro(E')$):
\begin{equation*}
\pi_0 : \Ppoint (E'') \to \Pro(E'_{\coh}){\quoi}.\tag{4.10.2}
\end{equation*}\etiquette[4.10.2]{VI.eq4.10.2}

\item[(iii)] 函子\pageoriginale{} $f$ 与 (小) 滤过归纳极限交换.

\item[(iii bis)] 函子 $g_0$
\hyperlink{sga4.exp.VI.equation.4.9.9}{(4.9.9)}
在同构意义下通过函子
\begin{equation*}
E''_{\coh} \to \Pro(E'_{PF}){\quoi},\tag{4.10.3}
\end{equation*}\etiquette[4.10.3]{VI.eq4.10.3}
分解, 借助由包含 $E'_{PF}\to E'$ 得到的全忠实函子
$\Pro(E'_{PF})\to\Pro(E')$.

\item[(iii ter)] 典范函子 $\pi$
\eqhyperref[VI.eq4.10.1]{4.10.1}
在同构意义下分解为
\begin{equation*}
\pi_1 : \Ppoint(E'') \to \Pro(E'_{PF}){\quoi}.\tag{4.10.4}
\end{equation*}\etiquette[4.10.4]{VI.eq4.10.4}

\item[(iv)] 函子 $g_0$
\hyperlink{sga4.exp.VI.equation.4.9.9}{(4.9.9)}
在同构意义下通过一个函子
\begin{equation*}
E''_{\coh} \to \Pro(E'_{\coh}){\quoi}.\tag{4.10.5}
\end{equation*}\etiquette[4.10.5]{VI.eq4.10.5}
分解. 此时有下列蕴含图:
\begin{equation*}
\text{\rm (iv) $\Rightarrow$ (i) $\Leftrightarrow$ (ii) $\Leftrightarrow$ (ii bis)
  $\Leftrightarrow$ (ii ter) $\Rightarrow$ (iii) $\Leftrightarrow$ (iii bis)
  $\Leftrightarrow$ (iii ter)}{\quoi}.\tag{4.10.6}
\end{equation*}\etiquette[4.10.6]{VI.eq4.10.6}
\end{itemize}
\end{enonce*}

立即记下如下推论.

\begin{enonce*}{推论 4.11}\etiquette[4.11]{VI.cor4.11}
\begin{enumerate}
\renewcommand{\theenumi}{\alph{enumi}}
\renewcommand{\labelenumi}{\theenumi)}
\item 假定 $E'$ 完美.
则 \hyperref[VI.prop4.10]{4.10} 中考察的一切条件都等价;
特别地, $E$ 凝聚, 当且仅当 $f$ 与滤过归纳极限交换,
或等价地, 当且仅当函子 $g_0$
\hyperlink{sga4.exp.VI.equation.4.9.9}{(4.9.9)}
在同构意义下\pageoriginale{}通过 $\Pro(E'_{\coh})$ 分解.

\item $E$ 完美的充要条件是 $E'$ 与 $E''$ 完美,
且 $f$ (\resp $g_0$) 满足 {\rm a)} 中所述条件.
\end{enumerate}
\end{enonce*}

事实上, a) 来自 $E'$ 完美意味着 $E'_{\coh}=E'_{PF}$,
从而 (iii bis) 蕴含 (iv) 这一事实.
断言 b) 则如 \hyperref[VI.subsec4.9]{4.9} 的预备说明所示.


\setcounter{subsection}{11}

% END INCLUDED SGA4_English_current_buildable_source_20260906/source/06/drafts/SGA4_VI_LINES_2419_2936_English.texfrag


% BEGIN INCLUDED SGA4_English_current_buildable_source_20260906/source/06/drafts/SGA4_VI_LINES_2937_3239_English.texfrag
\subsection{4.10 的证明}\etiquette[4.12]{VI.subsec4.12}

\begin{enumerate}
\renewcommand{\theenumi}{\alph{enumi}}
\renewcommand{\labelenumi}{\theenumi)}
\item 明确写出 \hyperref[VI.subsubsec4.9.3]{4.9.3}
中得到的与 (i) 等价的条件,
即 $E$ 的满子范畴 $E_0$ 生成:
\ie 对 $E$ 的每个对象
$X = (X', X'', u : X''\to f(X'))$,
满足 $Y \in \Ob E_0$,
\ie $Y'$ 与 $Y''$ 凝聚的所有态射
$$
v = (v', v'') : Y = (Y', Y'', v : Y''\to f(Y')) \to X {\quoi},
$$
构成满族.
由于函子对 $(i^{\ast}, j^{\ast})$ 保守,
这也就是说相应态射 $Y' \to X'$ 构成的族在 $E'$ 中为满族,
且相应态射 $Y''\to X''$ 构成的族在 $E''$ 中为满族.
对第一个族, 这显然成立:
只须取\emph{在 $U$ 上的} $Y \in \Ob E_0$,
\ie 满足 $Y'' = \varnothing_{E''}$ 的对象;
得到 $E'$ 中以 $X'$ 为靶且源凝聚的所有箭头构成的族,
它确实为满族,
因为 $E'$ 凝聚, 故子范畴 ${E'}_{\coh}$ 生成.
因此只须表达对每个给定的 $E$ 的对象 $X$,
族 $Y''\to X''$ 为满族的条件.

\item 再假定\pageoriginale{}对所考察的对象 $X'$, 条件 (ii) 满足.
由此可找到一个满态射族 $Y''_{\beta} \to X''$,
其与 $u : X''\to f(X')$ 的复合,
各自提升为适当 $\alpha = \phi(\beta)$ 的态射
$Y''_{\beta} \to f({Y'_{\alpha}})$.
由于 $E''_{\coh}$ 是 $E''$ 的生成子范畴,
还可以假定各 $Y''_{\beta}$ 凝聚.
交换图
\[
\xymatrix@C=2.25cm@R=1.25cm{
Y''_{\beta}  \ar[r] \ar[d]& f(Y'_{\varphi(\beta)}) \ar[d]\\
X'' \ar[r] & f(X')
}
\]
于是给出 $E$ 中所需的态射族,
形成满族 $Y''_{\beta} \to X''$.
因此 (ii) $\Rightarrow$ (i).

反过来, 假定 (i), 证明 (ii).
将 a) 中明确写出的条件应用于对象
$$
X = j_{\ast}(X') = (X', f(X'), \id = f(X') \xrightarrow{\sim} f(X')){\quoi}.
$$
由于各 $Y''\to X''\simeq f(X')$ 构成满族,
各 $f(v') : f(Y') \to f(X')$ 构成的族更为满族.
由于各 $Y'$ 由假设凝聚, 得到结论.
因此
$$
{\rm(i)} \Leftrightarrow {\rm(ii)}{\quoi}.
$$

\item 再注意\pageoriginale{},
条件 (iv) 也意味着对 $E''$ 的每个对象 $X''$,
$E'$ 上的 pro-可表示函子
$$
g(X'') : X'\longmapsto \Hom (X'', f(X'))
$$
由一个各分量属于 ${E'}_{\coh}$ 的 pro-对象 pro-表示.
等价地, 在由 $E'$ 的对象 $X'$
和元素 $u \in g(X'')(X') =\Hom(X'', f(X'))$,
$u : X''\to f(X')$, 构成的对所成的范畴 $E'_{/g(X'')}$ 中,
存在共尾对象集 $(X'_{\alpha},u_{\alpha})$,
其中 $u_{\alpha} : X'' \to f(X'_{\alpha})$
且 $X'_{\alpha} \in \Ob E'_{\coh}$.
由于范畴 $E'_{/g(X'')}$ 滤过,
共尾条件就是说对 $E'_{/g(X'')}$ 的每个对象 $(X',u)$,
$u : X''\to f(X')$,
都存在从某个 $(X'_{\alpha},u_{\alpha})$ 到此对象的态射,
\ie 一个使三角形
\[\vcenter{
\xymatrix@C=2cm@R=1.5cm{
                      &f(X'_{\alpha}) \ar[d]^{f(v'_\alpha)}\\
X''\ar@{-}[ru]^{u_{\alpha}} \ar[r]& f(X').
}}\leqno{(*)}
\]
交换的态射 $v'_{\alpha} : X'_{\alpha} \to X'$.
因此条件 (iv) 等价于下面的条件:

\begin{itemize}
\item[(iv bis)] 对每个箭头 $u : X''\to f(X')$,
$X'\in \Ob E'$, $X''\in \Ob E''_{\coh}$,
存在 $E'$ 中源 $Y'$ 凝聚的箭头 $v' : Y'\to X'$,
使 $u$ 通过 $f(v') : f(Y') \to f(X')$ 分解.

显然条件 (iv bis) 蕴含 (ii),
因为对固定的 $X'$, 源凝聚的各 $u : X''\to f(X')$ 构成满族.
因而各相应图 $(\ast)$ 中所有 $f(v'_{\alpha})$ 构成的族更为满族.
因此
$$
{\rm(iv)} \Rightarrow {\rm(i)}{\quoi}.
$$
\end{itemize}

\item 明确写出条件 (ii ter).
设 $p$ 是 $E''$ 的一个点.
由定义, $\pi(p)$ 是 pro-表示函子
$h : X'\longmapsto f(X')_{p}$ 的 $E'$ 的 pro-对象;
或者等价地, 是由典范函子 $E'_{/h} \to E'$ 定义的 pro-对象,
其中 $E'_{/h}$ 表示由 $X'\in \Ob E'$ 及 $u\in h(X')$,
\ie $u\in f(X')_p$, 构成的对 $(X',u)$ 的范畴.
说这个 pro-对象同构于一个来自 $\Pro(E'_{\coh})$ 的 pro-对象,
意味着 $E'_{/h}$ 中由 $X'\in\Ob E'_{\coh}$ 的各 $(X',u)$
构成的子范畴共尾.
换句话说, 对 $E'_{/h}$ 的每个对象 $(X',u)$,
$u\in f(X')_p$,
存在一个映向它的对象 $(X'_0,u_0)$, $X'_0$ 凝聚;
即存在态射 $X'_0\to X'$ ($X'_0$ 凝聚),
使 $u\in\text{Im}(f(X'_0)\to f(X'))$.
称这对每个点 $p$ 都成立,
意味着对 $E'$ 的每个对象 $X'$,
各 $X'_0$ 凝聚的映射 $f(X'_0)\to f(X')$ 所成的族在每个茎上满射.
我们将在附录 (10) 中看到凝聚拓扑斯 $E''$ 有足够多的点,
故所得条件也就是说对每个 $X'\in\Ob E'$,
族 $f(X'_0)\to f(X')$ 为满族.
这正是条件 (ii), 从而
$$
{\rm(ii)} \Leftrightarrow \text{(ii ter)}{\quoi}.
$$

\item $f$ 与滤过归纳极限交换的条件\pageoriginale{}
(条件 (iii)),
由于 $E''$ 的纤维函子族保守,
等价于复合函子 $h : X'\longmapsto f(X')_p$ 具有此性质.
下面证明, 这意味着 pro-表示此函子的 pro-对象 $\pi(p)$
属于 $\Pro(E'_{PF})$ 的本质像.
这就证明蕴含
$$
{\rm(iii)} \Leftrightarrow \text{(iii ter)} \Leftarrow \text{(ii ter)}
\Rightarrow \text{(iii bis)}{\quoi}.
$$
因此归结为证明下列引理.

\begin{enonce*}{引理 4.12.1}\etiquette[4.12.1]{VI.lem4.12.1}
设 $E'$ 是一个子范畴 $E'_{\coh}$ 生成的拓扑斯,
$X' \in \Ob \Pro(E')$.
则 $X'$ 属于 $\Pro(E'_{PF})$ 的本质像,
当且仅当它 pro-表示的函子
$h : E'\to (\Ens)$ 与滤过归纳极限交换.
\end{enonce*}

由于与滤过归纳极限交换的函子的滤过归纳极限仍具有相同性质,
充分性来自 $E'_{PF}$ 的定义.
证明必要性, 须证明若函子 $h$ 与滤过归纳极限交换,
则在由满足 $u\in h(X')$ 的对 $(X',u)$ 构成的滤过范畴 $E'_{/h}$ 中,
各 $X'\in\Ob E'_{PF}$ 的对 $(X',u)$ 构成的子范畴共尾.
就是说, 对 $E'_{/h}$ 的每个 $(X',u)$,
存在态射 $X'_0\to X'$, $X'_0\in\Ob E'_{PF}$,
使 $u\in\text{Im}(h(X'_0)\to h(X'))$.\pageoriginale{}
但已知 $X'$ 是 $E'_{PF}$ 的对象 $X'_i$ 的滤过归纳极限
(\hyperref[VI.cor1.24.2]{1.24.2} a)),
故 $h(X')$ 是各 $h(X'_i)$ 的滤过归纳极限,
这就证明了断言.

\item 只须再证明蕴含
$$
\text{(ii bis)} \Rightarrow \text{(ii)\quad 以及\quad (iii)}
\Leftrightarrow \text{(iii bis)}{\quoi}.
$$
第一个蕴含显然来自 $E'$ 的每个对象 $X'$
都是 $E'_{PF}$ 的对象 $X'_i$ 的滤过归纳极限这一事实.
对第二个蕴含, 只须注意对 $X''\in\Ob {E''}_{\coh}$ 的函子族
$$
Y'' \longmapsto \Hom(X'', Y'') {\quoi},
$$
保守, 且由与滤过归纳极限交换的函子构成
(\hyperref[VI.thm1.23]{1.23} (ii)).
因此, 函子 $f : E'\to E''$ 与滤过归纳极限交换,
当且仅当复合函子
$X'\longmapsto \Hom(X'',f(X'))=
\Hom_{\Pro(E')}(g_0(X''),X')$ 具有此性质.
由 \hyperref[VI.lem4.12.1]{4.12.1},
这等价于各 $g_0(X'')$ 属于 $\Pro(E'_{PF})$ 的本质像. \cqfd
\end{enumerate}

\begin{enonce*}[remark]{注 4.13}\etiquette[4.13]{VI.rem4.13}
用层论语言说明条件 \hyperref[VI.prop4.10]{4.10} (iv) 的含义很有用.
考察函子的交换图
\[\vcenter{
\xymatrix@C=2.5cm@R=1.5cm{
\Pro(E'_{\coh})^{\circ}  \ar@{^{(}->}[r]^{\alpha}  \ar[d]^{\beta_0}&
\Pro(E')^{\circ}_{\beta} \ar[d]^{\beta}\\
\text{$\Ev{\coh}$} & \text{$\Ev{}$}\ar[l]^{\gamma} ,
}}\leqno{(4.13.1)}
\]\etiquette[4.13.1]{VI.eq4.13.1}
其中\pageoriginale{}上标圆圈表示反范畴,
第二个横箭头 $\gamma$ 为函子的限制,
其余箭头都表示熟悉的全忠实函子.
注意, 包含箭头 $\alpha$ \emph{一般不与投射极限交换},
即使有限投射极限也不例外.
因此, 一般说来, 景 $C$ (例如 $E''$) 上的层 $F$
在同一景上定义取值于 $\Pro(E'')^{0}$ 的预层 $\alpha\circ F$,
它未必为层.
不过, 若 $C$ 上取值于 $\Pro(E'_{\coh})^{0}$ 的一个\emph{预层} $F$,
使相应的取值于 $\Pro(E')^{0}$ 的预层 $\alpha F$ 为层,
则 $F$ 本身为层.
事实上, 层条件为左正合性质
(II \hyperlink{sga4.exp.II.remark.6.2.item.1}{6.2.1})),
而函子 $\beta$、$\beta_0$ 及 $\gamma$ 与小投射极限交换.
\end{enonce*}

另一方面, \hyperref[VI.subsubsec4.9.10]{4.9.10}
对 $E''_{\coh}$ 上取值于任意范畴 $P$ 的层的刻画表明:
对每个使 $\alpha^{0}:P^0\to Q^0$
与\emph{有限余积}及\pageoriginale{}\emph{过渡到等价关系的商}交换的函子
$\alpha:P\to Q$,
从 $\Hom({E''}^0_{\coh},P)$ 到 $\Hom({E''}^0_{\coh},Q)$
的函子 $g\mapsto\alpha\circ g$ 将层映为层.
还知道, 在 $\Pro(E'_{\coh})$ 及 $\Pro(E')$ 中,
有限余积与等价关系的商可表示,
且包含 $E'_{\coh}\to E'$ 与它们交换
(\hyperref[VI.cor1.15]{1.15} 及 \hyperref[VI.cor1.17.1]{1.17.1}).
因此, 对范畴 $\Pro(E'_{\coh})$ 及 $\Pro(E')$
以及它们之间的包含, 同样成立
(I \hyperref[I.prop8.9.5]{8.9.5} b) 及 \hyperref[I.prop8.9.7]{8.9.7});
而这一包含正是 $\alpha^0$,
其中 $\alpha$ 为 \eqhyperref[VI.eq4.13.1]{4.13.1} 中的函子.
由此可知, 若
\begin{equation*}
\varphi:{E''}^0_{\coh}\to\Pro(E'_{\coh})^0\tag{4.13.2}
\end{equation*}\etiquette[4.13.2]{VI.eq4.13.2}
为函子, 则 \emph{$\varphi$ 是 $E''_{\coh}$ 上取值于
$\Pro(E'_{\coh})^0$ 的层,
当且仅当 $\alpha\circ\varphi$ 是 $E''_{\coh}$ 上取值于
$\Pro(E')^0$ 的层}.
因此, 满足条件 \hyperref[VI.prop4.10]{4.10} (iv)
的粘合函子 $f:E'\to E''$ 的范畴,
典范等价于 $E''_{\coh}$ 上取值于 $\Pro(E'_{\coh})^0$ 的层 $\varphi$ 的范畴.
考虑 \hyperref[VI.prop4.10]{4.10},
得到这样的层 $\varphi$ 典范地定义粘合函子
(在唯一同构意义下),
由后者得到的拓扑斯 $E$ \emph{凝聚}.
此外 (\hyperref[VI.cor4.11]{4.11}),
\emph{若 $E'$ 完美},
就由此得到可以由 $E'$ 与 $E''$ 经粘合得到的\emph{所有}凝聚拓扑斯.

还注意\pageoriginale{}, $E'_{\coh}$ 等价于小范畴,
并对有限投射极限稳定.
因此 (I \hyperref[I.sec8.10.14]{8.10.14}),
借助 \eqhyperref[VI.eq4.13.1]{4.13.1} 的全忠实函子 $\beta_0$,
$\Pro(E'_{\coh})^0$ 等价于
$\widehat{E'_{\coh}}=\Hom(E'_{\coh},(\Ens))$
中由左正合函子构成的满子范畴.
因此, 给定一个层 \eqhyperref[VI.eq4.13.2]{4.13.2},
等价于给定一个函子
\begin{equation*}
F:{E''}_{\coh}^0\times E'_{\coh}\to(\Ens)\tag{4.13.3}
\end{equation*}\etiquette[4.13.3]{VI.eq4.13.3}
它关于第一个自变量为层, 关于第二个自变量左正合;
或者等价地, 给定一个左正合函子
\begin{equation*}
f_0:E'_{\coh}\to E''\tag{4.13.4}
\end{equation*}\etiquette[4.13.4]{VI.eq4.13.4}
若 $f$ 是与 $\varphi$ 相伴且满足条件
\hyperref[VI.prop4.10]{4.10} (iv) 的粘合函子,
立即可验证, $f_0$ 典范同构于 $f:E'\to E''$ 的限制.
因此, \emph{在所考察的条件下, $f_0$ 确定 $f$}
(在典范同构意义下).
为概括 \hyperref[VI.cor4.11]{4.11} 的内容, 最后陈述:

\begin{enonce*}{推论 4.14}\etiquette[4.14]{VI.cor4.14}
当 $E'$ 是完美拓扑斯时,
由粘合给出凝聚拓扑斯 $E$ 的粘合函子 $f:E'\to E''$ 的范畴,
在 $f\mapsto f_0=f|E'_{\coh}$ 下,
等价于左正合函子 $f_0:E'_{\coh}\to E''$ 的范畴;
或者等价地, 等价于\pageoriginale{}景 $E''_{\coh}$ 上
(或同样地, 拓扑斯 $E''$ 上)
取值于范畴 $\Pro(E'_{\coh})^0$ 的层 $\varphi$ 的范畴
\emph{(\cf \hyperref[VI.subsubsec4.9.10]{4.9.10})}.
\end{enonce*}

第一个断言还可以由 \hyperref[VI.prop4.10]{4.10} 得到,
只须注意 $E'\approxeq\Ind(E'_{\coh})$
(\hyperref[VI.cor1.25]{1.25}).
事实上, 限制给出如下两个范畴的等价:
与滤过归纳极限交换的函子 $f:E'\to P$ 的范畴
(其中 $P$ 为小滤过归纳极限可表示的范畴),
以及任意函子 $f:E'_{\coh}\to P$ 的范畴.
立即可知, 在这些条件下,
若 $P$ 中有限投射极限可表示,
则 $f$ 左正合, 当且仅当 $f_0$ 左正合.

\begin{enonce*}[remark]{注 4.15}\etiquette[4.15]{VI.rem4.15}
\begin{enumerate}
\renewcommand{\theenumi}{\alph{enumi}}
\renewcommand{\labelenumi}{\theenumi)}
\item \hyperref[VI.prop4.10]{4.10} 的证明表明,
将范畴 $\Ppoint(E'')$ 换成一个定义\emph{保守}纤维函子族的子范畴,
可得与 (ii ter) \resp (iii ter) 等价的条件.

\item 取 $E''$ 为点拓扑斯
(IV \hyperref[IV.subsec2.2]{2.2}),
立即可见, 条件 (ii ter) 与 (iii ter)
等价\emph{仅当} $E'_{\coh}=E'_{PF}$,
\ie 仅当 $E'$ 完美.
因此, 条件“$E$ 凝聚”等价于条件“$f$ 与滤过归纳极限交换”,
仅当 $E'$ 完美.
另一方面, 更一般地取 $E''$ 为形如 $\widehat{C}$ 的凝聚拓扑斯,
可见\emph{在此情形, 对每个 $E'$, 条件} (i) \emph{等价于} (iv).
事实上, 在一般情形,
编者一直未能构造 $E$ 凝聚而条件 (iv) 不满足的例子.
实在惭愧!
\end{enumerate}
\end{enonce*}


% END INCLUDED SGA4_English_current_buildable_source_20260906/source/06/drafts/SGA4_VI_LINES_2937_3239_English.texfrag


% BEGIN INCLUDED SGA4_English_current_buildable_source_20260906/source/06/drafts/SGA4_VI_LINES_3240_3413_English.texfrag
\section{函子 $H^i(X,-)$ 与滤过归纳极限的交换}\etiquette[5]{VI.sec5}\pageoriginale

\begin{enonce*}{定理 5.1}\etiquette[5.1]{VI.thm5.1}
设 $E'$ 是一个代数拓扑斯
(\hyperref[VI.def2.3]{2.3}),
$E$ 是一个局部凝聚拓扑斯
(\hyperref[VI.def2.3]{2.3}),
$f:E'\to E$ 是凝聚拓扑斯态射
(\hyperref[VI.def3.1]{3.1}).
则对每个整数 $q$, 函子 $R^qf_{\ast}$
\emph{(V \hyperref[V.subsec5.0]{5.0})}
与 Abel 层的滤过归纳极限交换.
\end{enonce*}

\begin{enonce*}{推论 5.2}\etiquette[5.2]{VI.cor5.2}
设 $E'$ 是一个凝聚拓扑斯
(\hyperref[VI.def2.3]{2.3}).
对每个整数 $q$, 函子 $H^q(E',-)$
与 Abel 层的滤过归纳极限交换.
\end{enonce*}

\begin{enonce*}{推论 5.3}\etiquette[5.3]{VI.cor5.3}
设 $E$ 是一个拓扑斯 (\resp 代数拓扑斯)
(\hyperref[VI.def2.3]{2.3}),
$X$ 是 $E$ 的代数且凝聚 (\resp 凝聚) 对象
(\hyperref[VI.def2.3]{2.3}).
对每个整数 $q$, 函子 $H^q(X,-)$
与 Abel 层的滤过归纳极限交换.
\end{enonce*}

在 \hyperref[VI.thm5.1]{5.1} 中取 $E$ 为点拓扑斯,
$f:E'\to E$ 为唯一的态射
(IV \hyperref[IV.subsec2.2]{2.2}),
即得推论 \hyperref[VI.cor5.2]{5.2}.
证明 \hyperref[VI.cor5.2]{5.2} 蕴含 \hyperref[VI.cor5.3]{5.3}.
设 $j_X:E_{/X}\to E$ 为局部化态射.
有关于 Abel 层 $F$ 函子性的典范同构
(V \hyperref[V.2.2.1]{2.2.1}),
$H^q(X,F)\simeq H^q(E_{/X},j_X^{\ast}F)$.
函子 $j_X^{\ast}$ 与归纳极限交换,
且拓扑斯 $E_{/X}$ 凝聚
(\hyperref[VI.subsubsec2.4.6]{2.4.6}),
从而得到 \hyperref[VI.cor5.3]{5.3}.
证明 \hyperref[VI.cor5.3]{5.3} 蕴含 \hyperref[VI.thm5.1]{5.1}.
设 $E_{\text{cohalg}}$ 为 $E$ 中由 $E$ 的代数且凝聚的对象定义的满子范畴.
它是生成范畴 (\hyperref[VI.subsubsec2.4.5]{2.4.5}).
对每个 Abel 层 $F$,
$R^qf_{\ast}F$ 为预层
$X\mapsto H^q(f^{\ast}X,F)$ $(X\in\ob C)$ 的相伴层
(\hyperrefext[V][V.prop5.1]{5.1}).
由于 $f$ 凝聚,
对 $C$ 的每个对象 $X$, $f^{\ast}X$ 凝聚
(\hyperref[VI.def3.1]{3.1}).
由 \hyperref[VI.cor5.3]{5.3},
预层 $X\mapsto H^q(f^{\ast}X,F)$
关于自变量 $F$ 与滤过归纳极限交换;
而相伴层函子与归纳极限交换, 故得到 \hyperref[VI.thm5.1]{5.1}.
(这里只用到态射 $f$ 的下列性质:
$E$ 存在一个生成族 $S$,
使对 $S$ 的每个 $X$, $f^{\ast}(X)$ 为代数且凝聚的.)
因此只须证明 \hyperref[VI.cor5.2]{5.2}.
已经知道函子 $H^{\circ}(E',-)$
与 Abel 层的滤过归纳极限交换
\hyperref[VI.thm1.23]{(1.2.3)}.
由 \cite{sga4.exp.VI.bib.7} 的一个经典结果,
要证明 \hyperref[VI.cor5.2]{5.2},
只须证明对函子 $H^{\circ}(E',-)$ 无上同调的层的滤过归纳极限,
对 $H^{\circ}(E',-)$ 仍无上同调.
这来自下述引理, 其中取 $C=E'_{\coh}$.

\begin{enonce*}{引理 5.4}\etiquette[5.4]{VI.lem5.4}
设\pageoriginale{} $E$ 是一个局部凝聚拓扑斯
(\hyperref[VI.def2.3]{2.3}),
$C$ 是 $E$ 的一个对纤维积稳定、
所有对象均拟紧 (\hyperref[VI.def1.1]{1.1}) 的满生成子范畴.
$C$-无上同调层 \emph{(\hyperrefext[V][V.subsec4.2]{4.2})}
的滤过归纳极限仍为 $C$-无上同调层.
\end{enonce*}

设 $(F_i)_{i\in I}$ 是 $C$-无上同调层的一个滤过归纳系,
以 $F$ 表示其归纳极限.
$C$ 的每个对象 $Y$ 都凝聚 (\hyperref[VI.prop2.1]{2.1}),
因此 $F(Y)=\displaystyle\mathop{\varinjlim}_iF_i(Y)$
\hyperref[VI.thm1.23]{(1.2.3)}.
要证明 $F$ 为 $C$-无上同调的,
只须证明对每个 $q>0$ 及 $C$ 的每个 $Y$,
$\check H^q(Y,F)=0$
(\hyperrefext[V][V.prop4.3]{4.3}).
设 $X$ 是 $C$ 的对象,
$\Xx=(X_{\alpha}\to X)_{\alpha\in A}$
是由 $C$ 的对象构成的一个有限覆盖族.
有 $H^q(\Xx,F)=H^q(C^{\bigdot}(\Xx,F))$
(\hyperrefext[V][V.subsubsec2.4.3]{2.4.3}).
由前述结果及 $C^{\bigdot}(\Xx,F)$ 只涉及群的有限积这一事实
(\hyperrefext[V][V.subsubsec2.3.3]{2.3.3}), 有
$C^{\bigdot}(\Xx,F)=\varinjlim_iC^{\bigdot}(\Xx,F_i)$.
因此
$$
H^q(\Xx,F)
=H^q\bigl(\varinjlim_iC^{\bigdot}(\Xx,F_i)\bigr)
=\varinjlim_iH^q(C^{\bigdot}(\Xx,F_i))=0
$$
对 $q>0$ 成立 (\hyperrefext[V][V.prop4.3]{4.3}).
因而 $\check H^q(X,F)=\varinjlim_{\mathfrak X}H^q(\mathfrak X,F)=0$.

\begin{enonce*}{推论 5.5}\etiquette[5.5]{VI.cor5.5}
设 $E$ 是一个凝聚拓扑斯
(\hyperref[VI.def2.3]{2.3}),
${(X_i)_{i\in I}}$ 为终对象的凝聚子对象的一个递减滤过族.
以 $\Phi$ 表示 $E$ 的闭子拓扑斯
\emph{(IV 9)} 中那些包含于某个 $X_i$ 的闭补的子拓扑斯所成的族.
则 $\Phi$ 为 $E$ 的支族
\emph{(\hyperrefext[V][V.def6.12]{6.12})},
并且对每个整数 $q$,
函子 $H^q_{\Phi}(E,-)$ 与 $H^q_{\Phi}$
\emph{(\hyperrefext[V][V.subsec6.13]{6.13})}
同滤过归纳极限交换.
\end{enonce*}

设 $Z_i$ 是 $X_i$ 的闭补.
对每个 Abel 层 $F$, 有正合序列
(\hyperrefext[V][V.prop6.5]{6.5}):
$$
0 \to H^{0}_{Z_i} (E, F) \to H^{0} (E, F) \to H^{0} (X_i, F) \to
H^{1}_{Z_i} (E, F) \to \ldots\quad {\quoi}.
$$
函子 $H^q(E,F)$ 与 $H^q(X_i,F)$
关于 $F$ 与滤过归纳极限交换
(\hyperref[VI.cor5.2]{5.2}, \hyperref[VI.cor5.3]{5.3}).
因此 $H^q_{Z_i}(E,F)$ 也如此;
再对各 $Z_i$ 取归纳极限
(\hyperrefext[V][V.subsec6.13]{6.13}),
即得 $H^q_{\Phi}(E,F)$ 的相同性质.
关于函子 $\mathcal H^q_{\Phi}$ 的断言,
由局部化及过渡到相伴层得到.

\begin{enonce*}[remark]{定义 5.6}\etiquette[5.6]{VI.def5.6}
设\pageoriginale{} $E$ 是一个环为 $A$ 的赋环拓扑斯,
$q$ 为整数.
$A$-模 $G$ 称为有限 $q$-呈示的,
是指存在 $E$ 的一个对象族 $X_i$ 覆盖 $E$ 的终对象,
使对每个 $i$, $E_{/X_i}$ 上存在模的正合序列:
\begin{equation*}
L_q \to L_{q-1} \to \cdots L_1 \to L_0 \to G_{/X_i} \to 0\tag{5.5.1}
\end{equation*}\etiquette[5.5.1]{VI.eq5.5.1}
其中对每个 $p$, $L_p=A^{n_p}_{/X_i}$,
且 $n_p$ 为整数.
\end{enonce*}

\begin{enonce*}{命题 5.7}\etiquette[5.7]{VI.prop5.7}
设 $G$ 为有限 $q$-呈示的层.
对每个 $i\leq q-1$,
函子 $\Eext^i_A(G,-)$ 与 $A$-模的滤过归纳极限交换.
\end{enonce*}

问题在 $E$ 上局部 (\hyperrefext[V][V.prop6.1]{6.1}).
因此局部化到各 $E_{/X_i}$ 后,
可以假定存在形如 \eqhyperref[VI.eq5.5.1]{5.5.1} 的正合序列.
置 $L^{\bigdot}=L_q\to\cdots\to L_1\to L_0$.
有 $\Eext^i_A(G,F)\simeq\Hh^i(\Hhom_A(L^{\bigdot},F))$,
且函子 $F\mapsto\Hhom_A(L^{\bigdot},F)$ 与归纳极限交换,
这就证明了命题.

\begin{enonce*}{推论 5.8}\etiquette[5.8]{VI.cor5.8}
设 $(E,A)$ 为赋环拓扑斯,
$G$ 为有限 $q$-呈示的 $A$-模层,
$X$ 是 $E$ 的代数且凝聚的对象.
对 $\frac{i(i-1)}{2}\leq q-1$,
函子 $F\mapsto\Ext^i_A(X;G,F)$ 关于 $F$ 与滤过归纳极限交换.
\end{enonce*}

这由 \hyperref[VI.prop5.7]{5.7} 及 \hyperref[VI.cor5.3]{5.3},
通过 V \hyperref[V.6.1.3]{6.1.3} 的谱序列得到.

\setcounter{subsection}{8}
\subsection{}\etiquette[5.9]{VI.subsec5.9}

设 $(E,A)$ 为赋环拓扑斯,
$\Phi$ 是 $E$ 的支族
(\hyperrefext[V][V.def6.12]{6.12}),
$G$ 为有限 $r$-呈示的 $A$-模
(\hyperref[VI.def5.6]{5.6}),
$X$ 是 $E$ 的代数且凝聚的对象
(\hyperref[VI.def2.3]{2.3}).
分别对 V \hyperref[V.6.9.1]{6.9.1}
及 V \hyperref[V.6.9.2]{6.9.2} 的最后几个谱序列
中 $\Phi$ 内的闭子拓扑斯取归纳极限,
得到两个谱序列 (\hyperrefext[V][V.subsec6.13]{6.13}):
\begin{equation*}
\begin{cases}
'E^{pq}_2 = \varinjlim_{Z \in \Phi} \Ext^p_A (X; G, H^q_Z F)
\Rightarrow \Ext^{p+q}_{A, \Phi} (X; G, F){\quoi},\\
''E^{pq}_2 = \varinjlim_{Z \in \Phi} \Eext^p_A (G, H^q_Z F)
\Rightarrow \Eext^{p+q}_{A, \Phi} (G, F){\quoi}.
\end{cases}\tag{5.9.1}
\end{equation*}\etiquette[5.9.1]{VI.eq5.9.1}

由 \hyperref[VI.prop5.7]{5.7} 及 \hyperref[VI.cor5.8]{5.8},
有典范同构
\begin{equation*}
\begin{cases}
'E^{pq}_2 \simeq \Ext^p_A (X; G, H^q_\Phi F)&\quav
\frac{p(p-1)}{2} \leq r - 1{\quoi},\\
''E^{pq}_2 \simeq \Eext^p_A (G, H^q_\Phi F) &\quav p \leq r-1{\quoi}.
\end{cases}\tag{5.9.2}
\end{equation*}\etiquette[5.9.2]{VI.eq5.9.2}

\begin{enonce*}{推论 5.10}\etiquette[5.10]{VI.cor5.10}
采用\pageoriginale{} \hyperref[VI.cor5.5]{5.5} 的假设及记号.
设 $A$ 为 $E$ 的环,
$G$ 为有限 $q$-呈示的 $A$-模.
对满足 $\dfrac{i(i-1)}{2}\leq q-1$ 的每个整数 $i$,
函子 $F\mapsto\Ext^i_{A,\Phi}(E;G,F)$
及 $F\mapsto\Eext^i_{A,\Phi}(G,F)$
\emph{(\hyperrefext[V][V.subsec6.13]{6.13})}
与滤过归纳极限交换.
\end{enonce*}

这由 \hyperref[VI.cor5.5]{5.5} 及 \hyperref[VI.prop5.7]{5.7},
分别通过 V \hyperref[V.6.9.1]{6.9.1}
及 V \hyperref[V.6.9.2]{6.9.2} 的第一个谱序列得到.


% END INCLUDED SGA4_English_current_buildable_source_20260906/source/06/drafts/SGA4_VI_LINES_3240_3413_English.texfrag


% BEGIN INCLUDED SGA4_English_current_buildable_source_20260906/source/06/drafts/SGA4_VI_LINES_3414_3872_English.texfrag
